\chapter{Complexity — Primitives}
%%% generated by the blueprint pipeline from TCSlib.Complexity.TuringMachine.Build.Primitives
%%%%%%%%%%%%%%%%%%%%%%%%%%%%%%%%%%%%%%%%%%%%%%%%%%%%%%%%%%%%%%%%%%%%%%%%%%%%%%%
\section{Overview}

This module builds the primitive catalog of the machine-construction library:
small multi-tape Turing machines over the binary alphabet --- fixed-prefix
emission, duplication into the self-delimiting pair code, fixed-width
increment, pair-code validity testing and component extraction, unary
polynomial generation, and length comparison against a captured result ---
each verified by exact run invariants and packaged as an existential machine
contract with an explicit linear or polynomial time bound. Throughout, bits
are written $0$ and $1$, words are elements of $\{0,1\}^*$, and the
self-delimiting pair code of $(a,b)$ doubles each bit of $a$, then appends the
separator bits $0,1$, then $b$.

%%%%%%%%%%%%%%%%%%%%%%%%%%%%%%%%%%%%%%%%%%%%%%%%%%%%%%%%%%%%%%%%%%%%%%%%%%%%%%%
\section{Declarations}

\begin{definition}[Fixed-prefix emitter machine]
\lean{Turing.FinTM.catalogPrefixTM}
\label{Turing.FinTM.catalogPrefixTM}
\leanok
\uses{Turing.FinTM}
For a fixed word $w \in \{0,1\}^*$, a finite Turing machine with no work tapes
and $\abs{w}+1$ control states that first writes the bits of $w$ to its output
tape, one per step, and then copies its input to the output verbatim.
\end{definition}

\begin{definition}[Prefix-emitter configuration]
\lean{Turing.FinTM.catalogPrefixCfg}
\label{Turing.FinTM.catalogPrefixCfg}
\leanok
\uses{Turing.Cfg}
A configuration of the fixed-prefix emitter on input $x$, specified by an
optional control state, an input-head position, and the output written so far;
the machine has no work tapes, so the work-tape fields are trivial.
\end{definition}

\begin{lemma}[Emission phase of the prefix emitter]
\lean{Turing.FinTM.catalogPrefixTM_emit}
\label{Turing.FinTM.catalogPrefixTM_emit}
\leanok
\uses{Turing.Action.apply, Turing.Cfg.ext_zero_tapes, Turing.FinTM.catalogPrefixCfg, Turing.FinTM.catalogPrefixTM, Turing.MultiTapeTM.initCfg, Turing.MultiTapeTM.runFrom, Turing.MultiTapeTM.runFrom_succ_eq_step', Turing.MultiTapeTM.step}
\difficulty{4}
For every $i \le \abs{w}$, after $i$ steps from its initial configuration on
input $x$ the fixed-prefix emitter for $w$ is in control state $i$, the input
head has not moved, and the output is the first $i$ bits of $w$.
\end{lemma}

\begin{lemma}[Copy phase of the prefix emitter]
\lean{Turing.FinTM.catalogPrefixTM_copy}
\label{Turing.FinTM.catalogPrefixTM_copy}
\leanok
\uses{Turing.Action.apply, Turing.Cfg.ext_zero_tapes, Turing.Cfg.inputSymbol, Turing.FinTM.catalogPrefixCfg, Turing.FinTM.catalogPrefixTM, Turing.MultiTapeTM.runFrom, Turing.MultiTapeTM.runFrom_succ_eq_step', Turing.inputSymbolInner, Turing.moveInputPos, Turing.moveInputPos_pos_of_ne_right}
\difficulty{4}
For every $i \le \abs{x}$, after $i$ steps from the configuration in which the
whole fixed word $w$ has been emitted and the input head sits at the start of
$x$, the prefix emitter has advanced the input head by $i$ cells and its
output is $w$ followed by the first $i$ bits of $x$.
\end{lemma}

\begin{lemma}[Time bound for prefix emission]
\lean{Turing.FinTM.catalogPrefixTM_computes}
\label{Turing.FinTM.catalogPrefixTM_computes}
\leanok
\uses{Turing.Action.apply, Turing.Cfg.inputSymbol, Turing.FinTM.ComputesFunInTime, Turing.FinTM.catalogPrefixCfg, Turing.FinTM.catalogPrefixTM, Turing.FinTM.catalogPrefixTM_copy, Turing.FinTM.catalogPrefixTM_emit, Turing.FinTM.computesInTime_iff, Turing.MultiTapeTM.runFrom_add, Turing.MultiTapeTM.runFrom_succ_eq_step', Turing.MultiTapeTM.step}
\difficulty{3}
For every fixed word $w$ and every input $x$, the prefix emitter for $w$
halts on $x$ within $\abs{w} + \abs{x} + 1$ steps, and its output is the
concatenation $w x$; the final step is the halting transition on the right
blank.
\end{lemma}

\begin{definition}[Zero-work-tape scan configuration]
\lean{Turing.FinTM.scanCfg}
\label{Turing.FinTM.scanCfg}
\leanok
\uses{Turing.Cfg}
A configuration of a machine with no work tapes on input $x$, given by an
optional control state, the number $i \le \abs{x}$ of input cells already
passed, and the output written so far.
\end{definition}

\begin{lemma}[Input symbol of a scan configuration]
\lean{Turing.FinTM.scanCfg_read}
\label{Turing.FinTM.scanCfg_read}
\leanok
\uses{Turing.Cfg.inputSymbol, Turing.FinTM.scanCfg, Turing.inputSymbolInner}
\difficulty{3}
In a scan configuration indexed by $i \le \abs{x}$, the symbol under the input
head is the $i$-th bit of $x$ when $i < \abs{x}$, and the blank when
$i = \abs{x}$.
\end{lemma}

\begin{lemma}[Copy-state run invariant]
\lean{Turing.FinTM.scanCopy_run}
\label{Turing.FinTM.scanCopy_run}
\leanok
\uses{Turing.Action.apply, Turing.Cfg.ext_zero_tapes, Turing.Cfg.inputSymbol, Turing.FinTM.scanCfg, Turing.FinTM.scanCfg_read, Turing.MultiTapeTM, Turing.MultiTapeTM.runFrom, Turing.MultiTapeTM.runFrom_succ_eq_step', Turing.MultiTapeTM.step, Turing.moveInputPos_pos_of_ne_right}
\difficulty{4}
Suppose a control state $q$ of a machine without work tapes emits each scanned
bit and moves right, and halts silently on the blank. Then for every
$j \le \abs{x}$, $j$ steps from the scan configuration at position $0$ with
prior output $o$ lead to the scan configuration at position $j$ in state $q$
with output $o$ followed by the first $j$ bits of $x$.
\end{lemma}

\begin{lemma}[Completion of a copy state]
\lean{Turing.FinTM.scanCopy_finish}
\label{Turing.FinTM.scanCopy_finish}
\leanok
\uses{Turing.Action.apply, Turing.Cfg.ext_zero_tapes, Turing.Cfg.inputSymbol, Turing.FinTM.scanCfg, Turing.FinTM.scanCfg_read, Turing.FinTM.scanCopy_run, Turing.MultiTapeTM, Turing.MultiTapeTM.runFrom, Turing.MultiTapeTM.runFrom_succ_eq_step', Turing.MultiTapeTM.step}
\difficulty{3}
Suppose a control state $q$ of a machine without work tapes emits each scanned
bit and moves right, and halts silently on the blank. Then $\abs{x}+1$ steps
from the scan configuration at position $0$ with prior output $o$ reach a
halted configuration with output $o$ followed by all of $x$.
\end{lemma}

\begin{definition}[Input duplicator machine]
\lean{Turing.FinTM.pairDupTM}
\label{Turing.FinTM.pairDupTM}
\leanok
\uses{Turing.FinTM, Turing.FinTM.controlAction}
A five-state machine without work tapes that transforms its input $x$ into the
self-delimiting pair code of $(x,x)$: on a first pass it emits every input bit
twice, then it emits the separator bit $0$, rewinds to the start, emits the
separator bit $1$, and copies $x$ once more.
\end{definition}

\begin{lemma}[Doubling pass of the duplicator]
\lean{Turing.FinTM.pairDup_double}
\label{Turing.FinTM.pairDup_double}
\leanok
\uses{Turing.Action.apply, Turing.Cfg.ext_zero_tapes, Turing.FinTM.pairDupTM, Turing.FinTM.scanCfg, Turing.FinTM.scanCfg_read, Turing.MultiTapeTM.initCfg, Turing.MultiTapeTM.runFrom, Turing.MultiTapeTM.runFrom_succ_eq_step', Turing.MultiTapeTM.step, Turing.moveInputPos_pos_of_ne_right}
\difficulty{4}
For every $j \le \abs{x}$, after $2j$ steps from its initial configuration on
input $x$ the duplicator is in its scanning state at input position $j$ and
has emitted each of the first $j$ bits of $x$ twice.
\end{lemma}

\begin{lemma}[Time bound for duplication]
\lean{Turing.FinTM.pairDup_computes}
\label{Turing.FinTM.pairDup_computes}
\leanok
\uses{Turing.Action.apply, Turing.Cfg, Turing.Cfg.ext_zero_tapes, Turing.FinTM.ComputesInTime, Turing.FinTM.computesInTime_iff, Turing.FinTM.pairDupTM, Turing.FinTM.pairDup_double, Turing.FinTM.rewind_scan, Turing.FinTM.scanCfg, Turing.FinTM.scanCfg_read, Turing.FinTM.scanCopy_finish, Turing.MultiTapeTM.initCfg, Turing.MultiTapeTM.runFrom, Turing.MultiTapeTM.runFrom_add, Turing.MultiTapeTM.runFrom_succ_eq_step', Turing.MultiTapeTM.step, Turing.moveInputPos_neg_of_ne_left, Turing.pairEncode}
\difficulty{5}
On every input $x$ the duplicator halts within $4(\abs{x}+1)$ steps with
output the self-delimiting pair code of $(x,x)$: each bit of $x$ doubled, then
the separator bits $0,1$, then $x$.
\end{lemma}

\begin{lemma}[Suffix copy from an interior position]
\lean{Turing.FinTM.scanCopy_suffix}
\label{Turing.FinTM.scanCopy_suffix}
\leanok
\uses{Turing.Action.apply, Turing.Cfg.ext_zero_tapes, Turing.FinTM.scanCfg, Turing.FinTM.scanCfg_read, Turing.MultiTapeTM, Turing.MultiTapeTM.runFrom, Turing.MultiTapeTM.runFrom_succ_eq_step, Turing.MultiTapeTM.runFrom_zero, Turing.MultiTapeTM.step, Turing.moveInputPos_pos_of_ne_right}
\difficulty{4}
Suppose a control state $q$ of a machine without work tapes emits each scanned
bit and moves right, and halts silently on the blank. If $x = p s$ is a split
of the input into a prefix $p$ and suffix $s$, then $\abs{s}+1$ steps from the
scan configuration at position $\abs{p}$ in state $q$ with prior output $o$
reach a halted configuration with output $o$ followed by $s$.
\end{lemma}

\begin{lemma}[Scan over a run of ones]
\lean{Turing.FinTM.scanTrues_run}
\label{Turing.FinTM.scanTrues_run}
\leanok
\uses{Turing.Action.apply, Turing.Cfg.ext_zero_tapes, Turing.FinTM.scanCfg, Turing.FinTM.scanCfg_read, Turing.MultiTapeTM, Turing.MultiTapeTM.runFrom, Turing.MultiTapeTM.runFrom_succ_eq_step', Turing.MultiTapeTM.step, Turing.moveInputPos_pos_of_ne_right}
\difficulty{4}
Suppose a control state $q$ of a machine without work tapes moves right on
reading a $1$, emitting a $0$ when an emission flag is set and nothing
otherwise. If the first $j \le \abs{x}$ bits of $x$ are all $1$, then $j$
steps from the scan configuration at position $0$ in state $q$ with empty
output lead to position $j$ in state $q$, with output $j$ zeros when the flag
is set and empty output otherwise.
\end{lemma}

\begin{lemma}[Dichotomy for fixed-width increment]
\lean{Turing.FinTM.incFixed_cases}
\label{Turing.FinTM.incFixed_cases}
\leanok
\uses{Turing.incFixed}
\difficulty{4}
Every word $x \in \{0,1\}^*$, read least-significant bit first, satisfies
exactly one of the following: $x$ consists entirely of ones and its
fixed-width increment overflows (is undefined); or $x = 1^j 0 r$ for some
$j \ge 0$ and tail $r$, and the fixed-width increment of $x$ is $0^j 1 r$.
\end{lemma}

\begin{definition}[Fixed-width increment machine]
\lean{Turing.FinTM.incFixedTM}
\label{Turing.FinTM.incFixedTM}
\leanok
\uses{Turing.FinTM, Turing.FinTM.controlAction}
The fixed-width increment machine is a four-state machine without work
tapes computing the fixed-width binary increment of its input, read
least-significant bit first: it scans the leading ones silently, and on
reaching a $0$ it rewinds and re-scans, emitting a $0$ for each carried $1$,
a $1$ in place of the first $0$, and the remaining bits unchanged; an
all-ones input halts with empty output.
\end{definition}

\begin{lemma}[Time bound for fixed-width increment]
\lean{Turing.FinTM.incFixed_computes}
\label{Turing.FinTM.incFixed_computes}
\leanok
\uses{Turing.Action.apply, Turing.Cfg, Turing.Cfg.ext_zero_tapes, Turing.FinTM.ComputesInTime, Turing.FinTM.ComputesInTime.mono, Turing.FinTM.computesInTime_iff, Turing.FinTM.controlAction, Turing.FinTM.incFixedTM, Turing.FinTM.incFixed_cases, Turing.FinTM.rewind_scan, Turing.FinTM.scanCfg, Turing.FinTM.scanCfg_read, Turing.FinTM.scanCopy_suffix, Turing.FinTM.scanTrues_run, Turing.MultiTapeTM.initCfg, Turing.MultiTapeTM.runFrom, Turing.MultiTapeTM.runFrom_add, Turing.MultiTapeTM.runFrom_succ_eq_step', Turing.MultiTapeTM.step, Turing.incFixed, Turing.moveInputPos_neg_of_ne_left, Turing.moveInputPos_pos_of_ne_right}
\difficulty{5}
On every input $x$, including the empty word, the increment machine halts
within $3(\abs{x}+1)$ steps; its output is the fixed-width binary increment of
$x$ (least-significant bit first) when no overflow occurs, and is empty when
$x$ consists entirely of ones.
\end{lemma}

\begin{lemma}[Single right-moving scan step]
\lean{Turing.FinTM.scanStep_right}
\label{Turing.FinTM.scanStep_right}
\leanok
\uses{Turing.Action.apply, Turing.Cfg.ext_zero_tapes, Turing.FinTM.scanCfg, Turing.FinTM.scanCfg_read, Turing.MultiTapeTM, Turing.MultiTapeTM.step, Turing.moveInputPos_pos_of_ne_right}
\difficulty{3}
If at control state $q$ and input position $i < \abs{x}$ the transition of a
machine without work tapes moves right, optionally emits a bit $e$, and passes
to state $q'$, then one step from the scan configuration at $i$ with output
$o$ yields the scan configuration at $i+1$ in state $q'$ with output $o$
extended by $e$ when present.
\end{lemma}

\begin{definition}[Pair-code validity tester]
\lean{Turing.FinTM.pairValidTM}
\label{Turing.FinTM.pairValidTM}
\leanok
\uses{Turing.FinTM}
A machine without work tapes that checks membership in the self-delimiting
pair code: it scans its input in aligned two-bit blocks, remembering the first
bit of the current block in its finite control, continues silently while the
two bits of a block agree, and halts with a one-bit verdict at the first
disagreement or boundary blank; the verdict is $1$ exactly when the
disagreeing block is $01$.
\end{definition}

\begin{lemma}[Two steps of the validity tester]
\lean{Turing.FinTM.pairValid_block}
\label{Turing.FinTM.pairValid_block}
\leanok
\uses{Turing.FinTM.pairValidTM, Turing.FinTM.scanCfg, Turing.FinTM.scanStep_right, Turing.MultiTapeTM.runFrom, Turing.MultiTapeTM.step}
\difficulty{4}
If the input splits as $x = p\, b\, c\, r$ with bits $b, c$, then two steps of
the validity tester from its block state at position $\abs{p}$ continue
silently to the block state at position $\abs{p}+2$ when $b = c$, and
otherwise halt at position $\abs{p}+2$ with single-bit output that is $1$
exactly when $(b,c) = (0,1)$.
\end{lemma}

\begin{lemma}[Halting bound for the validity tester]
\lean{Turing.FinTM.pairValid_run}
\label{Turing.FinTM.pairValid_run}
\leanok
\uses{Turing.Action.apply, Turing.FinTM.pairValidTM, Turing.FinTM.pairValid_block, Turing.FinTM.scanCfg, Turing.FinTM.scanCfg_read, Turing.FinTM.scanStep_right, Turing.MultiTapeTM.runFrom, Turing.MultiTapeTM.runFrom_add, Turing.MultiTapeTM.runFrom_succ_eq_step, Turing.MultiTapeTM.runFrom_zero, Turing.MultiTapeTM.step, Turing.pairDecode}
\difficulty{5}
For any split $x = p s$ of the input, there is a time $t \le \abs{s}+1$ such
that after $t$ steps from the block state at position $\abs{p}$ with empty
output, the validity tester has halted and its output is the single bit
recording whether $s$ is a valid self-delimiting pair code.
\end{lemma}

\begin{lemma}[Time bound for the validity test]
\lean{Turing.FinTM.pairValid_computes}
\label{Turing.FinTM.pairValid_computes}
\leanok
\uses{Turing.Cfg.ext_zero_tapes, Turing.FinTM.ComputesInTime, Turing.FinTM.ComputesInTime.mono, Turing.FinTM.computesInTime_iff, Turing.FinTM.pairValidTM, Turing.FinTM.pairValid_run, Turing.FinTM.scanCfg, Turing.MultiTapeTM.initCfg, Turing.pairDecode}
\difficulty{3}
For every input $x$, the pair-validity tester halts on $x$ within
$\abs{x}+1$ steps, and its output is the single bit recording whether $x$
decodes under the self-delimiting pair code.
\end{lemma}

\begin{definition}[Pair component extractor]
\lean{Turing.FinTM.pairExtractTM}
\label{Turing.FinTM.pairExtractTM}
\leanok
\uses{Turing.FinTM}
A one-work-tape machine, parameterized by two Boolean flags, that parses an
input in the self-delimiting pair code of $(a,b)$: it buffers the decoded
first component $a$ on its work tape while validating the doubled prefix, and
after the separator it rewinds and replays the buffer, emitting $a$ when the
first flag is set, then scans the suffix, emitting $b$ when the second flag is
set; invalid inputs make it halt silently.
\end{definition}

\begin{definition}[Extractor configuration]
\lean{Turing.FinTM.extractCfg}
\label{Turing.FinTM.extractCfg}
\leanok
\uses{Turing.Cfg, Turing.FinTM.bufferTape}
Given an input $x$, an optional control state, an input position
$i \le \abs{x}$, a buffered word $a$, a work-head displacement $z$, and an
output word, the extractor configuration is the configuration of the pair
extractor with that control state, with its input head past the first $i$
cells, with $a$ stored on its single work tape and the work head at $z$, and
with the given output.
\end{definition}

\begin{lemma}[Input symbol of an extractor configuration]
\lean{Turing.FinTM.extractCfg_read}
\label{Turing.FinTM.extractCfg_read}
\leanok
\uses{Turing.Cfg.inputSymbol, Turing.FinTM.extractCfg, Turing.FinTM.scanCfg_read}
\difficulty{1}
In an extractor configuration at input position $i \le \abs{x}$, the symbol
under the input head is the $i$-th bit of $x$ when $i < \abs{x}$ and the blank
otherwise, regardless of the buffer contents.
\end{lemma}

\begin{lemma}[First half of an extractor block]
\lean{Turing.FinTM.extract_first}
\label{Turing.FinTM.extract_first}
\leanok
\uses{Turing.Action.apply, Turing.FinTM.extractCfg, Turing.FinTM.extractCfg_read, Turing.FinTM.pairExtractTM, Turing.MultiTapeTM.step, Turing.moveInputPos_pos_of_ne_right}
\difficulty{3}
If the input splits as $x = p\, b\, r$ with bit $b$, one step of the pair
extractor from its block-start state at position $\abs{p}$ records $b$ in the
finite control and advances the input head, leaving the buffer and the output
unchanged.
\end{lemma}

\begin{lemma}[One aligned block of the extractor]
\lean{Turing.FinTM.extract_block}
\label{Turing.FinTM.extract_block}
\leanok
\uses{Turing.Action.apply, Turing.FinTM.bufferTape_append, Turing.FinTM.extractCfg, Turing.FinTM.extractCfg_read, Turing.FinTM.extract_first, Turing.FinTM.pairExtractTM, Turing.MultiTapeTM.runFrom, Turing.MultiTapeTM.step, Turing.moveInputPos, Turing.moveInputPos_pos_of_ne_right}
\difficulty{4}
If the input splits as $x = p\, b\, c\, r$ with bits $b, c$, then two
extractor steps from the block-start state at position $\abs{p}$ with buffer
$a$ append $b$ to the buffer and continue when $b = c$; halt silently when
$(b,c) = (1,0)$; and enter the rewind phase with the work head one cell before
the end of the buffer when $(b,c) = (0,1)$. In all three cases nothing is
emitted.
\end{lemma}

\begin{lemma}[Buffer rewind of the extractor]
\lean{Turing.FinTM.extract_rewind}
\label{Turing.FinTM.extract_rewind}
\leanok
\uses{Turing.Action.apply, Turing.Cfg.workTapeSymbols, Turing.FinTM.bufferTape_left, Turing.FinTM.bufferTape_nat, Turing.FinTM.extractCfg, Turing.FinTM.pairExtractTM, Turing.MultiTapeTM.runFrom, Turing.MultiTapeTM.runFrom_succ_eq_step, Turing.MultiTapeTM.runFrom_zero, Turing.MultiTapeTM.step, Turing.moveInputPos_zero}
\difficulty{4}
From the extractor's rewind state with buffer $a$ and work head at
displacement $j-1$, for any $j \le \abs{a}$, exactly $j+1$ steps return the
work head to displacement $0$ and enter the replay state, with the input
position, the buffer, and the (empty) output unchanged.
\end{lemma}

\begin{lemma}[Buffer replay of the extractor]
\lean{Turing.FinTM.extract_replay}
\label{Turing.FinTM.extract_replay}
\leanok
\uses{Turing.Action.apply, Turing.Cfg.workTapeSymbols, Turing.FinTM.bufferTape_nat, Turing.FinTM.extractCfg, Turing.FinTM.pairExtractTM, Turing.MultiTapeTM.runFrom, Turing.MultiTapeTM.runFrom_succ_eq_step', Turing.MultiTapeTM.step, Turing.moveInputPos_zero}
\difficulty{4}
From the extractor's replay state with buffer $a$ and work head at
displacement $0$, for any $j \le \abs{a}$, exactly $j$ steps move the work
head to displacement $j$ and emit the first $j$ bits of $a$ when the
first-component flag is set, and nothing otherwise.
\end{lemma}

\begin{lemma}[End of buffer replay]
\lean{Turing.FinTM.extract_replay_finish}
\label{Turing.FinTM.extract_replay_finish}
\leanok
\uses{Turing.Action.apply, Turing.Cfg.workTapeSymbols, Turing.FinTM.bufferTape_nat, Turing.FinTM.extractCfg, Turing.FinTM.extract_replay, Turing.FinTM.pairExtractTM, Turing.MultiTapeTM.runFrom, Turing.MultiTapeTM.runFrom_succ_eq_step', Turing.MultiTapeTM.step, Turing.moveInputPos_zero}
\difficulty{3}
From the extractor's replay state with buffer $a$ and work head at
displacement $0$, exactly $\abs{a}+1$ steps reach the suffix state with the
work head at displacement $\abs{a}$, and the output is $a$ when the
first-component flag is set and empty otherwise.
\end{lemma}

\begin{lemma}[Suffix copy of the extractor]
\lean{Turing.FinTM.extract_suffix}
\label{Turing.FinTM.extract_suffix}
\leanok
\uses{Turing.Action.apply, Turing.FinTM.extractCfg, Turing.FinTM.extractCfg_read, Turing.FinTM.pairExtractTM, Turing.MultiTapeTM.runFrom, Turing.MultiTapeTM.runFrom_succ_eq_step, Turing.MultiTapeTM.runFrom_zero, Turing.MultiTapeTM.step, Turing.moveInputPos_pos_of_ne_right}
\difficulty{4}
With the second-component flag set and a split $x = p s$ of the input, from
the extractor's suffix state at position $\abs{p}$ with buffer $a$ and prior
output $o$, exactly $\abs{s}+1$ steps reach a halted configuration with output
$o$ followed by $s$; the buffer and its head are untouched.
\end{lemma}

\begin{lemma}[Post-validation ledger of the extractor]
\lean{Turing.FinTM.extract_finish}
\label{Turing.FinTM.extract_finish}
\leanok
\uses{Turing.Action.apply, Turing.FinTM.extractCfg, Turing.FinTM.extract_replay_finish, Turing.FinTM.extract_rewind, Turing.FinTM.extract_suffix, Turing.FinTM.pairExtractTM, Turing.MultiTapeTM.runFrom, Turing.MultiTapeTM.runFrom_add, Turing.MultiTapeTM.runFrom_succ_eq_step', Turing.MultiTapeTM.step}
\difficulty{4}
For a split $x = p s$ of the input and a validated buffer $a$, there is a time
$t \le 2\abs{a} + \abs{s} + 3$ such that after $t$ steps from the start of the
rewind phase at position $\abs{p}$ the extractor has halted, and its output is
$a$ (when the first flag is set) followed by $s$ (when the second flag is
set).
\end{lemma}

\begin{lemma}[Main run invariant of the extractor]
\lean{Turing.FinTM.extract_run}
\label{Turing.FinTM.extract_run}
\leanok
\uses{Turing.Action.apply, Turing.FinTM.extractCfg, Turing.FinTM.extractCfg_read, Turing.FinTM.extract_block, Turing.FinTM.extract_finish, Turing.FinTM.extract_first, Turing.FinTM.pairExtractTM, Turing.MultiTapeTM.runFrom, Turing.MultiTapeTM.runFrom_add, Turing.MultiTapeTM.runFrom_succ_eq_step, Turing.MultiTapeTM.runFrom_zero, Turing.MultiTapeTM.step, Turing.pairDecode}
\difficulty{5}
For any split $x = p s$ of the input and any buffered prefix $a$, there is a
time $t \le 3\abs{s} + 2\abs{a} + 5$ such that after $t$ steps from the
block-start state at position $\abs{p}$ the extractor has halted; if $s$
decodes under the self-delimiting pair code as $(b,c)$, the output is $a b$
(when the first flag is set) followed by $c$ (when the second flag is set),
and if $s$ is not a valid pair code the output is empty.
\end{lemma}

\begin{lemma}[Time bound for pair extraction]
\lean{Turing.FinTM.pairExtract_computes}
\label{Turing.FinTM.pairExtract_computes}
\leanok
\uses{Turing.Cfg.init, Turing.FinTM.ComputesInTime, Turing.FinTM.ComputesInTime.mono, Turing.FinTM.computesInTime_iff, Turing.FinTM.extractCfg, Turing.FinTM.extract_run, Turing.FinTM.pairExtractTM, Turing.MultiTapeTM.initCfg, Turing.pairDecode}
\difficulty{3}
On every input $x$ the pair extractor halts within $5(\abs{x}+1)$ steps; if
$x$ decodes under the self-delimiting pair code as $(a,b)$, the output is $a$
(when the first flag is set) followed by $b$ (when the second flag is set),
and if $x$ is not a valid pair code the output is empty.
\end{lemma}

\begin{definition}[Control states of the polynomial generator]
\lean{Turing.FinTM.CatalogPolyControl}
\label{Turing.FinTM.CatalogPolyControl}
\leanok
For parameters $c$ and $C$, the control-state type of the unary polynomial
generator has a copy state, a setup state, a loop state and a rewind state
for each of the $c+1$ loop levels, an advance state for each of $c+2$
indices (one per level and one past the last), and an emission state for
each remaining count from $0$ to $C$.
\end{definition}

\begin{definition}[Unary tape of a given length]
\lean{Turing.FinTM.catalogPolyTape}
\label{Turing.FinTM.catalogPolyTape}
\leanok
The tape content holding the unary word of length $q$: cell $z$ contains a $1$
for $0 \le z < q$ and is blank otherwise.
\end{definition}

\begin{definition}[Head-only move action]
\lean{Turing.FinTM.catalogPolyMove}
\label{Turing.FinTM.catalogPolyMove}
\leanok
\uses{Turing.Action, Turing.FinTM.CatalogPolyControl}
Given a work-tape index $i$, a direction $d$, and a control state $s$, the
head-only move action is the action that moves work head $i$ one cell in
direction $d$, keeps the input head and all other work heads in place,
writes no tape cell, emits no output, and enters $s$.
\end{definition}

\begin{definition}[Unary polynomial generator machine]
\lean{Turing.FinTM.catalogPolyUnaryTM}
\label{Turing.FinTM.catalogPolyUnaryTM}
\leanok
\uses{Turing.FinTM, Turing.FinTM.CatalogPolyControl, Turing.FinTM.catalogPolyMove}
A machine with $c+1$ work tapes that writes a run of $C\,(n+1)^{c+1}$ ones on
its output tape, where $n$ is the input length: in one input scan it copies
$n+1$ in unary onto every work tape in parallel, and then runs $c+1$ nested
loops, each of side $n+1$, emitting $C$ ones at each visit of the innermost
body; a completed inner loop is rewound at a cost charged to the iterations
that just completed.
\end{definition}

\begin{definition}[Generator loop configuration]
\lean{Turing.FinTM.catalogPolyCfg}
\label{Turing.FinTM.catalogPolyCfg}
\leanok
\uses{Turing.Cfg, Turing.FinTM.CatalogPolyControl, Turing.FinTM.catalogPolyTape}
A configuration of the unary polynomial generator in which every work tape
holds the unary word of length $q$, with a given control state, arbitrary
work-head positions, and accumulated output.
\end{definition}

\begin{lemma}[Effect of a head-only move]
\lean{Turing.FinTM.catalogPolyMove_apply}
\label{Turing.FinTM.catalogPolyMove_apply}
\leanok
\uses{Turing.Action.apply, Turing.FinTM.CatalogPolyControl, Turing.FinTM.catalogPolyCfg, Turing.FinTM.catalogPolyMove, Turing.moveInputPos_zero}
\difficulty{3}
For every work-tape index $i$ and direction $d$, applying the head-only move
action for $i$, $d$, and a target state to a generator loop configuration
yields the loop configuration in that target state whose $i$-th work-head
position is increased by $d$; all tapes, the input position, the other
heads, and the output are unchanged.
\end{lemma}

\begin{lemma}[Emission chain of the generator]
\lean{Turing.FinTM.catalogPoly_emit}
\label{Turing.FinTM.catalogPoly_emit}
\leanok
\uses{Turing.Action.apply, Turing.FinTM.catalogPolyCfg, Turing.FinTM.catalogPolyUnaryTM, Turing.MultiTapeTM.runFrom, Turing.MultiTapeTM.runFrom_succ_eq_step, Turing.MultiTapeTM.runFrom_zero, Turing.MultiTapeTM.step}
\difficulty{4}
From the generator's emission state with counter $j \le C$, exactly $j+1$
steps append $j$ ones to the output and enter the innermost advance state,
leaving every head position unchanged.
\end{lemma}

\begin{lemma}[Rewind of one loop head]
\lean{Turing.FinTM.catalogPoly_rewind}
\label{Turing.FinTM.catalogPoly_rewind}
\leanok
\uses{Turing.Action, Turing.Action.apply, Turing.Cfg.workTapeSymbols, Turing.FinTM.CatalogPolyControl, Turing.FinTM.catalogPolyCfg, Turing.FinTM.catalogPolyMove_apply, Turing.FinTM.catalogPolyTape, Turing.FinTM.catalogPolyUnaryTM, Turing.MultiTapeTM.runFrom, Turing.MultiTapeTM.runFrom_succ_eq_step, Turing.MultiTapeTM.runFrom_zero, Turing.MultiTapeTM.step}
\difficulty{4}
From the rewind state of loop level $i$ with head $i$ at displacement $j-1$,
for any $j \le q$, exactly $j+1$ steps reset head $i$ to displacement $0$ and
enter the advance state of level $i+1$; the other loop heads and the
accumulated output are unchanged.
\end{lemma}

\begin{lemma}[Advance into a loop level]
\lean{Turing.FinTM.catalogPoly_advance}
\label{Turing.FinTM.catalogPoly_advance}
\leanok
\uses{Turing.FinTM.catalogPolyCfg, Turing.FinTM.catalogPolyMove_apply, Turing.FinTM.catalogPolyUnaryTM, Turing.MultiTapeTM.step}
\difficulty{2}
For every loop level $i$, one step of the generator from the advance state
of level $i$ yields the loop state of level $i$ with head $i$ moved one cell
to the right, everything else being unchanged.
\end{lemma}

\begin{definition}[Exact cost of the loop nest]
\lean{Turing.FinTM.catalogPolyCost}
\label{Turing.FinTM.catalogPolyCost}
\leanok
For a side length $q$, an emission count $C$, and a number $r$ of loop
levels, the exact running time $T_r$ of a full nest of unary loops is
defined by $T_0 = C + 1$ and $T_{r+1} = q\,(T_r + 2) + q + 2$.
\end{definition}

\begin{lemma}[Main loop invariant of the generator]
\lean{Turing.FinTM.catalogPoly_loop}
\label{Turing.FinTM.catalogPoly_loop}
\leanok
\uses{Turing.Action.apply, Turing.Cfg.workTapeSymbols, Turing.FinTM.catalogPolyCfg, Turing.FinTM.catalogPolyCost, Turing.FinTM.catalogPolyMove_apply, Turing.FinTM.catalogPolyTape, Turing.FinTM.catalogPolyUnaryTM, Turing.FinTM.catalogPoly_advance, Turing.FinTM.catalogPoly_emit, Turing.FinTM.catalogPoly_rewind, Turing.MultiTapeTM.runFrom, Turing.MultiTapeTM.runFrom_add, Turing.MultiTapeTM.runFrom_succ_eq_step, Turing.MultiTapeTM.runFrom_zero, Turing.MultiTapeTM.step}
\difficulty{6}
Let $q > 0$ be the unary side length, let $i \le c$ be a loop level, and
suppose every head at a level $\le i$ is at displacement $0$ except that head
$i$ sits at $j$, where $j + r = q$. Then exactly $r\,(T_i + 2) + q + 2$ steps
(with $T_i$ the exact level-$i$ nest cost) execute the remaining $r$
iterations of loop $i$, append $r \cdot C q^i$ ones to the output, reset head
$i$ to $0$, and enter the advance state of level $i+1$; heads above level $i$
are arbitrary and remain unchanged.
\end{lemma}

\begin{lemma}[Extending a unary tape]
\lean{Turing.FinTM.catalogPolyTape_write}
\label{Turing.FinTM.catalogPolyTape_write}
\leanok
\uses{Turing.FinTM.catalogPolyTape}
\difficulty{3}
For every $q$, writing a $1$ at cell $q$ of the unary tape of length $q$
yields the unary tape of length $q+1$.
\end{lemma}

\begin{lemma}[Linear bound on the loop cost]
\lean{Turing.FinTM.catalogPolyCost_le}
\label{Turing.FinTM.catalogPolyCost_le}
\leanok
\uses{Turing.FinTM.catalogPolyCost}
\difficulty{4}
For every side length $q > 0$, emission count $C$, and number of levels $r$,
the exact nest cost satisfies $T_r \le (C + 1 + 5r)\, q^r$.
\end{lemma}

\begin{definition}[Copy-phase configuration of the generator]
\lean{Turing.FinTM.catalogPolyCopyCfg}
\label{Turing.FinTM.catalogPolyCopyCfg}
\leanok
\uses{Turing.Cfg, Turing.FinTM.CatalogPolyControl, Turing.FinTM.catalogPolyTape}
The generator's configuration after scanning $i \le \abs{x}$ input cells in
its copy phase: every work tape holds the unary word of length $i$ and the
input head has passed $i$ cells.
\end{definition}

\begin{lemma}[Unary copy of the input length]
\lean{Turing.FinTM.catalogPoly_copy}
\label{Turing.FinTM.catalogPoly_copy}
\leanok
\uses{Turing.Action.apply, Turing.Cfg.init, Turing.Cfg.inputSymbol, Turing.FinTM.catalogPolyCopyCfg, Turing.FinTM.catalogPolyTape, Turing.FinTM.catalogPolyTape_write, Turing.FinTM.catalogPolyUnaryTM, Turing.MultiTapeTM.initCfg, Turing.MultiTapeTM.runFrom, Turing.MultiTapeTM.runFrom_succ_eq_step', Turing.MultiTapeTM.step, Turing.inputSymbolInner, Turing.moveInputPos, Turing.moveInputPos_pos_of_ne_right}
\difficulty{4}
For every $i \le \abs{x}$, after $i$ steps from the generator's initial
configuration on input $x$ the machine is in the copy-phase configuration
with $i$ unary marks on every work tape.
\end{lemma}

\begin{lemma}[Setup rewind of the generator]
\lean{Turing.FinTM.catalogPoly_setup}
\label{Turing.FinTM.catalogPoly_setup}
\leanok
\uses{Turing.Action.apply, Turing.Cfg.workTapeSymbols, Turing.FinTM.catalogPolyCfg, Turing.FinTM.catalogPolyTape, Turing.FinTM.catalogPolyUnaryTM, Turing.MultiTapeTM.runFrom, Turing.MultiTapeTM.runFrom_succ_eq_step, Turing.MultiTapeTM.runFrom_zero, Turing.MultiTapeTM.step}
\difficulty{4}
From the generator's setup state with all work heads at displacement $j-1$,
for any $j \le q$, exactly $j+1$ steps move all heads to displacement $0$ in
lockstep and enter the outermost loop state, with empty output throughout.
\end{lemma}

\begin{lemma}[Startup of the generator]
\lean{Turing.FinTM.catalogPoly_start}
\label{Turing.FinTM.catalogPoly_start}
\leanok
\uses{Turing.Action.apply, Turing.Cfg.inputSymbol, Turing.FinTM.catalogPolyCfg, Turing.FinTM.catalogPolyCopyCfg, Turing.FinTM.catalogPolyTape_write, Turing.FinTM.catalogPolyUnaryTM, Turing.FinTM.catalogPoly_copy, Turing.FinTM.catalogPoly_setup, Turing.MultiTapeTM.initCfg, Turing.MultiTapeTM.runFrom, Turing.MultiTapeTM.runFrom_add, Turing.MultiTapeTM.runFrom_succ_eq_step', Turing.MultiTapeTM.step, Turing.moveInputPos_zero}
\difficulty{4}
For every input $x$, running the generator for exactly $2(\abs{x}+1)$ steps
from its initial configuration yields the configuration in which every work
tape holds the unary word of length $\abs{x}+1$, the control is in the
outermost loop state, all heads are at displacement $0$, and the output is
empty; the extra unary cell makes empty input need no special case.
\end{lemma}

\begin{lemma}[Time bound for the unary generator]
\lean{Turing.FinTM.catalogPoly_unary_computes}
\label{Turing.FinTM.catalogPoly_unary_computes}
\leanok
\uses{Turing.Action.apply, Turing.FinTM.ComputesFunInTime, Turing.FinTM.ComputesInTime, Turing.FinTM.ComputesInTime.mono, Turing.FinTM.catalogPolyCfg, Turing.FinTM.catalogPolyCost, Turing.FinTM.catalogPolyCost_le, Turing.FinTM.catalogPolyUnaryTM, Turing.FinTM.catalogPoly_loop, Turing.FinTM.catalogPoly_start, Turing.FinTM.computesInTime_iff, Turing.MultiTapeTM.runFrom, Turing.MultiTapeTM.runFrom_add, Turing.MultiTapeTM.runFrom_succ_eq_step', Turing.MultiTapeTM.step}
\difficulty{5}
For all parameters $c$ and $C$, the unary polynomial generator computes the
function sending a word of length $n$ to the run of $C\,(n+1)^{c+1}$ ones,
within $(C + 5(c+1) + 4)\,(n+1)^{c+1}$ steps; this covers coefficient $C = 0$
and empty input.
\end{lemma}

\begin{definition}[Administrative action of the length checker]
\lean{Turing.FinTM.lenAction}
\label{Turing.FinTM.lenAction}
\leanok
\uses{Turing.Action, Turing.FinTM}
For a machine $M$, the administrative action of the length checker built
over $M$, given an input direction, a countdown direction, an optional bit,
and a target state, is the action that moves the input head and the final
(countdown) work head by those directions, emits the optional bit, passes to
the target state, and writes no tape cell.
\end{definition}

\begin{definition}[Captured-length checker machine]
\lean{Turing.FinTM.pairCountTM}
\label{Turing.FinTM.pairCountTM}
\leanok
\uses{Turing.FinTM, Turing.FinTM.controlAction, Turing.FinTM.lenAction, Turing.captureAction}
Given a machine $M$, the captured-length checker over $M$ is a machine with
one additional work tape that first simulates $M$ on the input while
capturing the output of $M$ on the extra tape, rewinds the input head to the
start, parses the input's doubled prefix under the self-delimiting pair
grammar, and after the separator compares the length of the remaining suffix
against the length of the captured word, emitting a single Boolean verdict;
only the final comparison or a rejection emits output.
\end{definition}

\begin{definition}[Length-checker configuration]
\lean{Turing.FinTM.lenCfg}
\label{Turing.FinTM.lenCfg}
\leanok
\uses{Turing.Cfg, Turing.FinTM, Turing.FinTM.pairCountTM, Turing.captureCfg}
Given a machine $M$, an input $x$, a configuration $c$ of $M$ on $x$, a
checker control state, an input position $i \le \abs{x}$, and a counter
$r$, the length-checker configuration retains the tape bank of $c$ together
with the output of $c$ captured on the extra tape, has the given control
state and its input head past the first $i$ cells, and places the capture
head at displacement $r - 1$, so that $r$ countdown cells remain available.
\end{definition}

\begin{lemma}[Input symbol of a checker configuration]
\lean{Turing.FinTM.lenCfg_read}
\label{Turing.FinTM.lenCfg_read}
\leanok
\uses{Turing.Cfg, Turing.Cfg.inputSymbol, Turing.FinTM, Turing.FinTM.inputSymbol_at, Turing.FinTM.lenCfg, Turing.FinTM.pairCountTM}
\difficulty{1}
For every input $x$ and every position $i \le \abs{x}$, in a length-checker
configuration at input position $i$ the input head reads the $i$-th bit of
$x$ (the blank when $i = \abs{x}$), independently of the saved tape bank.
\end{lemma}

\begin{lemma}[Effect of an administrative action]
\lean{Turing.FinTM.lenAction_apply}
\label{Turing.FinTM.lenAction_apply}
\leanok
\uses{Turing.Action.apply, Turing.Cfg, Turing.FinTM, Turing.FinTM.lenAction, Turing.FinTM.lenCfg, Turing.FinTM.pairCountTM, Turing.moveInputPos}
\difficulty{3}
Applying an administrative action with compatible input-head and
countdown-head moves to a length-checker configuration yields the checker
configuration with the correspondingly updated input position and countdown
counter, with output equal to the action's optional emission and all work
tapes unchanged.
\end{lemma}

\begin{lemma}[Suffix length comparison]
\lean{Turing.FinTM.lenSuffix_run}
\label{Turing.FinTM.lenSuffix_run}
\leanok
\uses{Turing.Action.apply, Turing.Cfg, Turing.Cfg.workTapeSymbols, Turing.FinTM, Turing.FinTM.bufferTape, Turing.FinTM.bufferTape_left, Turing.FinTM.lenAction, Turing.FinTM.lenAction_apply, Turing.FinTM.lenCfg, Turing.FinTM.lenCfg_read, Turing.FinTM.pairCountTM, Turing.MultiTapeTM.runFrom, Turing.MultiTapeTM.runFrom_add, Turing.MultiTapeTM.runFrom_succ_eq_step, Turing.MultiTapeTM.runFrom_zero, Turing.MultiTapeTM.step, Turing.captureCfg, Turing.moveInputPos_pos_of_ne_right}
\difficulty{5}
For a split $x = p s$ of the input and a counter $r$ at most the captured
word's length, from the checker's comparison state at position $\abs{p}$
there is a time $t \le \abs{s} + 1$ such that after $t$ steps the checker has
halted with output the single bit recording whether $\abs{s} \le r$; the
empty suffix succeeds even with counter $0$.
\end{lemma}

\begin{lemma}[First half of a checker block]
\lean{Turing.FinTM.lenParse_first}
\label{Turing.FinTM.lenParse_first}
\leanok
\uses{Turing.Action.apply, Turing.Cfg, Turing.FinTM, Turing.FinTM.lenAction_apply, Turing.FinTM.lenCfg, Turing.FinTM.lenCfg_read, Turing.FinTM.pairCountTM, Turing.MultiTapeTM.step, Turing.moveInputPos_pos_of_ne_right}
\difficulty{3}
If the input splits as $x = p\, b\, s$ with bit $b$, one checker step from
the block-start parse state at position $\abs{p}$ records $b$ in the finite
control and advances the input head; the countdown cells and the output are
untouched.
\end{lemma}

\begin{lemma}[One aligned block of the checker]
\lean{Turing.FinTM.lenParse_block}
\label{Turing.FinTM.lenParse_block}
\leanok
\uses{Turing.Action.apply, Turing.Cfg, Turing.FinTM, Turing.FinTM.lenAction_apply, Turing.FinTM.lenCfg, Turing.FinTM.lenCfg_read, Turing.FinTM.lenParse_first, Turing.FinTM.pairCountTM, Turing.MultiTapeTM.runFrom, Turing.MultiTapeTM.step, Turing.moveInputPos, Turing.moveInputPos_pos_of_ne_right, Turing.moveInputPos_zero}
\difficulty{4}
If the input splits as $x = p\, b\, d\, s$ with bits $b, d$, two checker
steps from the block-start parse state at position $\abs{p}$ continue
silently at position $\abs{p}+2$ when $b = d$; halt with verdict $0$ when
$(b,d) = (1,0)$; and enter the suffix length comparison when $(b,d) = (0,1)$.
The countdown counter is unchanged in every case.
\end{lemma}

\begin{lemma}[Parse-and-compare run of the checker]
\lean{Turing.FinTM.lenParse_run}
\label{Turing.FinTM.lenParse_run}
\leanok
\uses{Turing.Action.apply, Turing.Cfg, Turing.FinTM, Turing.FinTM.lenAction, Turing.FinTM.lenCfg, Turing.FinTM.lenCfg_read, Turing.FinTM.lenParse_block, Turing.FinTM.lenParse_first, Turing.FinTM.lenSuffix_run, Turing.FinTM.pairCountTM, Turing.MultiTapeTM.runFrom, Turing.MultiTapeTM.runFrom_add, Turing.MultiTapeTM.runFrom_succ_eq_step, Turing.MultiTapeTM.runFrom_zero, Turing.MultiTapeTM.step, Turing.captureCfg, Turing.pairDecode}
\difficulty{5}
For a split $x = p s$ of the input and a counter $r$ at most the captured
word's length, from the checker's block-start parse state at position
$\abs{p}$ there is a time $t \le \abs{s} + 1$ after which the checker has
halted with a one-bit output: if $s$ decodes under the self-delimiting pair
code as $(a,b)$, the bit records whether $\abs{b} \le r$, and if $s$ is
malformed the bit is $0$.
\end{lemma}

\begin{lemma}[Quantitative input rewind]
\lean{Turing.FinTM.catalogRewind}
\label{Turing.FinTM.catalogRewind}
\leanok
\uses{Turing.Cfg, Turing.FinTM.controlAction, Turing.FinTM.controlAction_apply, Turing.FinTM.moveInputPos_neg_val, Turing.FinTM.rewind_scan, Turing.MultiTapeTM, Turing.MultiTapeTM.runFrom, Turing.MultiTapeTM.runFrom_add, Turing.MultiTapeTM.step, Turing.moveInputPos}
\difficulty{4}
Suppose a machine has a start state that unconditionally moves the input head
left into a scanning state, and the scanning state keeps moving left on
symbols and exits right into a destination state on the left blank. Then from
any configuration $c$ in the start state there is a time
$r \le i + 2$, where $i$ is the current input position, after which the
machine is in the destination state with the input head back at the first
cell, all tapes and the output unchanged.
\end{lemma}

\begin{lemma}[Capture-and-rewind startup of the checker]
\lean{Turing.FinTM.lenStart}
\label{Turing.FinTM.lenStart}
\leanok
\uses{Turing.Action.apply, Turing.Cfg, Turing.Cfg.init, Turing.FinTM, Turing.FinTM.ComputesInTime, Turing.FinTM.ComputesInTime.output_unique, Turing.FinTM.catalogRewind, Turing.FinTM.computesInTime_iff, Turing.FinTM.lenAction, Turing.FinTM.lenCfg, Turing.FinTM.pairCountTM, Turing.MultiTapeTM.initCfg, Turing.MultiTapeTM.runFrom, Turing.MultiTapeTM.runFrom_add, Turing.MultiTapeTM.runFrom_succ_eq_step, Turing.MultiTapeTM.runFrom_zero, Turing.MultiTapeTM.step, Turing.captureCfg, Turing.capture_run, Turing.moveInputPos_zero}
\difficulty{5}
Let $M$ be a machine that, on input $x$, halts within $T$ steps with output
$w$. Then within $T + \abs{x} + 4$ steps the length checker built over $M$
reaches, without emitting any physical output, the checker configuration at
input position $0$ in its parse state, retaining a halted bank of $M$ whose
captured output is $w$ and with all $\abs{w}$ countdown cells available.
\end{lemma}

\begin{lemma}[Time bound for the captured-length checker]
\lean{Turing.FinTM.pairCount_computes}
\label{Turing.FinTM.pairCount_computes}
\leanok
\uses{Turing.FinTM, Turing.FinTM.ComputesFunInTime, Turing.FinTM.ComputesInTime, Turing.FinTM.ComputesInTime.mono, Turing.FinTM.computesInTime_iff, Turing.FinTM.lenParse_run, Turing.FinTM.lenStart, Turing.FinTM.pairCountTM, Turing.MultiTapeTM.runFrom_add, Turing.pairDecode}
\difficulty{3}
Let $M$ compute the total function $g$ within time $T$. Then the length
checker built over $M$ computes, within $T(n) + 2n + 5$ steps on inputs of
length $n$, the one-bit function that on input $x$ outputs whether
$\abs{b} \le \abs{g(x)}$ when $x$ decodes under the self-delimiting pair code
as $(a,b)$, and outputs $0$ when $x$ is malformed.
\end{lemma}

\begin{definition}[Raw marker-stripping machine]
\lean{Turing.FinTM.rawStripTM}
\label{Turing.FinTM.rawStripTM}
\leanok
\uses{Turing.FinTM}
A four-state, one-work-tape machine that copies its input to the work tape,
then erases the trailing run of zeros and the last $1$ from the right, rewinds,
and replays the retained prefix to the output; an input containing no $1$
halts silently during the reverse scan.
\end{definition}

\begin{definition}[Stripping-machine configuration]
\lean{Turing.FinTM.stripCfg}
\label{Turing.FinTM.stripCfg}
\leanok
\uses{Turing.Cfg, Turing.FinTM.bufferTape}
Given an input $x$, an optional control state, an input position
$i \le \abs{x}$, a buffer word $w$, a head displacement $h$, and an output
word, the stripping-machine configuration is the configuration of the
marker-stripping machine with that control state, with its input head past
the first $i$ cells, with $w$ stored contiguously on its work tape and the
work head at $h$, and with the given output.
\end{definition}

\begin{lemma}[Erasing the last buffer cell]
\lean{Turing.FinTM.catalogBuffer_erase}
\label{Turing.FinTM.catalogBuffer_erase}
\leanok
\uses{Turing.FinTM.bufferTape, Turing.FinTM.bufferTape_append}
\difficulty{3}
Blanking the cell at position $\abs{w}$ of the tape holding the buffered word
$w b$ (for any final bit $b$) yields exactly the tape holding the buffered
word $w$.
\end{lemma}

\begin{lemma}[Forward copy of the stripping machine]
\lean{Turing.FinTM.rawStrip_copy}
\label{Turing.FinTM.rawStrip_copy}
\leanok
\uses{Turing.Action.apply, Turing.Cfg.inputSymbol, Turing.FinTM.bufferTape, Turing.FinTM.bufferTape_append, Turing.FinTM.rawStripTM, Turing.FinTM.stripCfg, Turing.MultiTapeTM.initCfg, Turing.MultiTapeTM.runFrom, Turing.MultiTapeTM.runFrom_succ_eq_step', Turing.MultiTapeTM.step, Turing.inputSymbolInner, Turing.moveInputPos_pos_of_ne_right}
\difficulty{4}
For every $j \le \abs{x}$, after $j$ steps from its initial configuration on
input $x$ the stripping machine holds the first $j$ bits of $x$ on its work
tape with both heads advanced by $j$, and has emitted nothing.
\end{lemma}

\begin{lemma}[Rewind of the stripping machine]
\lean{Turing.FinTM.rawStrip_rewind}
\label{Turing.FinTM.rawStrip_rewind}
\leanok
\uses{Turing.Action.apply, Turing.Cfg.workTapeSymbols, Turing.FinTM.bufferTape_left, Turing.FinTM.bufferTape_nat, Turing.FinTM.rawStripTM, Turing.FinTM.stripCfg, Turing.MultiTapeTM.runFrom, Turing.MultiTapeTM.runFrom_succ_eq_step, Turing.MultiTapeTM.runFrom_zero, Turing.MultiTapeTM.step, Turing.moveInputPos_zero}
\difficulty{4}
From the stripping machine's rewind state with buffer $a$ and work head at
displacement $j-1$, for any $j \le \abs{a}$, exactly $j+1$ steps return the
work head to displacement $0$ and enter the replay state, emitting nothing.
\end{lemma}

\begin{lemma}[Replay of the stripping machine]
\lean{Turing.FinTM.rawStrip_replay}
\label{Turing.FinTM.rawStrip_replay}
\leanok
\uses{Turing.Action.apply, Turing.Cfg.workTapeSymbols, Turing.FinTM.bufferTape_nat, Turing.FinTM.rawStripTM, Turing.FinTM.stripCfg, Turing.MultiTapeTM.runFrom, Turing.MultiTapeTM.runFrom_succ_eq_step', Turing.MultiTapeTM.step, Turing.moveInputPos_zero}
\difficulty{4}
From the stripping machine's replay state with buffer $a$ and work head at
displacement $0$, for any $j \le \abs{a}$, exactly $j$ steps move the work
head to displacement $j$ and emit the first $j$ bits of $a$, preserving the
buffer.
\end{lemma}

\begin{lemma}[Rewind-replay completion of the stripping machine]
\lean{Turing.FinTM.rawStrip_finish}
\label{Turing.FinTM.rawStrip_finish}
\leanok
\uses{Turing.Action.apply, Turing.Cfg.workTapeSymbols, Turing.FinTM.rawStripTM, Turing.FinTM.rawStrip_replay, Turing.FinTM.rawStrip_rewind, Turing.FinTM.stripCfg, Turing.MultiTapeTM.runFrom, Turing.MultiTapeTM.runFrom_add, Turing.MultiTapeTM.runFrom_succ_eq_step', Turing.MultiTapeTM.step}
\difficulty{3}
From the stripping machine's rewind state with retained buffer $a$ and work
head at displacement $\abs{a}-1$, exactly $2(\abs{a}+1)$ steps reach a halted
configuration whose output is exactly $a$.
\end{lemma}

\begin{lemma}[Reverse erasure step]
\lean{Turing.FinTM.rawStrip_erase}
\label{Turing.FinTM.rawStrip_erase}
\leanok
\uses{Turing.Action.apply, Turing.Cfg.workTapeSymbols, Turing.FinTM.bufferTape_nat, Turing.FinTM.catalogBuffer_erase, Turing.FinTM.rawStripTM, Turing.FinTM.stripCfg, Turing.MultiTapeTM.step, Turing.moveInputPos_zero}
\difficulty{3}
One step of the stripping machine in its reverse state with buffer $w b$ and
work head on the final cell erases that cell and moves left, entering the
rewind state when the erased bit $b$ is $1$ and staying in the reverse state
when $b$ is $0$.
\end{lemma}

\begin{lemma}[Reverse erasure computes the marker split]
\lean{Turing.FinTM.rawStrip_trim}
\label{Turing.FinTM.rawStrip_trim}
\leanok
\uses{Turing.Action.apply, Turing.Cfg.workTapeSymbols, Turing.FinTM.rawStripTM, Turing.FinTM.rawStrip_erase, Turing.FinTM.rawStrip_finish, Turing.FinTM.stripCfg, Turing.MultiTapeTM.runFrom, Turing.MultiTapeTM.runFrom_succ_eq_step, Turing.MultiTapeTM.step, Turing.splitAtLastTrue}
\difficulty{4}
From the stripping machine's reverse state with buffer $w$ and work head on
the final cell, there is a time $t \le 3(\abs{w}+1)$ after which the machine
has halted; its output is the prefix of $w$ strictly before the last $1$ bit
of $w$, and is empty when $w$ contains no $1$.
\end{lemma}

\begin{lemma}[Time bound for marker stripping]
\lean{Turing.FinTM.rawStrip_computes}
\label{Turing.FinTM.rawStrip_computes}
\leanok
\uses{Turing.Action.apply, Turing.Cfg.inputSymbol, Turing.FinTM.ComputesFunInTime, Turing.FinTM.ComputesInTime, Turing.FinTM.ComputesInTime.mono, Turing.FinTM.computesInTime_iff, Turing.FinTM.rawStripTM, Turing.FinTM.rawStrip_copy, Turing.FinTM.rawStrip_trim, Turing.FinTM.stripCfg, Turing.MultiTapeTM.initCfg, Turing.MultiTapeTM.runFrom, Turing.MultiTapeTM.runFrom_add, Turing.MultiTapeTM.runFrom_succ_eq_step', Turing.MultiTapeTM.step, Turing.splitAtLastTrue}
\difficulty{4}
For every $n$ and every input $x$ of length $n$, the marker-stripping
machine halts within $4(n+1)$ steps, and its output is the prefix of $x$
strictly before its last $1$ bit, or the empty word when $x$ contains no
$1$.
\end{lemma}

\begin{definition}[Marker-existence scanner]
\lean{Turing.FinTM.anyTrueTM}
\label{Turing.FinTM.anyTrueTM}
\leanok
\uses{Turing.FinTM}
A one-state machine without work tapes that scans its input left to right,
advancing silently over zeros, and halts with the one-bit verdict recording
whether the input contains a $1$ bit.
\end{definition}

\begin{lemma}[Run of the marker-existence scanner]
\lean{Turing.FinTM.anyTrue_run}
\label{Turing.FinTM.anyTrue_run}
\leanok
\uses{Turing.Action.apply, Turing.FinTM.anyTrueTM, Turing.FinTM.scanCfg, Turing.FinTM.scanCfg_read, Turing.FinTM.scanStep_right, Turing.MultiTapeTM.runFrom, Turing.MultiTapeTM.runFrom_succ_eq_step, Turing.MultiTapeTM.runFrom_zero, Turing.MultiTapeTM.step}
\difficulty{4}
For any split $x = p s$ of the input, there is a time $t \le \abs{s} + 1$
such that after $t$ steps from the scan configuration at position $\abs{p}$
with empty output, the marker-existence scanner has halted with output the
single bit recording whether $s$ contains a $1$.
\end{lemma}

\begin{lemma}[Time bound for the marker-existence test]
\lean{Turing.FinTM.anyTrue_computes}
\label{Turing.FinTM.anyTrue_computes}
\leanok
\uses{Turing.Cfg.ext_zero_tapes, Turing.FinTM.ComputesFunInTime, Turing.FinTM.ComputesInTime, Turing.FinTM.ComputesInTime.mono, Turing.FinTM.anyTrueTM, Turing.FinTM.anyTrue_run, Turing.FinTM.computesInTime_iff, Turing.FinTM.scanCfg, Turing.MultiTapeTM.initCfg}
\difficulty{3}
For every $n$ and every input of length $n$, the marker-existence scanner
halts within $n+1$ steps, and its output is the single bit recording whether
the input contains a $1$ bit.
\end{lemma}

\begin{lemma}[Decoding inverts the pair encoding]
\lean{Turing.FinTM.catalogPair_inverse}
\label{Turing.FinTM.catalogPair_inverse}
\leanok
\uses{Turing.pairDecode, Turing.pairEncode}
\difficulty{4}
If a word $x$ decodes under the self-delimiting pair code as the pair
$(a,v)$, then $x$ is exactly the pair encoding of $(a,v)$: each bit of $a$
doubled, then the separator bits $0,1$, then $v$.
\end{lemma}

\begin{lemma}[Dichotomy for the last-marker split]
\lean{Turing.FinTM.catalogMarker_cases}
\label{Turing.FinTM.catalogMarker_cases}
\leanok
\uses{Turing.splitAtLastTrue}
\difficulty{4}
Every word $v$ satisfies exactly one of: $v$ contains no $1$ and the
last-marker split of $v$ is undefined; or $v$ contains a $1$ and the
last-marker split of $v$ is some word $u$ (the prefix of $v$ strictly before
its last $1$) such that moreover, for every prefix $p$, the last-marker split
of $p v$ is $p u$.
\end{lemma}

\begin{lemma}[Extracted payload never lengthens the input]
\lean{Turing.FinTM.catalogPayload_length}
\label{Turing.FinTM.catalogPayload_length}
\leanok
\uses{Turing.FinTM.catalogPair_inverse, Turing.pairDecode, Turing.pairEncode}
\difficulty{3}
For every word $x$, the second component extracted from $x$ under the
self-delimiting pair code (taken empty when $x$ is malformed) has length at
most $\abs{x}$.
\end{lemma}

\begin{lemma}[Running a machine on the extracted payload]
\lean{Turing.FinTM.catalogPayload_computes}
\label{Turing.FinTM.catalogPayload_computes}
\leanok
\uses{Turing.Cfg.init, Turing.FinTM, Turing.FinTM.ComputesFunInTime, Turing.FinTM.ComputesInTime, Turing.FinTM.ComputesInTime.mono, Turing.FinTM.VirtualTag, Turing.FinTM.bufferedCompTM, Turing.FinTM.bufferedComp_start, Turing.FinTM.bufferedSecondCfg, Turing.FinTM.bufferedSecondCfg_run, Turing.FinTM.catalogPayload_length, Turing.FinTM.computesInTime_iff, Turing.FinTM.pairExtractTM, Turing.FinTM.pairExtract_computes, Turing.MultiTapeTM.initCfg, Turing.MultiTapeTM.runFrom_add, Turing.pairDecode}
\difficulty{5}
Let a machine compute the total function $g$ within a monotone time bound
$T_g$. Then the buffered composition of the second-component pair extractor
with that machine computes, within $6(n+1) + T_g(n) + 1$ steps on inputs of
length $n$, the function sending $x$ to $g(b)$, where $b$ is the second
component of $x$ under the self-delimiting pair code (the empty word when $x$
is malformed).
\end{lemma}

\begin{theorem}[Prepending a fixed word, machine contract]
\lean{Turing.FinTM.computesFunInTime_prepend}
\label{Turing.FinTM.computesFunInTime_prepend}
\leanok
\uses{Turing.FinTM, Turing.FinTM.ComputesFunInTime, Turing.FinTM.ComputesInTime.mono, Turing.FinTM.catalogPrefixTM, Turing.FinTM.catalogPrefixTM_computes}
\difficulty{2}
For every fixed word $w \in \{0,1\}^*$ there exist a finite Turing machine
$M$ and a constant $c$ (which may depend on $w$) such that $M$ computes the
function $x \mapsto w x$ (prepend $w$ to the input) within $c\,(n+1)$ steps
on inputs of length $n$.
\end{theorem}

\begin{theorem}[Input length in binary, machine contract]
\lean{Turing.FinTM.computesFunInTime_lengthBits}
\label{Turing.FinTM.computesFunInTime_lengthBits}
\leanok
\uses{Complexity.timeConstructible_id, Turing.FinTM, Turing.FinTM.ComputesFunInTime}
\difficulty{1}
There exist a finite Turing machine $M$ and a constant $c$ such that $M$
computes the function sending each input $x$ to the little-endian binary
representation of $\abs{x}$, within $c\,(n+1)$ steps on inputs of length $n$.
\end{theorem}

\begin{theorem}[Unary polynomial of the input length, machine contract]
\lean{Turing.FinTM.computesFunInTime_polyUnary}
\label{Turing.FinTM.computesFunInTime_polyUnary}
\leanok
\uses{Turing.FinTM, Turing.FinTM.ComputesFunInTime, Turing.FinTM.ComputesInTime.mono, Turing.FinTM.catalogPolyUnaryTM, Turing.FinTM.catalogPoly_unary_computes, Turing.FinTM.computesFunInTime_const}
\difficulty{3}
For all constants $C$ and $e$ there exist a finite Turing machine $M$ and a
constant $c$ such that $M$ computes the function sending each input of length
$n$ to the run of $C\,(n+1)^{e}$ ones, within $c\,(n+1)^{e+1}$ steps.
\end{theorem}

\begin{theorem}[Binary polynomial of the input length, machine contract]
\lean{Turing.FinTM.computesFunInTime_polyBits}
\label{Turing.FinTM.computesFunInTime_polyBits}
\leanok
\uses{Turing.FinTM, Turing.FinTM.ComputesFunInTime, Turing.FinTM.ComputesInTime.mono, Turing.FinTM.computesFunInTime_comp, Turing.FinTM.computesFunInTime_lengthBits, Turing.FinTM.computesFunInTime_polyUnary}
\difficulty{4}
For all constants $C$ and $e$ there exist a finite Turing machine $M$ and a
constant $c$ such that $M$ computes the function sending each input of length
$n$ to the little-endian binary representation of $C\,(n+1)^{e}$, within
$c\,(n+1)^{e+1}$ steps.
\end{theorem}

\begin{theorem}[Pairing with a fixed first component, machine contract]
\lean{Turing.FinTM.computesFunInTime_pairEncodeFixed}
\label{Turing.FinTM.computesFunInTime_pairEncodeFixed}
\leanok
\uses{Turing.FinTM, Turing.FinTM.ComputesFunInTime, Turing.FinTM.computesFunInTime_prepend, Turing.pairEncode}
\difficulty{1}
For every fixed word $\alpha$ there exist a finite Turing machine $M$ and a
constant $c$ such that $M$ computes the function sending $x$ to the
self-delimiting pair encoding of $(\alpha, x)$ --- each bit of $\alpha$
doubled, the separator bits $0,1$, then $x$ --- within $c\,(n+1)$ steps on
inputs of length $n$.
\end{theorem}

\begin{theorem}[First-component extraction, machine contract]
\lean{Turing.FinTM.computesFunInTime_pairFst}
\label{Turing.FinTM.computesFunInTime_pairFst}
\leanok
\uses{Turing.FinTM, Turing.FinTM.ComputesFunInTime, Turing.FinTM.pairExtractTM, Turing.FinTM.pairExtract_computes, Turing.pairDecode}
\difficulty{3}
There exist a finite Turing machine $M$ and a constant $c$ such that $M$
computes, within $c\,(n+1)$ steps on inputs of length $n$, the function
sending a word in the self-delimiting pair code to the first component of the
encoded pair, and any malformed word to the empty word.
\end{theorem}

\begin{theorem}[Second-component extraction, machine contract]
\lean{Turing.FinTM.computesFunInTime_pairSnd}
\label{Turing.FinTM.computesFunInTime_pairSnd}
\leanok
\uses{Turing.FinTM, Turing.FinTM.ComputesFunInTime, Turing.FinTM.pairExtractTM, Turing.FinTM.pairExtract_computes, Turing.pairDecode}
\difficulty{3}
There exist a finite Turing machine $M$ and a constant $c$ such that $M$
computes, within $c\,(n+1)$ steps on inputs of length $n$, the function
sending a word in the self-delimiting pair code to the second component of
the encoded pair, and any malformed word to the empty word.
\end{theorem}

\begin{theorem}[Pair-grammar validity, machine contract]
\lean{Turing.FinTM.computesFunInTime_pairValid}
\label{Turing.FinTM.computesFunInTime_pairValid}
\leanok
\uses{Turing.FinTM, Turing.FinTM.ComputesFunInTime, Turing.FinTM.pairValidTM, Turing.FinTM.pairValid_computes, Turing.pairDecode}
\difficulty{1}
There exist a finite Turing machine $M$ and a constant $c$ such that $M$
computes, within $c\,(n+1)$ steps on inputs of length $n$, the one-bit
function recording whether the input is a valid word of the self-delimiting
pair code.
\end{theorem}

\begin{theorem}[Pair to concatenation, machine contract]
\lean{Turing.FinTM.computesFunInTime_pairConcat}
\label{Turing.FinTM.computesFunInTime_pairConcat}
\leanok
\uses{Turing.FinTM, Turing.FinTM.ComputesFunInTime, Turing.FinTM.pairExtractTM, Turing.FinTM.pairExtract_computes, Turing.pairDecode}
\difficulty{3}
There exist a finite Turing machine $M$ and a constant $c$ such that $M$
computes, within $c\,(n+1)$ steps on inputs of length $n$, the function
sending a word that encodes the pair $(a,b)$ under the self-delimiting pair
code to the concatenation $a b$, and any malformed word to the empty word.
\end{theorem}

\begin{theorem}[Duplication into a pair, machine contract]
\lean{Turing.FinTM.computesFunInTime_pairDup}
\label{Turing.FinTM.computesFunInTime_pairDup}
\leanok
\uses{Turing.FinTM, Turing.FinTM.ComputesFunInTime, Turing.FinTM.pairDupTM, Turing.FinTM.pairDup_computes, Turing.pairEncode}
\difficulty{1}
There exist a finite Turing machine $M$ and a constant $c$ such that $M$
computes, within $c\,(n+1)$ steps on inputs of length $n$, the function
sending $x$ to the self-delimiting pair encoding of $(x,x)$: each bit of $x$
doubled, the separator bits $0,1$, then $x$.
\end{theorem}

\begin{definition}[Administrative action of the threaded map]
\lean{Turing.FinTM.mapAction}
\label{Turing.FinTM.mapAction}
\leanok
\uses{Turing.Action, Turing.FinTM}
For a machine $M$, an action of the threaded-map controller built over $M$:
it moves the input head and the final (capture) work head by the given
directions, optionally emits one output bit, and passes to the given control
state, writing no tape cell.
\end{definition}

\begin{definition}[Threaded-map controller machine]
\lean{Turing.FinTM.pairMapTM}
\label{Turing.FinTM.pairMapTM}
\leanok
\uses{Turing.FinTM, Turing.FinTM.controlAction, Turing.FinTM.mapAction, Turing.captureAction}
Given a machine $M$, the threaded-map controller over $M$ is a machine with
one extra work tape that first simulates $M$ on the input while capturing
the output of $M$ on the extra tape, rewinds the capture head and the input,
silently validates the input's membership in the self-delimiting pair code,
and only on success rewinds the input once more and emits the input's
doubled prefix and separator followed by the captured word; malformed inputs
halt without any emission.
\end{definition}

\begin{definition}[Threaded-map configuration]
\lean{Turing.FinTM.mapCfg}
\label{Turing.FinTM.mapCfg}
\leanok
\uses{Turing.Cfg, Turing.FinTM, Turing.FinTM.pairMapTM, Turing.captureCfg}
A configuration of the threaded-map controller that preserves the completed
tape bank of the simulated machine and its captured word, and records the
controller's state, the input position, the capture-head displacement, and
the output written so far.
\end{definition}

\begin{lemma}[Effect of a threaded-map administrative action]
\lean{Turing.FinTM.mapAction_apply}
\label{Turing.FinTM.mapAction_apply}
\leanok
\uses{Turing.Action.apply, Turing.Cfg, Turing.FinTM, Turing.FinTM.mapAction, Turing.FinTM.mapCfg, Turing.FinTM.pairMapTM, Turing.moveInputPos}
\difficulty{3}
Applying an administrative action of the threaded-map controller moves the
input position and the capture-head displacement by the specified directions,
installs the designated control state, and appends the optional emission to
the output; no work-tape cell changes.
\end{lemma}

\begin{lemma}[Capture-head rewind of the threaded map]
\lean{Turing.FinTM.mapBuffer_rewind}
\label{Turing.FinTM.mapBuffer_rewind}
\leanok
\uses{Turing.Action.apply, Turing.Cfg, Turing.Cfg.workTapeSymbols, Turing.FinTM, Turing.FinTM.bufferTape_left, Turing.FinTM.mapAction, Turing.FinTM.mapAction_apply, Turing.FinTM.mapCfg, Turing.FinTM.pairMapTM, Turing.MultiTapeTM.runFrom, Turing.MultiTapeTM.runFrom_succ_eq_step, Turing.MultiTapeTM.runFrom_zero, Turing.MultiTapeTM.step, Turing.captureCfg, Turing.moveInputPos_zero}
\difficulty{4}
From the controller's capture-rewind state with the capture head at
displacement $j-1$, for any $j$ at most the captured word's length, exactly
$j+1$ silent steps return the capture head to displacement $0$ and enter the
input-rewind phase, with the input position unchanged.
\end{lemma}

\begin{lemma}[Capture-and-rewind startup of the threaded map]
\lean{Turing.FinTM.mapStart}
\label{Turing.FinTM.mapStart}
\leanok
\uses{Turing.Action.apply, Turing.Cfg, Turing.Cfg.init, Turing.FinTM, Turing.FinTM.ComputesInTime, Turing.FinTM.ComputesInTime.output_unique, Turing.FinTM.catalogRewind, Turing.FinTM.computesInTime_iff, Turing.FinTM.mapAction, Turing.FinTM.mapBuffer_rewind, Turing.FinTM.mapCfg, Turing.FinTM.pairMapTM, Turing.MultiTapeTM.initCfg, Turing.MultiTapeTM.runFrom, Turing.MultiTapeTM.runFrom_add, Turing.MultiTapeTM.runFrom_succ_eq_step', Turing.MultiTapeTM.step, Turing.captureCfg, Turing.capture_run, Turing.moveInputPos_zero}
\difficulty{5}
Let $M$ be a machine that, on input $x$, halts within $T$ steps with output
$w$. Then within $T + \abs{w} + \abs{x} + 5$ steps the threaded-map
controller built over $M$ reaches, without emitting any physical output, its
silent validation state at input position $0$ with the capture head at $0$,
retaining a halted bank of $M$ whose captured output is $w$.
\end{lemma}

\begin{lemma}[Input symbol of a threaded-map configuration]
\lean{Turing.FinTM.mapCfg_read}
\label{Turing.FinTM.mapCfg_read}
\leanok
\uses{Turing.Cfg, Turing.Cfg.inputSymbol, Turing.FinTM, Turing.FinTM.inputSymbol_at, Turing.FinTM.mapCfg, Turing.FinTM.pairMapTM}
\difficulty{1}
In a threaded-map configuration whose input head has passed $i \le \abs{x}$
cells, the symbol read is the $i$-th bit of $x$ (the blank when
$i = \abs{x}$), independently of the preserved bank.
\end{lemma}

\begin{lemma}[First half of a threaded-map block]
\lean{Turing.FinTM.mapParse_first}
\label{Turing.FinTM.mapParse_first}
\leanok
\uses{Turing.Action.apply, Turing.Cfg, Turing.FinTM, Turing.FinTM.mapAction_apply, Turing.FinTM.mapCfg, Turing.FinTM.mapCfg_read, Turing.FinTM.pairMapTM, Turing.MultiTapeTM.step, Turing.moveInputPos_pos_of_ne_right}
\difficulty{3}
If the input splits as $x = p\, b\, s$ with bit $b$, one step of the
threaded-map controller from its block-start parse state at position
$\abs{p}$ records $b$ in the finite control and advances the input head,
emitting $b$ exactly when the controller is in its emitting (replay) pass.
\end{lemma}

\begin{lemma}[One aligned block of the threaded map]
\lean{Turing.FinTM.mapParse_block}
\label{Turing.FinTM.mapParse_block}
\leanok
\uses{Turing.Action.apply, Turing.Cfg, Turing.FinTM, Turing.FinTM.mapAction_apply, Turing.FinTM.mapCfg, Turing.FinTM.mapCfg_read, Turing.FinTM.mapParse_first, Turing.FinTM.pairMapTM, Turing.MultiTapeTM.runFrom, Turing.MultiTapeTM.step, Turing.moveInputPos_pos_of_ne_right}
\difficulty{4}
If the input splits as $x = p\, b\, d\, s$ with bits $b, d$, two steps of the
threaded-map controller from its block-start parse state at position
$\abs{p}$ continue parsing when $b = d$, halt when $(b,d) = (1,0)$, and pass
to the post-separator state when $(b,d) = (0,1)$; both bits are emitted
exactly when the controller is in its emitting pass. The captured word and
its head are preserved.
\end{lemma}

\begin{lemma}[Silent validation pass of the threaded map]
\lean{Turing.FinTM.mapValidate}
\label{Turing.FinTM.mapValidate}
\leanok
\uses{Turing.Action.apply, Turing.Cfg, Turing.FinTM, Turing.FinTM.mapAction, Turing.FinTM.mapCfg, Turing.FinTM.mapCfg_read, Turing.FinTM.mapParse_block, Turing.FinTM.mapParse_first, Turing.FinTM.pairMapTM, Turing.MultiTapeTM.runFrom, Turing.MultiTapeTM.runFrom_add, Turing.MultiTapeTM.runFrom_succ_eq_step, Turing.MultiTapeTM.runFrom_zero, Turing.MultiTapeTM.step, Turing.captureCfg, Turing.pairDecode}
\difficulty{5}
For a split $x = p s$ of the input, there is a time $t \le \abs{s} + 1$ from
the controller's silent validation state at position $\abs{p}$ such that: if
$s$ is a valid word of the self-delimiting pair code, after $t$ steps the
controller is at its post-validation state with nothing emitted; and if $s$
is malformed, after $t$ steps the controller has halted with empty output.
\end{lemma}

\begin{lemma}[Encoded-prefix replay of the threaded map]
\lean{Turing.FinTM.mapPrefix_replay}
\label{Turing.FinTM.mapPrefix_replay}
\leanok
\uses{Turing.Cfg, Turing.FinTM, Turing.FinTM.mapCfg, Turing.FinTM.mapParse_block, Turing.FinTM.pairMapTM, Turing.MultiTapeTM.runFrom, Turing.MultiTapeTM.runFrom_add, Turing.pairEncode}
\difficulty{4}
If the input splits as $x = p$ followed by the pair encoding of $(a,b)$, then
from the controller's emitting parse state at position $\abs{p}$ with prior
output $o$, exactly $2\abs{a} + 2$ steps emit the doubled word $a$ and the
separator bits $0,1$, and enter the captured-payload replay state with the
capture head still at displacement $0$.
\end{lemma}

\begin{lemma}[Captured-payload replay of the threaded map]
\lean{Turing.FinTM.mapPayload_replay}
\label{Turing.FinTM.mapPayload_replay}
\leanok
\uses{Turing.Action.apply, Turing.Cfg, Turing.Cfg.workTapeSymbols, Turing.FinTM, Turing.FinTM.mapAction_apply, Turing.FinTM.mapCfg, Turing.FinTM.pairMapTM, Turing.MultiTapeTM.runFrom, Turing.MultiTapeTM.runFrom_succ_eq_step', Turing.MultiTapeTM.step, Turing.captureCfg, Turing.moveInputPos_zero}
\difficulty{4}
For every $j$ at most the captured word's length and every prior output
$o$, running the controller for exactly $j$ steps from its payload-replay
state with the capture head at displacement $0$ yields the payload-replay
configuration with the capture head at displacement $j$ and output $o$
followed by the first $j$ captured bits; the bank is preserved.
\end{lemma}

\begin{lemma}[End of captured-payload replay]
\lean{Turing.FinTM.mapPayload_finish}
\label{Turing.FinTM.mapPayload_finish}
\leanok
\uses{Turing.Action.apply, Turing.Cfg, Turing.Cfg.workTapeSymbols, Turing.FinTM, Turing.FinTM.mapAction, Turing.FinTM.mapCfg, Turing.FinTM.mapPayload_replay, Turing.FinTM.pairMapTM, Turing.MultiTapeTM.runFrom, Turing.MultiTapeTM.runFrom_succ_eq_step', Turing.MultiTapeTM.step, Turing.captureCfg}
\difficulty{2}
From the controller's payload-replay state with the capture head at
displacement $0$ and prior output $o$, after one more step than the captured
word's length the controller has halted with output $o$ followed by the
entire captured word.
\end{lemma}

\begin{lemma}[Length of a pair encoding]
\lean{Turing.FinTM.catalogPair_length}
\label{Turing.FinTM.catalogPair_length}
\leanok
\uses{Turing.pairEncode}
\difficulty{1}
For all words $a$ and $b$, the self-delimiting pair encoding of $(a,b)$ has
length exactly $2\abs{a} + 2 + \abs{b}$.
\end{lemma}

\begin{lemma}[Time bound for the threaded-map controller]
\lean{Turing.FinTM.pairMap_computes}
\label{Turing.FinTM.pairMap_computes}
\leanok
\uses{Turing.FinTM, Turing.FinTM.ComputesFunInTime, Turing.FinTM.ComputesInTime, Turing.FinTM.ComputesInTime.mono, Turing.FinTM.catalogPair_inverse, Turing.FinTM.catalogPair_length, Turing.FinTM.catalogRewind, Turing.FinTM.computesInTime_iff, Turing.FinTM.mapCfg, Turing.FinTM.mapPayload_finish, Turing.FinTM.mapPrefix_replay, Turing.FinTM.mapStart, Turing.FinTM.mapValidate, Turing.FinTM.pairMapTM, Turing.MultiTapeTM.initCfg, Turing.MultiTapeTM.output_length_le, Turing.MultiTapeTM.runFrom, Turing.MultiTapeTM.runFrom_add, Turing.pairDecode, Turing.pairEncode}
\difficulty{5}
Let $M$ compute the total function $f$ within time $T$. Then the threaded-map
controller built over $M$ computes, within $4(T(n) + n + 3)$ steps on inputs
of length $n$, the function sending a word $x$ that decodes under the
self-delimiting pair code with first component $a$ to the pair encoding of
$(a, f(x))$, and any malformed word to the empty word.
\end{lemma}

\begin{theorem}[Threaded map on the second component, machine contract]
\lean{Turing.FinTM.computesFunInTime_pairMapSnd}
\label{Turing.FinTM.computesFunInTime_pairMapSnd}
\leanok
\uses{Turing.FinTM, Turing.FinTM.ComputesFunInTime, Turing.FinTM.ComputesInTime.mono, Turing.FinTM.bufferedCompTM, Turing.FinTM.catalogPayload_computes, Turing.FinTM.pairExtractTM, Turing.FinTM.pairMapTM, Turing.FinTM.pairMap_computes, Turing.pairDecode, Turing.pairEncode}
\difficulty{3}
Let a machine compute the total function $g$ within a monotone time bound
$T_g$. Then there exist a finite Turing machine $M$ and a constant $c$ such
that $M$ computes, within $c\,(n + 1 + T_g(n))$ steps on inputs of length
$n$, the function sending a word that encodes the pair $(a,b)$ under the
self-delimiting pair code to the pair encoding of $(a, g(b))$, and any
malformed word to the empty word.
\end{theorem}

\begin{theorem}[Threaded length-bound check, machine contract]
\lean{Turing.FinTM.computesFunInTime_pairLenCheck}
\label{Turing.FinTM.computesFunInTime_pairLenCheck}
\leanok
\uses{Turing.FinTM, Turing.FinTM.ComputesFunInTime, Turing.FinTM.ComputesInTime, Turing.FinTM.ComputesInTime.mono, Turing.FinTM.computesFunInTime_comp, Turing.FinTM.computesFunInTime_pairFst, Turing.FinTM.computesFunInTime_polyUnary, Turing.FinTM.pairCountTM, Turing.FinTM.pairCount_computes, Turing.pairDecode}
\difficulty{4}
For all constants $C$ and $e$ there exist a finite Turing machine $M$ and a
constant $c$ such that $M$ computes, within $c\,(n+1)^{e+1}$ steps on inputs
of length $n$, the one-bit function that on a word encoding the pair $(a,b)$
under the self-delimiting pair code outputs whether
$\abs{b} \le C\,(\abs{a}+1)^{e}$, and on any malformed word outputs $0$.
\end{theorem}

\begin{theorem}[Stripping at the last marker, machine contract]
\lean{Turing.FinTM.computesFunInTime_stripLast}
\label{Turing.FinTM.computesFunInTime_stripLast}
\leanok
\uses{Turing.FinTM, Turing.FinTM.ComputesFunInTime, Turing.FinTM.ComputesInTime, Turing.FinTM.ComputesInTime.mono, Turing.FinTM.anyTrue_computes, Turing.FinTM.catalogMarker_cases, Turing.FinTM.catalogPair_inverse, Turing.FinTM.computesFunInTime_comp, Turing.FinTM.computesFunInTime_cond, Turing.FinTM.computesFunInTime_const, Turing.FinTM.computesFunInTime_pairSnd, Turing.FinTM.rawStrip_computes, Turing.pairDecode, Turing.pairEncode, Turing.splitAtLastTrue}
\difficulty{5}
There exist a finite Turing machine $M$ and a constant $c$ such that $M$
computes, within $c\,(n+1)^{2}$ steps on inputs of length $n$, the following
function: a word encoding the pair $(a,v)$ under the self-delimiting pair
code, where $v$ contains a $1$ and $u$ is the prefix of $v$ strictly before
its last $1$, is sent to the pair encoding of $(a,u)$; a word whose second
component contains no $1$, or any malformed word, is sent to the empty word.
\end{theorem}

\begin{definition}[One step of the split search]
\lean{Turing.FinTM.splitStep}
\label{Turing.FinTM.splitStep}
\leanok
For a word $w$, the split-search step sends a candidate word $s$ with
$\abs{s} \le \abs{w}$ to $s$ followed by one $1$ bit, and sends any longer
candidate to itself.
\end{definition}

\begin{definition}[Split-search acceptance test]
\lean{Turing.FinTM.splitAccept}
\label{Turing.FinTM.splitAccept}
\leanok
The acceptance test of the split search with parameters $C$ and $e$: a
candidate $s$ is accepted for the word $w$ exactly when
$\abs{s} + C\,(\abs{s}+1)^{e} = \abs{w}$.
\end{definition}

\begin{lemma}[Length invariant of the split step]
\lean{Turing.FinTM.splitStep_inv}
\label{Turing.FinTM.splitStep_inv}
\leanok
\uses{Turing.FinTM.splitStep}
\difficulty{2}
If a candidate $s$ has length at most $\abs{w} + 1$, then so does the result
of one split-search step on $s$.
\end{lemma}

\begin{lemma}[Orbit of the split step]
\lean{Turing.FinTM.splitStep_orbit}
\label{Turing.FinTM.splitStep_orbit}
\leanok
\uses{Turing.FinTM.splitStep}
\difficulty{3}
For every $i \le \abs{w} + 1$, the $i$-th iterate of the split-search step on
the empty candidate is the unary word $1^i$.
\end{lemma}

\begin{lemma}[Search is extensional in its predicate]
\lean{Turing.FinTM.catalogFind_congr}
\label{Turing.FinTM.catalogFind_congr}
\leanok
\difficulty{3}
If two Boolean predicates agree on every member of a list, then searching the
list for the first element satisfying the one returns the same result as
searching for the first element satisfying the other.
\end{lemma}

\begin{lemma}[Orbit search agrees with the arithmetic split search]
\lean{Turing.FinTM.splitFind_eq}
\label{Turing.FinTM.splitFind_eq}
\leanok
\uses{Turing.FinTM.catalogFind_congr, Turing.FinTM.splitAccept, Turing.FinTM.splitStep, Turing.FinTM.splitStep_orbit, Turing.solveSplit}
\difficulty{3}
For all $C$ and $e$ and every word $w$, searching the indices
$0, \dots, \abs{w}$ for the first $i$ at which the split-acceptance test
succeeds on the $i$-th iterate of the split step returns exactly the
arithmetic split search at parameters $C$, $e$, and $\abs{w}$, i.e.\ the
least $i$ with $i + C\,(i+1)^{e} = \abs{w}$ when one exists and nothing
otherwise.
\end{lemma}

\begin{lemma}[Failure of the split search]
\lean{Turing.FinTM.splitFind_none}
\label{Turing.FinTM.splitFind_none}
\leanok
\uses{Turing.FinTM.splitAccept, Turing.FinTM.splitFind_eq, Turing.FinTM.splitStep, Turing.solveSplit}
\difficulty{1}
The arithmetic split search at parameters $C$, $e$, and $\abs{w}$ fails if
and only if the split-acceptance test rejects the $i$-th iterate of the split
step on the empty candidate for every $i \le \abs{w}$.
\end{lemma}

\begin{lemma}[Result of the split-search loop]
\lean{Turing.FinTM.splitLoop_result}
\label{Turing.FinTM.splitLoop_result}
\leanok
\uses{Turing.FinTM.splitAccept, Turing.FinTM.splitFind_eq, Turing.FinTM.splitStep, Turing.FinTM.splitStep_orbit, Turing.pairEncode, Turing.solveSplit}
\difficulty{3}
Splitting $w$ at the length of the first accepted orbit candidate and pair
encoding the two halves gives the same word as splitting $w$ at the index
returned by the arithmetic split search; when both searches fail, both sides
are the empty word.
\end{lemma}

\begin{lemma}[Polynomial bound for the search loop overhead]
\lean{Turing.FinTM.splitLoop_bound}
\label{Turing.FinTM.splitLoop_bound}
\leanok
\difficulty{3}
For all naturals $c$, $A$, $e$, $n$,
\[
  c\,\bigl(A\,(n+1)^{e+1} + 1\bigr)(n+2) \;\le\; 2c\,(A+1)\,(n+1)^{e+2}.
\]
\end{lemma}

\begin{definition}[Saturated input position]
\lean{Turing.FinTM.splitPos}
\label{Turing.FinTM.splitPos}
\leanok
For an input $w$ and a count $j$, the saturated input position is
$\min(j, \abs{w}) + 1$: the physical input-head position after consuming
$j$ cells, saturated at the right boundary.
\end{definition}

\begin{lemma}[Reading at a saturated position]
\lean{Turing.FinTM.splitPos_read}
\label{Turing.FinTM.splitPos_read}
\leanok
\uses{Turing.Cfg, Turing.Cfg.inputSymbol, Turing.FinTM.splitPos, Turing.inputSymbolInner}
\difficulty{3}
In any configuration whose input head is at the saturated position for $j$
on input $w$, the symbol read is the $j$-th bit of $w$ when $j < \abs{w}$ and
the blank otherwise.
\end{lemma}

\begin{lemma}[Advancing a saturated position]
\lean{Turing.FinTM.splitPos_succ}
\label{Turing.FinTM.splitPos_succ}
\leanok
\uses{Turing.FinTM.splitPos, Turing.moveInputPos, Turing.moveInputPos_pos_of_ne_right, Turing.moveInputPos_rightBoundary}
\difficulty{3}
For every input $w$ and every count $j$, moving the input head one cell to
the right from the saturated position for $j$ on $w$ yields the saturated
position for $j+1$.
\end{lemma}

\begin{definition}[Partially cleared unary scratch tape]
\lean{Turing.FinTM.splitScratch}
\label{Turing.FinTM.splitScratch}
\leanok
The tape content of a unary scratch word of length $q$ whose first $j$ cells
have been cleared: cell $z$ holds a $1$ for $j \le z < q$ and is blank
otherwise.
\end{definition}

\begin{lemma}[Clearing one scratch cell]
\lean{Turing.FinTM.splitScratch_erase}
\label{Turing.FinTM.splitScratch_erase}
\leanok
\uses{Turing.FinTM.splitScratch}
\difficulty{3}
For all $q$ and $j$, blanking cell $j$ of the unary scratch tape of length
$q$ with $j$ cells already cleared yields the scratch tape with $j+1$ cells
cleared.
\end{lemma}

\begin{definition}[Rejection-cleanup machine of the split search]
\lean{Turing.FinTM.splitRestoreTM}
\label{Turing.FinTM.splitRestoreTM}
\leanok
\uses{Turing.FinTM, Turing.FinTM.controlAction}
For $k$ scratch tapes, a machine with $k+1$ work tapes implementing the
rejection cleanup of the split search: it scans the candidate word on the
first work tape while clearing every unary scratch tape, appends one $1$ bit
to the candidate exactly when the candidate's length does not exceed the
input length, rewinds all work heads along the candidate and then rewinds the
input head, and finally stalls in an absorbing return state.
\end{definition}

\begin{definition}[Cleanup scan configuration]
\lean{Turing.FinTM.splitRestoreScan}
\label{Turing.FinTM.splitRestoreScan}
\leanok
\uses{Turing.Cfg, Turing.FinTM.bufferTape, Turing.FinTM.splitPos, Turing.FinTM.splitRestoreTM, Turing.FinTM.splitScratch}
A configuration of the cleanup machine in which $j$ candidate cells have been
consumed, exactly that prefix has been erased on each scratch tape, and the
physical input head tracks the same count at its saturated position.
\end{definition}

\begin{lemma}[Silent cleanup scan]
\lean{Turing.FinTM.splitRestore_scan}
\label{Turing.FinTM.splitRestore_scan}
\leanok
\uses{Turing.Action.apply, Turing.Cfg.inputSymbol, Turing.Cfg.workTapeSymbols, Turing.FinTM.splitPos_read, Turing.FinTM.splitPos_succ, Turing.FinTM.splitRestoreScan, Turing.FinTM.splitRestoreTM, Turing.FinTM.splitScratch_erase, Turing.MultiTapeTM.runFrom, Turing.MultiTapeTM.runFrom_succ_eq_step', Turing.MultiTapeTM.step}
\difficulty{4}
For every $j$ at most the candidate's length, $j$ steps of the cleanup
machine from the start of its scan phase reach the cleanup scan configuration
for $j$, with nothing emitted; the overflow flag records exactly whether more
candidate cells than input cells have been consumed.
\end{lemma}

\begin{definition}[Cleaned configuration]
\lean{Turing.FinTM.splitRestoreClean}
\label{Turing.FinTM.splitRestoreClean}
\leanok
\uses{Turing.Cfg, Turing.FinTM.bufferTape, Turing.FinTM.splitRestoreTM}
A configuration of the cleanup machine in which only the candidate word
remains, on the first work tape; all work heads sit at a common displacement
and the physical output is empty.
\end{definition}

\begin{lemma}[End-of-scan append step]
\lean{Turing.FinTM.splitRestore_append}
\label{Turing.FinTM.splitRestore_append}
\leanok
\uses{Turing.Action.apply, Turing.Cfg.workTapeSymbols, Turing.FinTM.bufferTape_append, Turing.FinTM.splitPos, Turing.FinTM.splitRestoreClean, Turing.FinTM.splitRestoreScan, Turing.FinTM.splitRestoreTM, Turing.FinTM.splitScratch, Turing.FinTM.splitScratch_erase, Turing.FinTM.splitStep, Turing.MultiTapeTM.step, Turing.moveInputPos_zero}
\difficulty{4}
One step of the cleanup machine from the end of its scan phase clears the
final extra scratch cell and appends a $1$ to the candidate exactly when the
candidate's length is at most the input length; the result is the cleaned
configuration holding one split-search step applied to the candidate.
\end{lemma}

\begin{lemma}[Candidate-guided head rewind]
\lean{Turing.FinTM.splitRestore_rewind}
\label{Turing.FinTM.splitRestore_rewind}
\leanok
\uses{Turing.Action.apply, Turing.Cfg.workTapeSymbols, Turing.FinTM.bufferTape_left, Turing.FinTM.bufferTape_nat, Turing.FinTM.splitRestoreClean, Turing.FinTM.splitRestoreScan, Turing.FinTM.splitRestoreTM, Turing.MultiTapeTM.runFrom, Turing.MultiTapeTM.runFrom_succ_eq_step, Turing.MultiTapeTM.runFrom_zero, Turing.MultiTapeTM.step, Turing.moveInputPos_zero}
\difficulty{4}
For every $j$ at most the candidate's length, exactly $j+1$ steps of the
cleanup machine from the cleaned configuration with all work heads at
displacement $j-1$ yield the cleaned configuration of the next phase with
every work head at displacement $0$; no candidate bit is altered.
\end{lemma}

\begin{lemma}[Complete rejection cleanup]
\lean{Turing.FinTM.splitRestore_run}
\label{Turing.FinTM.splitRestore_run}
\leanok
\uses{Turing.Cfg.ofWords, Turing.FinTM.catalogRewind, Turing.FinTM.splitPos, Turing.FinTM.splitRestoreClean, Turing.FinTM.splitRestoreScan, Turing.FinTM.splitRestoreTM, Turing.FinTM.splitRestore_append, Turing.FinTM.splitRestore_rewind, Turing.FinTM.splitRestore_scan, Turing.FinTM.splitStep, Turing.MultiTapeTM.runFrom, Turing.MultiTapeTM.runFrom_add, Turing.MultiTapeTM.runFrom_succ_eq_step', Turing.stateWord}
\difficulty{4}
For any input $w$ and candidate $s$, there is a time
$t \le 2\abs{s} + \abs{w} + 5$ after which the cleanup machine, started at
the beginning of its scan phase, reaches the canonical anchored
configuration, in its return state, whose stored words carry one
split-search step applied to $s$; every transition on the way is silent.
\end{lemma}

\begin{definition}[Emission-counting action transform]
\lean{Turing.FinTM.splitCountAction}
\label{Turing.FinTM.splitCountAction}
\leanok
\uses{Turing.Action}
The transform that converts an action of a source machine into an action of
the counting host: each output emission of the source becomes one rightward
move of the native input head, the first work tape (holding the candidate) is
untouched, the source's tape writes and moves are reproduced on the
successor-indexed tapes, nothing is emitted, and a finite flag records
consumption past the native right boundary.
\end{definition}

\begin{definition}[Counting-host configuration]
\lean{Turing.FinTM.splitCountCfg}
\label{Turing.FinTM.splitCountCfg}
\leanok
\uses{Turing.Cfg, Turing.FinTM.bufferTape, Turing.FinTM.splitPos}
The configuration of the counting host that embeds a source configuration
over the empty virtual input: the candidate $s$ sits on the first work tape,
the source's work tapes occupy the remaining tapes, and the native input head
is at the saturated position for $\abs{s}$ plus the length of the source's
output so far.
\end{definition}

\begin{lemma}[Overflow-flag update]
\lean{Turing.FinTM.splitCount_over}
\label{Turing.FinTM.splitCount_over}
\leanok
\uses{Turing.Cfg, Turing.Cfg.inputSymbol, Turing.FinTM.splitPos, Turing.FinTM.splitPos_read}
\difficulty{3}
For a configuration whose input head is at the saturated position for $j$ on
input $w$, the disjunction of the flag recording $\abs{w} < j$ with the test
that the symbol read is blank equals the flag recording $\abs{w} < j + 1$.
\end{lemma}

\begin{lemma}[One-step correspondence of the counting host]
\lean{Turing.FinTM.splitCount_apply}
\label{Turing.FinTM.splitCount_apply}
\leanok
\uses{Turing.Action, Turing.Action.apply, Turing.Cfg, Turing.Cfg.inputSymbol, Turing.FinTM.splitCountAction, Turing.FinTM.splitCountCfg, Turing.FinTM.splitCount_over, Turing.FinTM.splitPos_succ}
\difficulty{4}
Applying the transformed action to the counting-host embedding of a source
configuration yields exactly the counting-host embedding of the source
configuration after the original action: one native input cell is consumed
per source emission, the candidate is preserved, and all source-bank writes
and moves are reproduced.
\end{lemma}

\begin{lemma}[Counted simulation of a source run]
\lean{Turing.FinTM.splitCount_run}
\label{Turing.FinTM.splitCount_run}
\leanok
\uses{Turing.Cfg, Turing.Cfg.Halted, Turing.Cfg.inputSymbol, Turing.Cfg.workTapeSymbols, Turing.FinTM.splitCountAction, Turing.FinTM.splitCountCfg, Turing.FinTM.splitCount_apply, Turing.MultiTapeTM, Turing.MultiTapeTM.runFrom, Turing.MultiTapeTM.runFrom_succ_eq_step', Turing.MultiTapeTM.step}
\difficulty{5}
Suppose the host machine's transitions on embedded states agree with the
emission-counting transform of the source machine's transitions, and the
source, started from a configuration $c$ over the empty input, does not halt
before time $t$. Then $t$ steps of the host from the embedding of $c$ land
exactly on the embedding of the source configuration after $t$ steps; the
host's physical output stays empty throughout.
\end{lemma}

\begin{lemma}[End of the generator loop over an empty input]
\lean{Turing.FinTM.splitPoly_loop_end}
\label{Turing.FinTM.splitPoly_loop_end}
\leanok
\uses{Turing.Action.apply, Turing.FinTM.catalogPolyCfg, Turing.FinTM.catalogPolyCost, Turing.FinTM.catalogPolyUnaryTM, Turing.FinTM.catalogPoly_loop, Turing.MultiTapeTM.runFrom, Turing.MultiTapeTM.runFrom_succ_eq_step', Turing.MultiTapeTM.step}
\difficulty{4}
On the empty input, with every unary loop tape of side $q > 0$ already
installed and all heads at $0$, the unary polynomial generator runs for
exactly one more step than the exact nest cost at level $c+1$ and halts,
having emitted $C\,q^{c+1}$ ones, with every loop head back at $0$.
\end{lemma}

\begin{lemma}[First entry into an absorbing state]
\lean{Turing.FinTM.catalogFirstEntry}
\label{Turing.FinTM.catalogFirstEntry}
\leanok
\uses{Turing.Cfg, Turing.MultiTapeTM, Turing.MultiTapeTM.runFrom, Turing.MultiTapeTM.runFrom_add, Turing.MultiTapeTM.step}
\difficulty{4}
Let $q$ be a control state on which the machine stalls (every configuration
in state $q$ is a fixed point of the step function), and suppose the run from
$c$ reaches a configuration $d$ in state $q$ after $T$ steps. Then there is a
time $t \le T$ such that no configuration strictly before time $t$ is in
state $q$ and the configuration at time $t$ is already $d$.
\end{lemma}

\begin{lemma}[First-entry form of the rejection cleanup]
\lean{Turing.FinTM.splitRestore_first}
\label{Turing.FinTM.splitRestore_first}
\leanok
\uses{Turing.Action.apply, Turing.Cfg, Turing.Cfg.ofWords, Turing.FinTM.catalogFirstEntry, Turing.FinTM.controlAction, Turing.FinTM.controlAction_apply, Turing.FinTM.splitRestoreScan, Turing.FinTM.splitRestoreTM, Turing.FinTM.splitRestore_run, Turing.FinTM.splitStep, Turing.MultiTapeTM.runFrom, Turing.MultiTapeTM.runFrom_zero, Turing.MultiTapeTM.step, Turing.moveInputPos_zero, Turing.stateWord}
\difficulty{4}
For any input $w$ and candidate $s$, there is a time $t$ with
$0 < t \le 2\abs{s} + \abs{w} + 5$ such that the cleanup machine never visits
its return state before time $t$ and at time $t$ sits in the canonical
anchored return configuration carrying one split-search step applied to $s$.
\end{lemma}

\begin{lemma}[Acceptance read of the counting host]
\lean{Turing.FinTM.splitCount_accept}
\label{Turing.FinTM.splitCount_accept}
\leanok
\uses{Turing.Cfg, Turing.Cfg.inputSymbol, Turing.FinTM.splitCountCfg, Turing.FinTM.splitPos_read}
\difficulty{3}
In the counting-host embedding of a source configuration, the conjunction of
the negated overflow flag with the test that the native input read is blank
holds exactly when the candidate's length plus the source's output length
equals the input length.
\end{lemma}

\begin{lemma}[Stopping the counted simulation at the first halt]
\lean{Turing.FinTM.splitCount_firstHalt}
\label{Turing.FinTM.splitCount_firstHalt}
\leanok
\uses{Turing.Cfg, Turing.FinTM.splitCountAction, Turing.FinTM.splitCountCfg, Turing.FinTM.splitCount_run, Turing.MultiTapeTM, Turing.MultiTapeTM.runFrom, Turing.MultiTapeTM.runFrom_add, Turing.MultiTapeTM.runFrom_of_halt}
\difficulty{4}
Suppose the host's transitions on embedded states agree with the
emission-counting transform of the source's transitions, and the source run
from $c$ reaches a halted configuration $d$ after $T$ steps. Then there is a
time $t \le T$ at which the host run from the embedding of $c$ equals the
embedding of $d$.
\end{lemma}

\begin{definition}[Preparation machine of the split search]
\lean{Turing.FinTM.splitPrepareTM}
\label{Turing.FinTM.splitPrepareTM}
\leanok
\uses{Turing.FinTM, Turing.FinTM.controlAction}
For $k$ scratch tapes, a machine with $k+1$ work tapes that prepares the
polynomial loop bank of the split search: it scans the candidate on the first
work tape while writing its length in unary onto every scratch tape in
parallel, adds one extra side-length cell, rewinds all work heads along the
untouched candidate, and stalls in an absorbing return state.
\end{definition}

\begin{definition}[Preparation scan configuration]
\lean{Turing.FinTM.splitPrepareScan}
\label{Turing.FinTM.splitPrepareScan}
\leanok
\uses{Turing.Cfg, Turing.FinTM.bufferTape, Turing.FinTM.catalogPolyTape, Turing.FinTM.splitPos, Turing.FinTM.splitPrepareTM}
A configuration of the preparation machine in which $j$ candidate cells have
been scanned and every scratch tape holds the unary word of length $j$, with
the input head at its saturated position for $j$.
\end{definition}

\begin{lemma}[Silent preparation scan]
\lean{Turing.FinTM.splitPrepare_scan}
\label{Turing.FinTM.splitPrepare_scan}
\leanok
\uses{Turing.Action.apply, Turing.Cfg.workTapeSymbols, Turing.FinTM.catalogPolyTape_write, Turing.FinTM.splitCount_over, Turing.FinTM.splitPos_succ, Turing.FinTM.splitPrepareScan, Turing.FinTM.splitPrepareTM, Turing.MultiTapeTM.runFrom, Turing.MultiTapeTM.runFrom_succ_eq_step', Turing.MultiTapeTM.step}
\difficulty{4}
For every $j$ at most the candidate's length, $j$ steps of the preparation
machine from the start of its scan yield the preparation scan configuration
for $j$; the candidate is unchanged and nothing is emitted, and this holds
also for candidates longer than the native input.
\end{lemma}

\begin{definition}[Prepared configuration]
\lean{Turing.FinTM.splitPrepareReady}
\label{Turing.FinTM.splitPrepareReady}
\leanok
\uses{Turing.Cfg, Turing.FinTM.bufferTape, Turing.FinTM.catalogPolyTape, Turing.FinTM.splitPos, Turing.FinTM.splitPrepareTM}
A configuration of the preparation machine in which every scratch tape holds
the unary word of length one more than the candidate's length, all work
heads sit at a common displacement, and the finite flag records whether the
candidate is longer than the input.
\end{definition}

\begin{lemma}[Adding the extra side-length cell]
\lean{Turing.FinTM.splitPrepare_extra}
\label{Turing.FinTM.splitPrepare_extra}
\leanok
\uses{Turing.Action.apply, Turing.Cfg.workTapeSymbols, Turing.FinTM.catalogPolyTape_write, Turing.FinTM.splitPrepareReady, Turing.FinTM.splitPrepareScan, Turing.FinTM.splitPrepareTM, Turing.MultiTapeTM.step, Turing.moveInputPos_zero}
\difficulty{3}
For every input $w$ and candidate $s$, one step of the preparation machine
from the end of its scan yields the prepared configuration in its rewind
phase, in which every scratch tape holds the unary word of length
$\abs{s}+1$ and all work heads are at displacement $\abs{s}-1$; the uniform
extra cell makes the empty candidate need no special case.
\end{lemma}

\begin{lemma}[Rewind of the prepared bank]
\lean{Turing.FinTM.splitPrepare_rewind}
\label{Turing.FinTM.splitPrepare_rewind}
\leanok
\uses{Turing.Action.apply, Turing.Cfg.workTapeSymbols, Turing.FinTM.bufferTape_left, Turing.FinTM.bufferTape_nat, Turing.FinTM.splitPrepareReady, Turing.FinTM.splitPrepareTM, Turing.MultiTapeTM.runFrom, Turing.MultiTapeTM.runFrom_succ_eq_step, Turing.MultiTapeTM.runFrom_zero, Turing.MultiTapeTM.step, Turing.moveInputPos_zero}
\difficulty{4}
For every $j$ at most the candidate's length, exactly $j+1$ steps of the
preparation machine from the prepared configuration in its rewind phase with
all work heads at displacement $j-1$ yield the prepared configuration in its
return phase with every head at displacement $0$; each scratch tape keeps
its extra unary cell.
\end{lemma}

\begin{lemma}[Exact cost of preparation]
\lean{Turing.FinTM.splitPrepare_run}
\label{Turing.FinTM.splitPrepare_run}
\leanok
\uses{Turing.Cfg.ofWords, Turing.FinTM.catalogPolyTape, Turing.FinTM.splitPos, Turing.FinTM.splitPrepareReady, Turing.FinTM.splitPrepareScan, Turing.FinTM.splitPrepareTM, Turing.FinTM.splitPrepare_extra, Turing.FinTM.splitPrepare_rewind, Turing.FinTM.splitPrepare_scan, Turing.MultiTapeTM.runFrom, Turing.MultiTapeTM.runFrom_add, Turing.MultiTapeTM.runFrom_succ_eq_step', Turing.stateWord}
\difficulty{4}
For every input $w$ and candidate $s$, exactly $2(\abs{s}+1)$ silent steps
of the preparation machine from the canonical anchored configuration whose
stored words carry $s$ yield the prepared configuration, in which every
scratch tape holds the unary word of length $\abs{s}+1$ and every head is at
displacement $0$.
\end{lemma}

\begin{lemma}[First-entry form of preparation]
\lean{Turing.FinTM.splitPrepare_first}
\label{Turing.FinTM.splitPrepare_first}
\leanok
\uses{Turing.Action.apply, Turing.Cfg.ofWords, Turing.FinTM.catalogFirstEntry, Turing.FinTM.controlAction, Turing.FinTM.controlAction_apply, Turing.FinTM.splitPrepareReady, Turing.FinTM.splitPrepareTM, Turing.FinTM.splitPrepare_run, Turing.MultiTapeTM.runFrom, Turing.MultiTapeTM.step, Turing.moveInputPos_zero, Turing.stateWord}
\difficulty{3}
For any input $w$ and candidate $s$, there is a time $t \le 2(\abs{s}+1)$
such that the preparation machine, started from the canonical anchored
configuration carrying $s$, never visits its return state before time $t$
and at time $t$ sits in the prepared configuration with all heads at $0$.
\end{lemma}

\begin{definition}[Anchor-free trace]
\lean{Turing.FinTM.splitSafe}
\label{Turing.FinTM.splitSafe}
\leanok
\uses{Turing.Cfg, Turing.MultiTapeTM, Turing.MultiTapeTM.runFrom}
The property of a machine, an anchor control state, a start configuration
$c$, and a duration $t$, that none of the configurations of the run from $c$
at times $0$ through $t$ inclusive is in the anchor state.
\end{definition}

\begin{lemma}[Concatenation of anchor-free traces]
\lean{Turing.FinTM.splitSafe_add}
\label{Turing.FinTM.splitSafe_add}
\leanok
\uses{Turing.Cfg, Turing.FinTM.splitSafe, Turing.MultiTapeTM, Turing.MultiTapeTM.runFrom, Turing.MultiTapeTM.runFrom_add}
\difficulty{3}
If the run from $c$ is anchor-free for $u$ steps and the run continued from
its configuration at time $u$ is anchor-free for $v$ further steps, then the
run from $c$ is anchor-free for $u + v$ steps.
\end{lemma}

\begin{lemma}[Cutting an absorbing phase at first exit]
\lean{Turing.FinTM.splitEmbed_cut}
\label{Turing.FinTM.splitEmbed_cut}
\leanok
\uses{Turing.Action.mapState, Turing.Cfg, Turing.Cfg.mapState, Turing.FinTM.splitSafe, Turing.MultiTapeTM, Turing.MultiTapeTM.runFrom, Turing.MultiTapeTM.runFrom_add, Turing.MultiTapeTM.runFrom_succ_eq_step', Turing.MultiTapeTM.step}
\difficulty{5}
Let a source machine have a class of terminal control states on which it
stalls, and let a host machine agree, on the images of non-terminal states
under a state embedding, with the mapped source transitions, where no
embedded state equals a distinguished host anchor. If the source run from $c$
reaches a terminal configuration $d$ within $T$ steps, then there is a
$t \le T$ such that the host run from the embedded $c$ reaches the embedded
$d$ at time $t$ and the entire host trace up to time $t$ avoids the anchor.
\end{lemma}

\begin{definition}[Standalone input-rewind machine]
\lean{Turing.FinTM.splitRewindTM}
\label{Turing.FinTM.splitRewindTM}
\leanok
\uses{Turing.FinTM, Turing.FinTM.controlAction}
A three-state machine, with $k+1$ work tapes it never writes, that moves its
input head left to the left boundary and then one cell right, stalling in an
absorbing return state.
\end{definition}

\begin{definition}[Split emitter machine]
\lean{Turing.FinTM.splitEmitTM}
\label{Turing.FinTM.splitEmitTM}
\leanok
\uses{Turing.FinTM, Turing.FinTM.controlAction}
A four-state machine with $k+1$ work tapes that emits the self-delimiting
pair encoding of a split of its native input: using the first work tape only
as a length counter, it emits each native bit twice while the counter has
cells left, then emits the separator bits $0,1$, then copies the remaining
native suffix verbatim; no counter bit is ever emitted.
\end{definition}

\begin{definition}[Control phases of the split-search round]
\lean{Turing.FinTM.SplitBodyState}
\label{Turing.FinTM.SplitBodyState}
\leanok
The finite control of the combined split-search round body over a source
state type: an anchor state, a preparation phase, a counting phase carrying a
source state and an overflow flag, a check phase, an input-rewind phase
carrying the acceptance bit, a restoration phase, and an emission phase; only
restoration can return to the anchor.
\end{definition}

\begin{definition}[Combined round controller]
\lean{Turing.FinTM.splitBodyTM}
\label{Turing.FinTM.splitBodyTM}
\leanok
\uses{Turing.Action.mapState, Turing.FinTM, Turing.FinTM.SplitBodyState, Turing.FinTM.controlAction, Turing.FinTM.splitCountAction, Turing.FinTM.splitEmitTM, Turing.FinTM.splitPrepareTM, Turing.FinTM.splitRestoreTM, Turing.FinTM.splitRewindTM}
Given a source machine $M$ and a start state, the machine executing one round
of the split search: from the anchor it dispatches to preparation, runs the
embedded preparation machine, starts $M$ on the prepared bank with its
emissions counted against the native input, checks whether the consumed
length matches the input length, rewinds the input carrying the verdict, and
then either runs the split emitter (acceptance) or the rejection cleanup
followed by an explicit return to the anchor.
\end{definition}

\begin{definition}[Source bank configuration]
\lean{Turing.FinTM.splitBank}
\label{Turing.FinTM.splitBank}
\leanok
\uses{Turing.Cfg, Turing.FinTM, Turing.FinTM.catalogPolyTape}
A configuration of the source machine over the empty virtual input in which
every work tape holds the unary word of length one more than the candidate's
length, all heads are at displacement $0$, and the control state and output
are given parameters.
\end{definition}

\begin{lemma}[Initial configuration of the round body]
\lean{Turing.FinTM.splitBody_start}
\label{Turing.FinTM.splitBody_start}
\leanok
\uses{Turing.Cfg.ofWords, Turing.FinTM, Turing.FinTM.splitBodyTM, Turing.MultiTapeTM.initCfg, Turing.initCfg_ofWords, Turing.stateWord}
\difficulty{2}
On every input $w$, the initial configuration of the round body equals the
canonical anchored configuration whose stored words carry the empty
candidate.
\end{lemma}

\begin{lemma}[Full-run state embedding]
\lean{Turing.FinTM.splitEmbed_run}
\label{Turing.FinTM.splitEmbed_run}
\leanok
\uses{Turing.Action.mapState, Turing.Cfg, Turing.Cfg.mapState, Turing.MultiTapeTM, Turing.MultiTapeTM.runFrom, Turing.MultiTapeTM.runFrom_comm_of_step, Turing.MultiTapeTM.step}
\difficulty{3}
If a host machine's transitions on all embedded states equal the mapped
transitions of a source machine, then for every start configuration $c$ and
time $t$, the host run from the embedded $c$ equals the embedding of the
source run from $c$, including through halting.
\end{lemma}

\begin{lemma}[Preparation phase of the round]
\lean{Turing.FinTM.splitBody_prepare}
\label{Turing.FinTM.splitBody_prepare}
\leanok
\uses{Turing.Action.apply, Turing.Cfg.mapState, Turing.Cfg.ofWords, Turing.FinTM, Turing.FinTM.controlAction, Turing.FinTM.controlAction_apply, Turing.FinTM.splitBank, Turing.FinTM.splitBodyTM, Turing.FinTM.splitCountCfg, Turing.FinTM.splitEmbed_cut, Turing.FinTM.splitPrepareReady, Turing.FinTM.splitPrepareTM, Turing.FinTM.splitPrepare_run, Turing.FinTM.splitSafe, Turing.MultiTapeTM.runFrom, Turing.MultiTapeTM.runFrom_succ_eq_step', Turing.MultiTapeTM.step, Turing.moveInputPos_zero, Turing.stateWord}
\difficulty{5}
For any input $w$ and candidate $s$, there is a time $t \le 2(\abs{s}+1)$
such that $t+1$ steps of the round body from its preparation seam reach the
counting-host embedding of the prepared source bank at the source's start
state, and the whole trace avoids the anchor.
\end{lemma}

\begin{lemma}[Counting phase of the round]
\lean{Turing.FinTM.splitBody_count}
\label{Turing.FinTM.splitBody_count}
\leanok
\uses{Turing.FinTM, Turing.FinTM.splitBank, Turing.FinTM.splitBodyTM, Turing.FinTM.splitCountCfg, Turing.FinTM.splitCount_run, Turing.FinTM.splitSafe, Turing.MultiTapeTM.runFrom, Turing.MultiTapeTM.runFrom_add, Turing.MultiTapeTM.runFrom_of_halt}
\difficulty{4}
If the source machine, started on the prepared bank, reaches the halted bank
with output $o$ after $T$ steps, then there is a time $t \le T$ at which the
round body, started from the counting-host embedding of the start bank,
reaches the embedding of the halted bank with output $o$; the whole trace
avoids the anchor.
\end{lemma}

\begin{lemma}[Rewind phase of the round]
\lean{Turing.FinTM.splitBody_rewind}
\label{Turing.FinTM.splitBody_rewind}
\leanok
\uses{Turing.Action.apply, Turing.Cfg, Turing.Cfg.mapState, Turing.FinTM, Turing.FinTM.catalogRewind, Turing.FinTM.controlAction, Turing.FinTM.controlAction_apply, Turing.FinTM.splitBodyTM, Turing.FinTM.splitEmbed_cut, Turing.FinTM.splitRewindTM, Turing.FinTM.splitSafe, Turing.MultiTapeTM.runFrom, Turing.MultiTapeTM.step, Turing.moveInputPos_zero}
\difficulty{4}
For a configuration $c$ of the rewind machine in its start state, there is a
time $t \le i + 2$, with $i$ the current input position, such that the round
body run from the rewind-phase embedding of $c$ (carrying the acceptance bit)
reaches the embedding of $c$ with its input head back at the first cell and
in the rewind machine's return state; the whole trace avoids the anchor and
preserves the source bank and output.
\end{lemma}

\begin{lemma}[Restoration phase of the round]
\lean{Turing.FinTM.splitBody_restore}
\label{Turing.FinTM.splitBody_restore}
\leanok
\uses{Turing.Action.apply, Turing.Cfg.mapState, Turing.Cfg.ofWords, Turing.FinTM, Turing.FinTM.controlAction, Turing.FinTM.controlAction_apply, Turing.FinTM.splitBodyTM, Turing.FinTM.splitEmbed_cut, Turing.FinTM.splitRestoreScan, Turing.FinTM.splitRestoreTM, Turing.FinTM.splitRestore_run, Turing.FinTM.splitSafe, Turing.FinTM.splitStep, Turing.MultiTapeTM.runFrom, Turing.MultiTapeTM.step, Turing.moveInputPos_zero, Turing.stateWord}
\difficulty{4}
For any input $w$ and candidate $s$, there is a time
$t \le 2\abs{s} + \abs{w} + 5$ such that the round body, started from the
restoration-phase embedding of the cleanup scan, reaches the canonical
anchored-word configuration in the restoration phase's return state carrying
one split-search step applied to $s$; the whole trace avoids the anchor.
\end{lemma}

\begin{definition}[Emitter configuration]
\lean{Turing.FinTM.splitEmitCfg}
\label{Turing.FinTM.splitEmitCfg}
\leanok
\uses{Turing.Cfg, Turing.FinTM.bufferTape, Turing.FinTM.catalogPolyTape, Turing.FinTM.splitPos}
Given an input $w$, a candidate $s$, an optional control state, a count $j$,
a counter-head position $h$, and an output word, the emitter configuration
is the configuration of the split emitter with that control state, with its
input head at the saturated position for $j$, with $s$ on its first work
tape and the head there at $h$, with every scratch tape holding the unary
word of length $\abs{s}+1$ and its head at $0$, and with the given output.
\end{definition}

\begin{lemma}[Doubling pass of the split emitter]
\lean{Turing.FinTM.splitEmit_double}
\label{Turing.FinTM.splitEmit_double}
\leanok
\uses{Turing.Action.apply, Turing.Cfg.inputSymbol, Turing.Cfg.workTapeSymbols, Turing.FinTM.splitEmitCfg, Turing.FinTM.splitEmitTM, Turing.FinTM.splitPos_read, Turing.FinTM.splitPos_succ, Turing.MultiTapeTM.runFrom, Turing.MultiTapeTM.runFrom_succ_eq_step', Turing.MultiTapeTM.step, Turing.moveInputPos_zero}
\difficulty{4}
Let the candidate's length be at most $\abs{w}$. For every $j$ at most the
candidate's length, $2j$ steps of the split emitter from its start emit each
of the first $j$ bits of $w$ twice and advance both the native head and the
candidate counter to $j$; candidate bits are read only for their presence.
\end{lemma}

\begin{lemma}[Separator emission of the split emitter]
\lean{Turing.FinTM.splitEmit_separator}
\label{Turing.FinTM.splitEmit_separator}
\leanok
\uses{Turing.Action.apply, Turing.Cfg.workTapeSymbols, Turing.FinTM.splitEmitCfg, Turing.FinTM.splitEmitTM, Turing.MultiTapeTM.runFrom, Turing.MultiTapeTM.runFrom_succ_eq_step, Turing.MultiTapeTM.runFrom_zero, Turing.MultiTapeTM.step, Turing.moveInputPos_zero}
\difficulty{3}
Once the candidate counter is exhausted, two steps of the split emitter
append the separator bits $0,1$ to the output without moving the native head
from the start of the suffix.
\end{lemma}

\begin{lemma}[Suffix copy of the split emitter]
\lean{Turing.FinTM.splitEmit_suffix}
\label{Turing.FinTM.splitEmit_suffix}
\leanok
\uses{Turing.Action.apply, Turing.Cfg.inputSymbol, Turing.FinTM.controlAction, Turing.FinTM.splitEmitCfg, Turing.FinTM.splitEmitTM, Turing.FinTM.splitPos_read, Turing.FinTM.splitPos_succ, Turing.MultiTapeTM.runFrom, Turing.MultiTapeTM.runFrom_succ_eq_step, Turing.MultiTapeTM.runFrom_zero, Turing.MultiTapeTM.step}
\difficulty{4}
If the native input splits as $w = p r$, then from the split emitter's copy
state at position $\abs{p}$ with prior output $o$, exactly $\abs{r}+1$ steps
copy $r$ to the output and halt, preserving all work tapes; this includes the
empty suffix with its final blank-reading halt.
\end{lemma}

\begin{lemma}[Complete run of the split emitter]
\lean{Turing.FinTM.splitEmit_run}
\label{Turing.FinTM.splitEmit_run}
\leanok
\uses{Turing.FinTM.splitEmitCfg, Turing.FinTM.splitEmitTM, Turing.FinTM.splitEmit_double, Turing.FinTM.splitEmit_separator, Turing.FinTM.splitEmit_suffix, Turing.MultiTapeTM.runFrom, Turing.MultiTapeTM.runFrom_add, Turing.pairEncode}
\difficulty{3}
If the candidate's length $\ell$ is at most $\abs{w}$, then exactly
$\ell + \abs{w} + 3$ steps of the split emitter from its start produce the
self-delimiting pair encoding of the split of $w$ at position $\ell$ (the
first $\ell$ bits paired with the rest) and halt; the candidate may contain
any bit pattern.
\end{lemma}

\begin{lemma}[A single anchor-free step]
\lean{Turing.FinTM.splitSafe_one}
\label{Turing.FinTM.splitSafe_one}
\leanok
\uses{Turing.Cfg, Turing.FinTM.splitSafe, Turing.MultiTapeTM, Turing.MultiTapeTM.runFrom, Turing.MultiTapeTM.step}
\difficulty{2}
If one machine step carries $c$ to $d$ and neither $c$ nor $d$ is in the
anchor state, then the one-step run from $c$ ends at $d$ and is anchor-free.
\end{lemma}

\begin{lemma}[Joining two exact anchor-free runs]
\lean{Turing.FinTM.splitSafe_join}
\label{Turing.FinTM.splitSafe_join}
\leanok
\uses{Turing.Cfg, Turing.FinTM.splitSafe, Turing.FinTM.splitSafe_add, Turing.MultiTapeTM, Turing.MultiTapeTM.runFrom, Turing.MultiTapeTM.runFrom_add}
\difficulty{2}
If the run from $c$ reaches $d$ in exactly $u$ anchor-free steps and the run
from $d$ reaches $f$ in exactly $v$ anchor-free steps, then the run from $c$
reaches $f$ in exactly $u+v$ anchor-free steps.
\end{lemma}

\begin{lemma}[One complete round of the split search]
\lean{Turing.FinTM.splitBody_round}
\label{Turing.FinTM.splitBody_round}
\leanok
\uses{Turing.Action.apply, Turing.Cfg, Turing.Cfg.inputSymbol, Turing.Cfg.mapState, Turing.Cfg.ofWords, Turing.FinTM, Turing.FinTM.SplitBodyState, Turing.FinTM.bufferTape, Turing.FinTM.catalogPolyTape, Turing.FinTM.controlAction, Turing.FinTM.controlAction_apply, Turing.FinTM.splitBank, Turing.FinTM.splitBodyTM, Turing.FinTM.splitBody_count, Turing.FinTM.splitBody_prepare, Turing.FinTM.splitBody_restore, Turing.FinTM.splitBody_rewind, Turing.FinTM.splitCountCfg, Turing.FinTM.splitCount_accept, Turing.FinTM.splitEmbed_run, Turing.FinTM.splitEmitCfg, Turing.FinTM.splitEmitTM, Turing.FinTM.splitEmit_run, Turing.FinTM.splitPos, Turing.FinTM.splitRestoreScan, Turing.FinTM.splitSafe, Turing.FinTM.splitSafe_join, Turing.FinTM.splitSafe_one, Turing.FinTM.splitScratch, Turing.FinTM.splitStep, Turing.MultiTapeTM.runFrom, Turing.MultiTapeTM.runFrom_add, Turing.MultiTapeTM.runFrom_succ_eq_step', Turing.MultiTapeTM.step, Turing.moveInputPos_zero, Turing.pairEncode, Turing.stateWord}
\difficulty{6}
Suppose the source machine, started on the prepared bank for candidate $s$,
reaches the halted bank with output $o$ after $T$ steps. Then there is a time
$t$ with $0 < t \le T + 5\abs{s} + 3\abs{w} + 20$ such that the round body,
started from the canonical anchored configuration carrying $s$ on input $w$,
never revisits the anchor strictly between times $0$ and $t$, and at time
$t$: if $\abs{s} + \abs{o} = \abs{w}$ the body has halted with output the
pair encoding of the split of $w$ at position $\abs{s}$; otherwise the body
is back at the canonical anchored configuration carrying one split-search
step applied to $s$.
\end{lemma}

\begin{lemma}[Split solver from a round contract]
\lean{Turing.FinTM.splitSolve_of_body}
\label{Turing.FinTM.splitSolve_of_body}
\leanok
\uses{Turing.Cfg.ofWords, Turing.FinTM, Turing.FinTM.ComputesFunInTime, Turing.FinTM.ComputesInTime.mono, Turing.FinTM.computesFunInTime_lengthBits, Turing.FinTM.exists_loopFindTM, Turing.FinTM.splitAccept, Turing.FinTM.splitLoop_bound, Turing.FinTM.splitLoop_result, Turing.FinTM.splitStep, Turing.FinTM.splitStep_inv, Turing.MultiTapeTM.initCfg, Turing.MultiTapeTM.runFrom, Turing.pairEncode, Turing.solveSplit, Turing.stateWord}
\difficulty{5}
Let a body machine with a distinguished anchor state satisfy a startup
contract (within $A(n+1)^{e+1}$ steps it reaches the anchored empty-candidate
configuration without an earlier anchor visit) and a round contract (from the
anchored configuration carrying any candidate of length at most $n+1$, within
$A(n+1)^{e+1}$ steps and with no interior anchor visit, it either halts with
the encoded split when the candidate satisfies the length equation, or
returns to the anchor with the stepped candidate). Then there exist a machine
$M$ and a constant $c$ computing, within $c\,(n+1)^{e+2}$ steps, the function
that sends $w$ to the pair encoding of the split of $w$ at the least $i$ with
$i + C\,(i+1)^{e} = \abs{w}$, and to the empty word when no such $i$ exists.
\end{lemma}

\begin{lemma}[Polynomial source on the prepared bank]
\lean{Turing.FinTM.splitSource_poly}
\label{Turing.FinTM.splitSource_poly}
\leanok
\uses{Turing.FinTM.catalogPolyCfg, Turing.FinTM.catalogPolyCost, Turing.FinTM.catalogPolyUnaryTM, Turing.FinTM.splitBank, Turing.FinTM.splitPoly_loop_end, Turing.MultiTapeTM.runFrom}
\difficulty{1}
For all $c$ and $C$ and every candidate $s$, the unary polynomial generator
with $c+1$ loop levels and $C$ emissions per innermost visit, started in its
outermost loop state on the prepared bank for $s$, halts after exactly one
more step than the exact nest cost at side length $\abs{s}+1$, and its
captured output is a run of $C\,(\abs{s}+1)^{c+1}$ ones.
\end{lemma}

\begin{lemma}[Constant source for exponent zero]
\lean{Turing.FinTM.splitSource_constant}
\label{Turing.FinTM.splitSource_constant}
\leanok
\uses{Turing.Action.apply, Turing.Cfg.ext_zero_tapes, Turing.Cfg.inputSymbol, Turing.FinTM.catalogPrefixCfg, Turing.FinTM.catalogPrefixTM, Turing.FinTM.catalogPrefixTM_emit, Turing.FinTM.splitBank, Turing.MultiTapeTM.initCfg, Turing.MultiTapeTM.runFrom, Turing.MultiTapeTM.runFrom_succ_eq_step', Turing.MultiTapeTM.step}
\difficulty{3}
The fixed-prefix emitter for the word $1^C$, started on the empty virtual
input in the bank for candidate $s$ (it has no scratch tapes), halts after
exactly $C+1$ steps with captured output $1^C$, the final step being the
blank-reading halt.
\end{lemma}

\begin{lemma}[Common envelope for one round]
\lean{Turing.FinTM.splitBody_envelope}
\label{Turing.FinTM.splitBody_envelope}
\leanok
\difficulty{3}
For all naturals $C$, $e$, $l$, $n$, $T$ with $l \le n + 1$ and
$T \le (C + 1 + 5e)(l+1)^{e} + 1$,
\[
  T + 5l + 3n + 20 \;\le\; \bigl((C + 1 + 5e)\,2^{e} + 40\bigr)(n+1)^{e+1}.
\]
\end{lemma}

\begin{lemma}[Split solver from an exact unary source]
\lean{Turing.FinTM.splitSolve_source}
\label{Turing.FinTM.splitSolve_source}
\leanok
\uses{Turing.FinTM, Turing.FinTM.ComputesFunInTime, Turing.FinTM.splitAccept, Turing.FinTM.splitBank, Turing.FinTM.splitBodyTM, Turing.FinTM.splitBody_envelope, Turing.FinTM.splitBody_round, Turing.FinTM.splitBody_start, Turing.FinTM.splitSolve_of_body, Turing.MultiTapeTM.runFrom, Turing.pairEncode, Turing.solveSplit}
\difficulty{3}
Let a source machine, started on the prepared bank for any candidate $s$,
halt within $B(\abs{s})$ steps with captured output a run of
$C\,(\abs{s}+1)^{e}$ ones, where $B(l) \le (C+1+5e)(l+1)^{e} + 1$. Then there
exist a machine $N$ and a constant $c$ computing, within $c\,(n+1)^{e+2}$
steps, the function sending $w$ to the pair encoding of the split of $w$ at
the least $i$ with $i + C\,(i+1)^{e} = \abs{w}$, and to the empty word when
no such $i$ exists.
\end{lemma}

\begin{lemma}[Split solver, both exponent cases]
\lean{Turing.FinTM.splitSolve_closed}
\label{Turing.FinTM.splitSolve_closed}
\leanok
\uses{Turing.FinTM, Turing.FinTM.ComputesFunInTime, Turing.FinTM.catalogPolyCost, Turing.FinTM.catalogPolyCost_le, Turing.FinTM.catalogPolyUnaryTM, Turing.FinTM.catalogPrefixTM, Turing.FinTM.splitSolve_source, Turing.FinTM.splitSource_constant, Turing.FinTM.splitSource_poly, Turing.pairEncode, Turing.solveSplit}
\difficulty{3}
For all constants $C$ and $e$ there exist a machine $M$ and a constant $c$
computing, within $c\,(n+1)^{e+2}$ steps on inputs of length $n$, the
function sending $w$ to the pair encoding of the split of $w$ at the least
$i$ with $i + C\,(i+1)^{e} = \abs{w}$, and to the empty word when no such $i$
exists.
\end{lemma}

\begin{theorem}[Padding split search, machine contract]
\lean{Turing.FinTM.computesFunInTime_splitSolve}
\label{Turing.FinTM.computesFunInTime_splitSolve}
\leanok
\uses{Turing.FinTM, Turing.FinTM.ComputesFunInTime, Turing.FinTM.splitSolve_closed, Turing.pairEncode, Turing.solveSplit}
\difficulty{1}
For all constants $C$ and $e$ there exist a finite Turing machine $M$ and a
constant $c$ such that $M$ computes, within $c\,(n+1)^{e+2}$ steps on inputs
of length $n$, the function that on input $w$ searches for the least
$i \le \abs{w}$ with $i + C\,(i+1)^{e} = \abs{w}$ and, on success, outputs
the self-delimiting pair encoding of the first $i$ bits of $w$ paired with
the rest of $w$, and on failure outputs the empty word.
\end{theorem}

\begin{theorem}[Fixed-width increment, machine contract]
\lean{Turing.FinTM.computesFunInTime_incFixed}
\label{Turing.FinTM.computesFunInTime_incFixed}
\leanok
\uses{Turing.FinTM, Turing.FinTM.ComputesFunInTime, Turing.FinTM.incFixedTM, Turing.FinTM.incFixed_computes, Turing.incFixed}
\difficulty{1}
There exist a finite Turing machine $M$ and a constant $c$ such that $M$
computes, within $c\,(n+1)$ steps on inputs of length $n$, the fixed-width
binary successor of its input read least-significant bit first, with empty
output on overflow (an all-ones input).
\end{theorem}

\begin{definition}[Width-parametric split acceptance]
\lean{Turing.FinTM.emitterSplitAccept}
\label{Turing.FinTM.emitterSplitAccept}
\leanok
The acceptance test of the width-parametric split search: for a width
function $f$, a candidate $s$ is accepted for the word $w$ exactly when
$\abs{s} + f(\abs{s}) = \abs{w}$; no monotonicity of $f$ is assumed.
\end{definition}

\begin{lemma}[Orbit search agrees with the width-parametric search]
\lean{Turing.FinTM.emitterSplit_find}
\label{Turing.FinTM.emitterSplit_find}
\leanok
\uses{Turing.FinTM.catalogFind_congr, Turing.FinTM.emitterSplitAccept, Turing.FinTM.splitStep, Turing.FinTM.splitStep_orbit, Turing.solveSplitWith}
\difficulty{3}
For every function $f$ and every word $w$, searching the indices
$0, \dots, \abs{w}$ for the first $i$ at which the width-parametric
acceptance test succeeds on the $i$-th iterate of the split step (the unary
candidate $1^i$) returns exactly the width-parametric split search for $f$
at length $\abs{w}$.
\end{lemma}

\begin{lemma}[Result of the width-parametric search loop]
\lean{Turing.FinTM.emitterSplit_result}
\label{Turing.FinTM.emitterSplit_result}
\leanok
\uses{Turing.FinTM.emitterSplitAccept, Turing.FinTM.emitterSplit_find, Turing.FinTM.splitStep, Turing.FinTM.splitStep_orbit, Turing.pairEncode, Turing.solveSplitWith}
\difficulty{3}
Splitting $w$ at the length of the first orbit candidate accepted by the
width-parametric test and pair encoding the two halves gives the same word
as splitting $w$ at the index returned by the width-parametric split search
for $f$; when both searches fail, both sides are the empty word.
\end{lemma}

\begin{lemma}[Loop bound for the width-parametric search]
\lean{Turing.FinTM.emitterSplit_loop_bound}
\label{Turing.FinTM.emitterSplit_loop_bound}
\leanok
\difficulty{3}
For all naturals $c$, $A$, $n$, $t$,
\[
  c\,\bigl(A(t + n + 2) + 1\bigr)(n+2) \;\le\;
  2c\,(A+1)\,(n+1)\,(t + n + 2).
\]
\end{lemma}

\begin{lemma}[Width-parametric split solver from a round contract]
\lean{Turing.FinTM.emitterSplit_of_body}
\label{Turing.FinTM.emitterSplit_of_body}
\leanok
\uses{Turing.Cfg.ofWords, Turing.FinTM, Turing.FinTM.ComputesFunInTime, Turing.FinTM.ComputesInTime.mono, Turing.FinTM.computesFunInTime_lengthBits, Turing.FinTM.emitterSplitAccept, Turing.FinTM.emitterSplit_loop_bound, Turing.FinTM.emitterSplit_result, Turing.FinTM.exists_loopFindTM, Turing.FinTM.splitStep, Turing.FinTM.splitStep_inv, Turing.MultiTapeTM.initCfg, Turing.MultiTapeTM.runFrom, Turing.pairEncode, Turing.solveSplitWith, Turing.stateWord}
\difficulty{5}
Let a body machine with an anchor state satisfy a startup contract and a
round contract with common envelope $A\,(T_E(n+1) + n + 2)$, where the round
from a candidate $s$ of length at most $n+1$ either halts with the encoded
split of the input at $\abs{s}$ when $\abs{s} + f(\abs{s})$ equals the input
length, or returns to the anchor with the stepped candidate, never visiting
the anchor strictly inside. Then there exist a machine $M$ and a constant
$c$ computing, within $c\,(n+1)(T_E(n+1) + n + 2)$ steps, the function
sending $w$ to the pair encoding of the split of $w$ at the least $i$ with
$i + f(i) = \abs{w}$, and to the empty word when no such $i$ exists.
\end{lemma}

\begin{definition}[Idle placeholder machine]
\lean{Turing.FinTM.emitterIdleTM}
\label{Turing.FinTM.emitterIdleTM}
\leanok
\uses{Turing.FinTM, Turing.FinTM.controlAction}
A single-state machine with no work tapes, used only to supply the empty
inactive bank of the virtual-input wrapper; its own transition is never
entered by the phase that uses it.
\end{definition}

\begin{definition}[Candidate evaluator machine]
\lean{Turing.FinTM.emitterEvalTM}
\label{Turing.FinTM.emitterEvalTM}
\leanok
\uses{Turing.FinTM, Turing.FinTM.bufferedCompTM, Turing.FinTM.controlAction, Turing.FinTM.emitterIdleTM, Turing.captureAction}
Given a machine $M$, the candidate evaluator obtained by wrapping the
virtual-input phase of the buffered composition over $M$ with the capture
transformer: $M$ runs on a buffered candidate as virtual input, every
emission lands on a final capture tape rather than the physical output, and
completion passes to a live return state.
\end{definition}

\begin{definition}[Evaluator configuration]
\lean{Turing.FinTM.emitterEvalCfg}
\label{Turing.FinTM.emitterEvalCfg}
\leanok
\uses{Turing.Cfg, Turing.FinTM, Turing.FinTM.bufferedSecondCfg, Turing.FinTM.emitterEvalTM, Turing.FinTM.emitterIdleTM, Turing.captureCfg}
A configuration of the candidate evaluator in which the preserved candidate
occupies the first work tape, the simulated machine's work bank follows it,
and the final tape captures every emission of the simulated machine; the
physical input head stays at a fixed position and the physical output is
empty.
\end{definition}

\begin{lemma}[Guarded run of the candidate evaluator]
\lean{Turing.FinTM.emitter_eval_run}
\label{Turing.FinTM.emitter_eval_run}
\leanok
\uses{Turing.Cfg, Turing.Cfg.Halted, Turing.FinTM, Turing.FinTM.VirtualTag, Turing.FinTM.bufferedCompTM, Turing.FinTM.bufferedSecondCfg, Turing.FinTM.bufferedSecondCfg_run, Turing.FinTM.emitterEvalCfg, Turing.FinTM.emitterEvalTM, Turing.FinTM.emitterIdleTM, Turing.MultiTapeTM.runFrom, Turing.capture_run}
\difficulty{4}
Let the simulated machine, started from a configuration $c$ over the
candidate with a valid virtual-input tag, stay live for all times before
$t$. Then $t$ steps of the candidate evaluator from the evaluator
configuration of $c$ land exactly on the evaluator configuration of the
simulated machine's configuration at time $t$, with a valid tag.
\end{lemma}

\begin{lemma}[Initial configuration of the candidate evaluator]
\lean{Turing.FinTM.emitter_eval_initial}
\label{Turing.FinTM.emitter_eval_initial}
\leanok
\uses{Turing.Cfg.init, Turing.Cfg.ofWords, Turing.FinTM, Turing.FinTM.bufferTape, Turing.FinTM.bufferedSecondCfg, Turing.FinTM.emitterEvalCfg, Turing.FinTM.emitterEvalTM, Turing.FinTM.emitterIdleTM, Turing.FinTM.tapeBlocks, Turing.MultiTapeTM.initCfg, Turing.captureCfg, Turing.stateWord}
\difficulty{4}
The evaluator configuration of the simulated machine's initial configuration
on a candidate $s$ --- candidate buffered, source work bank blank, capture
tape empty, all heads at zero --- equals the canonical anchored configuration
whose stored words carry $s$, at the evaluator's initial virtual state.
\end{lemma}

\begin{lemma}[First halt of the candidate evaluator]
\lean{Turing.FinTM.emitter_eval_first}
\label{Turing.FinTM.emitter_eval_first}
\leanok
\uses{Turing.Cfg.init, Turing.FinTM, Turing.FinTM.ComputesInTime, Turing.FinTM.ComputesInTime.output_unique, Turing.FinTM.VirtualTag, Turing.FinTM.bufferedSecondCfg, Turing.FinTM.computesInTime_iff, Turing.FinTM.emitterEvalCfg, Turing.FinTM.emitterEvalTM, Turing.FinTM.emitter_eval_run, Turing.MultiTapeTM.initCfg, Turing.MultiTapeTM.runFrom, Turing.captureCfg}
\difficulty{5}
Let $M$, on the candidate $s$, halt within $T$ steps with output $o$. Then
there is a time $t \le T$ at which $M$'s run on $s$ has first halted, with
output $o$ and a valid virtual tag, such that the candidate evaluator never
visits its return state before time $t$ and at time $t$ sits in the
evaluator configuration of $M$'s halted configuration, its capture tape
holding exactly $o$.
\end{lemma}

\begin{lemma}[Prefix equality extended by one bit]
\lean{Turing.FinTM.emitter_take_succ_eq}
\label{Turing.FinTM.emitter_take_succ_eq}
\leanok
\difficulty{3}
For words $u, v$ and an index $j$, the length-$(j+1)$ prefixes of $u$ and
$v$ agree if and only if their length-$j$ prefixes agree and their optional
$j$-th entries agree; in particular a missing bit differs from a present
bit.
\end{lemma}

\begin{definition}[Two-word comparator machine]
\lean{Turing.FinTM.emitterCompareTM}
\label{Turing.FinTM.emitterCompareTM}
\leanok
\uses{Turing.FinTM, Turing.FinTM.controlAction}
A machine with two read-only word tapes that scans both in lockstep through
their common right blank, accumulating in its finite control the Boolean
record of whether all cells so far agree, then rewinds both heads to zero;
the phase never emits, and empty words follow the same positive-time path.
\end{definition}

\begin{definition}[Comparator configuration]
\lean{Turing.FinTM.emitterCompareCfg}
\label{Turing.FinTM.emitterCompareCfg}
\leanok
\uses{Turing.Cfg, Turing.FinTM.bufferTape, Turing.FinTM.emitterCompareTM}
Given a physical input $w$, words $u$ and $v$, an input position $p$, a
control phase $q$, a register bit $b$, and a displacement $h$, the
comparator configuration is the configuration of the two-word comparator in
state $(q, b)$ with input head at $p$, with $u$ on its first tape and $v$ on
its second, both heads at displacement $h$, and empty physical output.
\end{definition}

\begin{lemma}[No false blank inside the compared region]
\lean{Turing.FinTM.emitter_compare_nonblank}
\label{Turing.FinTM.emitter_compare_nonblank}
\leanok
\difficulty{2}
For words $u, v$ and an index $j < \max(\abs{u}, \abs{v})$, it is not the
case that both the $j$-th entry of $u$ and the $j$-th entry of $v$ are
missing.
\end{lemma}

\begin{lemma}[Forward scan of the comparator]
\lean{Turing.FinTM.emitter_compare_scan}
\label{Turing.FinTM.emitter_compare_scan}
\leanok
\uses{Turing.Action.apply, Turing.Cfg.workTapeSymbols, Turing.FinTM.bufferTape_nat, Turing.FinTM.emitterCompareCfg, Turing.FinTM.emitterCompareTM, Turing.FinTM.emitter_compare_nonblank, Turing.FinTM.emitter_take_succ_eq, Turing.MultiTapeTM.runFrom, Turing.MultiTapeTM.runFrom_succ_eq_step', Turing.MultiTapeTM.step, Turing.moveInputPos_zero}
\difficulty{4}
For every $j \le \max(\abs{u}, \abs{v})$, after $j$ forward steps of the
comparator the equality register records exactly whether the length-$j$
prefixes of $u$ and $v$ agree, with both heads at $j$; this includes
mismatches caused by unequal lengths.
\end{lemma}

\begin{lemma}[Rewind of the comparator]
\lean{Turing.FinTM.emitter_compare_rewind}
\label{Turing.FinTM.emitter_compare_rewind}
\leanok
\uses{Turing.Action.apply, Turing.Cfg.workTapeSymbols, Turing.FinTM.bufferTape_left, Turing.FinTM.bufferTape_nat, Turing.FinTM.emitterCompareCfg, Turing.FinTM.emitterCompareTM, Turing.FinTM.emitter_compare_nonblank, Turing.MultiTapeTM.runFrom, Turing.MultiTapeTM.runFrom_succ_eq_step, Turing.MultiTapeTM.runFrom_zero, Turing.MultiTapeTM.step, Turing.moveInputPos_zero}
\difficulty{4}
From the comparator's rewind state with both heads at displacement $j-1$,
for any $j \le \max(\abs{u}, \abs{v})$, exactly $j+1$ steps restore both
heads to displacement $0$ and enter the return state, preserving the
comparison register.
\end{lemma}

\begin{lemma}[Complete run of the comparator]
\lean{Turing.FinTM.emitter_compare_run}
\label{Turing.FinTM.emitter_compare_run}
\leanok
\uses{Turing.Action.apply, Turing.Cfg.workTapeSymbols, Turing.FinTM.bufferTape_nat, Turing.FinTM.emitterCompareCfg, Turing.FinTM.emitterCompareTM, Turing.FinTM.emitter_compare_rewind, Turing.FinTM.emitter_compare_scan, Turing.MultiTapeTM.runFrom, Turing.MultiTapeTM.runFrom_add, Turing.MultiTapeTM.runFrom_succ_eq_step', Turing.MultiTapeTM.step, Turing.moveInputPos_zero}
\difficulty{4}
For all words $u$ and $v$, whole-word comparison of $u$ and $v$ completes
silently in exactly $2(\max(\abs{u}, \abs{v}) + 1)$ steps, ending in the
return state with the register recording whether $u = v$, both word tapes
unchanged, and both heads back at zero.
\end{lemma}

\begin{definition}[Visited-interval marker]
\lean{Turing.FinTM.emitterInterval}
\label{Turing.FinTM.emitterInterval}
\leanok
The marker tape of a contiguous visited interval of a given left endpoint
and width: exactly the cells of the interval are marked, independently of
the simulated data.
\end{definition}

\begin{definition}[Partially cleared data tape]
\lean{Turing.FinTM.emitterCleared}
\label{Turing.FinTM.emitterCleared}
\leanok
For a data tape, a left endpoint $m$, and a count $j$, the partially cleared
tape is blank at every cell $z$ with $m \le z < m + j$ and agrees with the
data tape at every other cell.
\end{definition}

\begin{lemma}[One step of interval clearing]
\lean{Turing.FinTM.emitter_cleared_step}
\label{Turing.FinTM.emitter_cleared_step}
\leanok
\uses{Turing.FinTM.emitterCleared}
\difficulty{3}
Blanking the cell at offset $j$ from the left endpoint of a tape whose first
$j$ interval cells are already cleared yields the tape with $j+1$ interval
cells cleared.
\end{lemma}

\begin{definition}[Native interval cleaner machine]
\lean{Turing.FinTM.emitterClearTM}
\label{Turing.FinTM.emitterClearTM}
\leanok
\uses{Turing.FinTM, Turing.FinTM.controlAction}
A four-state machine with three work tapes --- data, interval marker, and
origin marker --- that clears a finite marked work interval natively: it
scans left to the marked interval's left end, clears data and marker in
lockstep moving right, rewinds using the origin marker, and erases that
marker at the origin, stalling silently in its return state. Blank holes in
the data are permitted; all heads stay aligned.
\end{definition}

\begin{definition}[Cleaner configuration]
\lean{Turing.FinTM.emitterClearCfg}
\label{Turing.FinTM.emitterClearCfg}
\leanok
\uses{Turing.Cfg, Turing.FinTM.emitterClearTM}
A configuration of the interval cleaner on physical input $w$: the three
tapes hold the data, the interval marker, and the origin marker, with all
three heads at a common displacement; the physical input head and output are
untouched.
\end{definition}

\begin{lemma}[Left scan of the cleaner]
\lean{Turing.FinTM.emitter_clear_left}
\label{Turing.FinTM.emitter_clear_left}
\leanok
\uses{Turing.Action.apply, Turing.Cfg.workTapeSymbols, Turing.FinTM.bufferTape, Turing.FinTM.emitterClearCfg, Turing.FinTM.emitterClearTM, Turing.FinTM.emitterInterval, Turing.MultiTapeTM.runFrom, Turing.MultiTapeTM.runFrom_succ_eq_step, Turing.MultiTapeTM.step, Turing.moveInputPos_zero}
\difficulty{4}
Let a marked interval have left endpoint $m$ and width $W$, and let
$j < W$. Then $j+2$ steps of the cleaner, started in its initial scan state
at offset $j$ from the left end, yield its clearing state at the left end
$m$, regardless of blank holes in the data; the data, interval marker, and
origin marker are unchanged.
\end{lemma}

\begin{lemma}[Clearing nothing changes nothing]
\lean{Turing.FinTM.emitter_cleared_zero}
\label{Turing.FinTM.emitter_cleared_zero}
\leanok
\uses{Turing.FinTM.emitterCleared}
\difficulty{1}
For every data tape and every left endpoint, the tape obtained by clearing
zero cells of the visited interval is the original data tape.
\end{lemma}

\begin{lemma}[Clearing the whole interval]
\lean{Turing.FinTM.emitter_cleared_all}
\label{Turing.FinTM.emitter_cleared_all}
\leanok
\uses{Turing.FinTM.emitterCleared}
\difficulty{2}
If the data tape is blank outside the visited interval, then clearing the
entire interval produces the all-blank tape; no assumption is made about
holes or values inside the interval.
\end{lemma}

\begin{lemma}[Right clearing scan]
\lean{Turing.FinTM.emitter_clear_scan}
\label{Turing.FinTM.emitter_clear_scan}
\leanok
\uses{Turing.Action.apply, Turing.Cfg.workTapeSymbols, Turing.FinTM.bufferTape, Turing.FinTM.emitterClearCfg, Turing.FinTM.emitterClearTM, Turing.FinTM.emitterCleared, Turing.FinTM.emitterInterval, Turing.FinTM.emitter_cleared_step, Turing.FinTM.emitter_cleared_zero, Turing.MultiTapeTM.runFrom, Turing.MultiTapeTM.runFrom_succ_eq_step', Turing.MultiTapeTM.step, Turing.moveInputPos_zero}
\difficulty{4}
Let a marked interval have left endpoint $m$ and width $W$, and let
$j \le W$. Then $j$ steps of the cleaner from its clearing state at $m$
yield the configuration in which the first $j$ cells of both the data and
the interval marker are cleared and all heads are at offset $j$, while the
origin marker is intact.
\end{lemma}

\begin{lemma}[Erasing the origin marker]
\lean{Turing.FinTM.emitter_origin_erase}
\label{Turing.FinTM.emitter_origin_erase}
\leanok
\uses{Turing.FinTM.bufferTape}
\difficulty{3}
The tape obtained by blanking cell $0$ of the tape holding the single-cell
marker word is the all-blank tape.
\end{lemma}

\begin{lemma}[Origin-guided rewind of the cleaner]
\lean{Turing.FinTM.emitter_clear_origin}
\label{Turing.FinTM.emitter_clear_origin}
\leanok
\uses{Turing.Action, Turing.Action.apply, Turing.Cfg.workTapeSymbols, Turing.FinTM.bufferTape, Turing.FinTM.bufferTape_nat, Turing.FinTM.emitterClearCfg, Turing.FinTM.emitterClearTM, Turing.FinTM.emitter_origin_erase, Turing.MultiTapeTM.runFrom, Turing.MultiTapeTM.runFrom_succ_eq_step, Turing.MultiTapeTM.step, Turing.moveInputPos_zero}
\difficulty{4}
For every $n$, if the data and interval-marker tapes are already blank and
the cleaner is in its rewind state at displacement $n$, then exactly $n+1$
steps return all heads to the origin and erase the origin marker, reaching
the return state with all three tapes blank.
\end{lemma}

\begin{lemma}[Complete native cleanup of a marked interval]
\lean{Turing.FinTM.emitter_clear_run}
\label{Turing.FinTM.emitter_clear_run}
\leanok
\uses{Turing.Action.apply, Turing.Cfg.workTapeSymbols, Turing.FinTM.bufferTape, Turing.FinTM.emitterClearCfg, Turing.FinTM.emitterClearTM, Turing.FinTM.emitterInterval, Turing.FinTM.emitter_clear_left, Turing.FinTM.emitter_clear_origin, Turing.FinTM.emitter_clear_scan, Turing.FinTM.emitter_cleared_all, Turing.MultiTapeTM.runFrom, Turing.MultiTapeTM.runFrom_add, Turing.MultiTapeTM.runFrom_succ_eq_step', Turing.MultiTapeTM.step, Turing.moveInputPos_zero}
\difficulty{5}
Let a marked interval of width $W$ contain the origin ($\mathrm{left} \le 0
< \mathrm{left} + W$), let the data vanish outside it, and start the cleaner
anywhere inside it. Then there is a time $t$ with $0 < t \le 3W + 4$ after
which the cleaner is in its return state with all three tapes blank and all
heads at the origin.
\end{lemma}

\begin{definition}[Closed visited span]
\lean{Turing.FinTM.emitterSpan}
\label{Turing.FinTM.emitterSpan}
\leanok
The marker tape of a closed interval of visited cells from a lower to an
upper endpoint; the data it tracks may have arbitrary blank holes inside the
span.
\end{definition}

\begin{lemma}[Extending a visited span]
\lean{Turing.FinTM.emitter_span_extend}
\label{Turing.FinTM.emitter_span_extend}
\leanok
\uses{Turing.FinTM.emitterSpan}
\difficulty{3}
If $\mathrm{lo} \le \mathrm{hi}$ and a head position $h$ lies within one
cell of the closed span $[\mathrm{lo}, \mathrm{hi}]$, then marking cell $h$
turns the span marker into the marker of
$[\min(\mathrm{lo}, h), \max(\mathrm{hi}, h)]$.
\end{lemma}

\begin{definition}[Three-bank tape layout]
\lean{Turing.FinTM.emitterSlots}
\label{Turing.FinTM.emitterSlots}
\leanok
Given three families $D$, $V$, $O$, each indexed by $k$ positions, the
three-bank layout is the family indexed by $k + (k + k)$ positions that
lists $D$ first, then $V$, then $O$; it is used to place simulated data,
visited-cell markers, and origin markers in corresponding slots.
\end{definition}

\begin{definition}[Tracked evaluator machine]
\lean{Turing.FinTM.emitterTrackTM}
\label{Turing.FinTM.emitterTrackTM}
\leanok
\uses{Turing.FinTM, Turing.FinTM.emitterSlots}
Given a machine $M$, the tracked evaluator that simulates $M$ using two
native steps per source step: an initialization step marks each origin in
both marker banks, the first microstep of each round performs $M$'s action
(moving the corresponding marker heads along), and the second stamps the
newly reached head cells in the visited bank before possibly halting; the
physical output is exactly $M$'s output.
\end{definition}

\begin{definition}[Tracked configuration]
\lean{Turing.FinTM.emitterTrackCfg}
\label{Turing.FinTM.emitterTrackCfg}
\leanok
\uses{Turing.Cfg, Turing.FinTM, Turing.FinTM.bufferTape, Turing.FinTM.emitterSlots, Turing.FinTM.emitterSpan, Turing.FinTM.emitterTrackTM}
The configuration of the tracked evaluator after a completed source step:
the data bank holds the source configuration's work tapes unchanged, each
visited-span marker covers the corresponding current source head, and each
origin tape carries its single marker.
\end{definition}

\begin{definition}[Intermediate stamp configuration]
\lean{Turing.FinTM.emitterTrackMid}
\label{Turing.FinTM.emitterTrackMid}
\leanok
\uses{Turing.Cfg, Turing.FinTM, Turing.FinTM.emitterTrackCfg, Turing.FinTM.emitterTrackTM}
The configuration of the tracked evaluator between the two microsteps of a
round: the source data, input position, heads, and emitted output already
reflect the source action, while the visited-interval markers are still the
previous ones.
\end{definition}

\begin{lemma}[Source-action microstep]
\lean{Turing.FinTM.emitter_track_action}
\label{Turing.FinTM.emitter_track_action}
\leanok
\uses{Turing.Action.apply, Turing.Cfg, Turing.Cfg.inputSymbol, Turing.Cfg.workTapeSymbols, Turing.FinTM, Turing.FinTM.emitterSlots, Turing.FinTM.emitterTrackCfg, Turing.FinTM.emitterTrackMid, Turing.FinTM.emitterTrackTM, Turing.MultiTapeTM.step}
\difficulty{4}
For a live source configuration $c$, one step of the tracked evaluator from
the tracked configuration of $c$ yields the intermediate stamp
configuration of the source configuration after one source step: the data
simulation is exact (including any halting emission), and both marker heads
move with the data head without writing.
\end{lemma}

\begin{lemma}[Stamping microstep]
\lean{Turing.FinTM.emitter_track_stamp}
\label{Turing.FinTM.emitter_track_stamp}
\leanok
\uses{Turing.Action.apply, Turing.Cfg, Turing.FinTM, Turing.FinTM.emitterSlots, Turing.FinTM.emitterTrackCfg, Turing.FinTM.emitterTrackMid, Turing.FinTM.emitterTrackTM, Turing.FinTM.emitter_span_extend, Turing.MultiTapeTM.step, Turing.moveInputPos_zero}
\difficulty{4}
Suppose each visited span $[\mathrm{lo}_i, \mathrm{hi}_i]$ is nonempty and
each current source head lies within one cell of it. Then one step of the
tracked evaluator from the intermediate stamp configuration marks every
current head cell, producing the tracked configuration whose spans are
extended to $[\min(\mathrm{lo}_i, h_i), \max(\mathrm{hi}_i, h_i)]$, and only
then dispatches the stored successor state.
\end{lemma}

\begin{definition}[Leftmost visited position]
\lean{Turing.FinTM.emitterLo}
\label{Turing.FinTM.emitterLo}
\leanok
\uses{Turing.FinTM, Turing.MultiTapeTM.initCfg, Turing.MultiTapeTM.runFrom}
For a machine $M$, an input $x$, and a work tape, the leftmost visited
position at time $0$ is $0$, and at time $t+1$ it is the minimum of its
value at time $t$ and the position of that work head after $t+1$ steps of
$M$ on $x$.
\end{definition}

\begin{definition}[Rightmost visited position]
\lean{Turing.FinTM.emitterHi}
\label{Turing.FinTM.emitterHi}
\leanok
\uses{Turing.FinTM, Turing.MultiTapeTM.initCfg, Turing.MultiTapeTM.runFrom}
For a machine $M$, an input $x$, and a work tape, the rightmost visited
position at time $0$ is $0$, and at time $t+1$ it is the maximum of its
value at time $t$ and the position of that work head after $t+1$ steps of
$M$ on $x$.
\end{definition}

\begin{lemma}[Extent of the visited interval]
\lean{Turing.FinTM.emitter_track_extent}
\label{Turing.FinTM.emitter_track_extent}
\leanok
\uses{Turing.Cfg.init, Turing.FinTM, Turing.FinTM.emitterHi, Turing.FinTM.emitterLo, Turing.MultiTapeTM.initCfg, Turing.MultiTapeTM.runFrom, Turing.MultiTapeTM.runFrom_succ_eq_step', Turing.MultiTapeTM.workTapePos_step_le}
\difficulty{4}
For every time $t$ and work tape $i$, the visited interval
$[\mathrm{lo}, \mathrm{hi}]$ of the run contains $0$ and the current head
position, and satisfies $-t \le \mathrm{lo}$ and $\mathrm{hi} \le t$; in
particular its width is at most $2t + 1$.
\end{lemma}

\begin{lemma}[Support of the written cells]
\lean{Turing.FinTM.emitter_track_support}
\label{Turing.FinTM.emitter_track_support}
\leanok
\uses{Turing.Action.apply, Turing.Cfg.inputSymbol, Turing.Cfg.workTapeSymbols, Turing.FinTM, Turing.FinTM.emitterHi, Turing.FinTM.emitterLo, Turing.FinTM.emitter_track_extent, Turing.MultiTapeTM.initCfg, Turing.MultiTapeTM.runFrom, Turing.MultiTapeTM.runFrom_succ_eq_step', Turing.MultiTapeTM.step}
\difficulty{4}
At every time $t$, every work-tape cell lying outside the visited interval
of its tape is blank; blank cells inside the interval are permitted.
\end{lemma}

\begin{lemma}[Zero-width span is the origin marker]
\lean{Turing.FinTM.emitter_span_zero}
\label{Turing.FinTM.emitter_span_zero}
\leanok
\uses{Turing.FinTM.bufferTape, Turing.FinTM.emitterSpan}
\difficulty{3}
The closed visited span from $0$ to $0$ equals the tape holding the
single-cell marker word.
\end{lemma}

\begin{lemma}[Initialization of the tracked evaluator]
\lean{Turing.FinTM.emitter_track_initial}
\label{Turing.FinTM.emitter_track_initial}
\leanok
\uses{Turing.Action.apply, Turing.Cfg.init, Turing.FinTM, Turing.FinTM.bufferTape_append, Turing.FinTM.bufferTape_nil, Turing.FinTM.emitterSlots, Turing.FinTM.emitterTrackCfg, Turing.FinTM.emitterTrackTM, Turing.FinTM.emitter_span_zero, Turing.MultiTapeTM.initCfg, Turing.MultiTapeTM.runFrom, Turing.MultiTapeTM.step, Turing.moveInputPos_zero}
\difficulty{4}
For every machine $M$ and input $x$, one step of the tracked evaluator for
$M$ from its initial configuration on $x$ yields the tracked configuration
of the initial configuration of $M$ on $x$ with every visited span equal to
$[0, 0]$: both origin markers are installed, the source work tapes are
blank, and the output is empty.
\end{lemma}

\begin{lemma}[Main run of the tracked evaluator]
\lean{Turing.FinTM.emitter_track_run}
\label{Turing.FinTM.emitter_track_run}
\leanok
\uses{Turing.FinTM, Turing.FinTM.emitterHi, Turing.FinTM.emitterLo, Turing.FinTM.emitterTrackCfg, Turing.FinTM.emitterTrackTM, Turing.FinTM.emitter_track_action, Turing.FinTM.emitter_track_extent, Turing.FinTM.emitter_track_initial, Turing.FinTM.emitter_track_stamp, Turing.MultiTapeTM.initCfg, Turing.MultiTapeTM.runFrom, Turing.MultiTapeTM.runFrom_add, Turing.MultiTapeTM.runFrom_of_halt, Turing.MultiTapeTM.runFrom_succ_eq_step', Turing.MultiTapeTM.step, Turing.MultiTapeTM.step_of_halt, Turing.MultiTapeTM.workTapePos_step_le}
\difficulty{5}
For every time $t$, the tracked evaluator's configuration after $1 + 2t$
native steps is exactly the tracked configuration of the source machine's
configuration after $t$ steps, with visited intervals given by the leftmost
and rightmost visited positions through time $t$; this holds also after the
source has halted.
\end{lemma}

\begin{lemma}[Cost of tracking]
\lean{Turing.FinTM.emitter_track_computes}
\label{Turing.FinTM.emitter_track_computes}
\leanok
\uses{Turing.FinTM, Turing.FinTM.ComputesInTime, Turing.FinTM.computesInTime_iff, Turing.FinTM.emitterTrackCfg, Turing.FinTM.emitterTrackTM, Turing.FinTM.emitter_track_run}
\difficulty{2}
If a machine halts on input $x$ with output $o$ within $T$ steps, then its
tracked evaluator halts on $x$ with the same output $o$ within $1 + 2T$
steps.
\end{lemma}

\begin{lemma}[Spans as cleaner intervals]
\lean{Turing.FinTM.emitter_span_interval}
\label{Turing.FinTM.emitter_span_interval}
\leanok
\uses{Turing.FinTM.emitterInterval, Turing.FinTM.emitterSpan}
\difficulty{3}
For $\mathrm{lo} \le \mathrm{hi}$, the closed visited span
$[\mathrm{lo}, \mathrm{hi}]$ equals the cleaner's half-open interval marker
with left endpoint $\mathrm{lo}$ and width $\mathrm{hi} - \mathrm{lo} + 1$.
\end{lemma}

\begin{lemma}[Tracked tapes meet the cleaner contract]
\lean{Turing.FinTM.emitter_track_clearable}
\label{Turing.FinTM.emitter_track_clearable}
\leanok
\uses{Turing.FinTM, Turing.FinTM.bufferTape, Turing.FinTM.emitterClearCfg, Turing.FinTM.emitterClearTM, Turing.FinTM.emitterHi, Turing.FinTM.emitterLo, Turing.FinTM.emitterSpan, Turing.FinTM.emitter_clear_run, Turing.FinTM.emitter_span_interval, Turing.FinTM.emitter_track_extent, Turing.FinTM.emitter_track_support, Turing.MultiTapeTM.initCfg, Turing.MultiTapeTM.runFrom}
\difficulty{4}
For each work tape of a source machine run for $T$ steps, there is a time
$t$ with $0 < t \le 6T + 7$ after which the interval cleaner, started on
that tape's data together with its tracked visited span and origin marker at
the tape's current head position, reaches its return state with all three
tapes blank and all heads at the origin.
\end{lemma}

\begin{lemma}[First entry into a stopping class]
\lean{Turing.FinTM.emitter_first_entry}
\label{Turing.FinTM.emitter_first_entry}
\leanok
\uses{Turing.Cfg, Turing.MultiTapeTM, Turing.MultiTapeTM.runFrom, Turing.MultiTapeTM.runFrom_add, Turing.MultiTapeTM.step}
\difficulty{4}
Let a class of control states be absorbing (any configuration whose state
lies in the class is a fixed point of the step function), and suppose the
run from $c$ reaches a configuration $d$ with state in the class after $T$
steps. Then there is a time $t \le T$ with no earlier visit to the class
such that the configuration at time $t$ is already $d$.
\end{lemma}

\begin{lemma}[First return of the comparator]
\lean{Turing.FinTM.emitter_compare_first}
\label{Turing.FinTM.emitter_compare_first}
\leanok
\uses{Turing.Action.apply, Turing.FinTM.controlAction, Turing.FinTM.controlAction_apply, Turing.FinTM.emitterCompareCfg, Turing.FinTM.emitterCompareTM, Turing.FinTM.emitter_compare_run, Turing.FinTM.emitter_first_entry, Turing.MultiTapeTM.runFrom, Turing.MultiTapeTM.step, Turing.moveInputPos_zero}
\difficulty{4}
For words $u, v$, there is a time $t$ with
$0 < t \le 2(\max(\abs{u},\abs{v}) + 1)$ such that the comparator never
visits either Boolean return state before time $t$ and at time $t$ sits in
its return configuration with the register recording whether $u = v$ and
both heads at zero.
\end{lemma}

\begin{lemma}[First return of a tracked-tape cleanup]
\lean{Turing.FinTM.emitter_clear_first}
\label{Turing.FinTM.emitter_clear_first}
\leanok
\uses{Turing.Action.apply, Turing.FinTM, Turing.FinTM.bufferTape, Turing.FinTM.controlAction, Turing.FinTM.controlAction_apply, Turing.FinTM.emitterClearCfg, Turing.FinTM.emitterClearTM, Turing.FinTM.emitterHi, Turing.FinTM.emitterLo, Turing.FinTM.emitterSpan, Turing.FinTM.emitter_first_entry, Turing.FinTM.emitter_track_clearable, Turing.MultiTapeTM.initCfg, Turing.MultiTapeTM.runFrom, Turing.MultiTapeTM.step, Turing.moveInputPos_zero}
\difficulty{4}
For each work tape of a source machine run for $T$ steps, there is a time
$t$ with $0 < t \le 6T + 7$ such that the interval cleaner, started on that
tape's data with its tracked span and origin marker, never visits its return
state before time $t$ and at time $t$ has reached the all-blank return
configuration.
\end{lemma}

\begin{lemma}[Injectivity of the binary representation]
\lean{Turing.FinTM.emitter_bits_injective}
\label{Turing.FinTM.emitter_bits_injective}
\leanok
\difficulty{4}
The map sending a natural number to its little-endian binary representation
is injective: distinct natural numbers have distinct binary digit words.
\end{lemma}

\begin{definition}[Selecting one bank triple]
\lean{Turing.FinTM.emitterBankSymbols}
\label{Turing.FinTM.emitterBankSymbols}
\leanok
For a work-tape index $i$, the bank-triple projection sends a symbol vector
over the tracked three-bank layout to the triple of its data, visited-marker,
and origin-marker symbols at index $i$.
\end{definition}

\begin{definition}[One component of simultaneous cleanup]
\lean{Turing.FinTM.emitterBankPart}
\label{Turing.FinTM.emitterBankPart}
\leanok
\uses{Turing.Action, Turing.FinTM.controlAction, Turing.FinTM.emitterBankSymbols, Turing.FinTM.emitterClearTM}
The action performed by one component of the simultaneous cleanup: a live
component takes its interval cleaner's transition on its own triple of
symbols, and a stopped component performs the stationary silent action; the
native input is ignored.
\end{definition}

\begin{definition}[Simultaneous cleaner bank machine]
\lean{Turing.FinTM.emitterBankTM}
\label{Turing.FinTM.emitterBankTM}
\leanok
\uses{Turing.FinTM, Turing.FinTM.emitterBankPart, Turing.FinTM.emitterSlots}
The machine running $k$ interval cleaners simultaneously on disjoint
data/marker/origin triples: its finite control stores the vector of all
component cleaner states, and no source-time clock is present.
\end{definition}

\begin{definition}[Assembled bank configuration]
\lean{Turing.FinTM.emitterBankCfg}
\label{Turing.FinTM.emitterBankCfg}
\leanok
\uses{Turing.Cfg, Turing.FinTM.emitterBankTM, Turing.FinTM.emitterSlots}
Given $k$ interval-cleaner configurations on a physical input $w$ and an
input-head position $p$, the assembled bank configuration is the
configuration of the simultaneous cleaner bank whose control state is the
vector of the $k$ cleaner states, whose data, visited-marker, and
origin-marker banks carry the corresponding tapes and head positions of the
$k$ cleaners, whose input head is at $p$, and whose physical output is empty.
\end{definition}

\begin{lemma}[Component projection of the bank]
\lean{Turing.FinTM.emitterBank_part}
\label{Turing.FinTM.emitterBank_part}
\leanok
\uses{Turing.Cfg, Turing.Cfg.inputSymbol, Turing.Cfg.workTapeSymbols, Turing.FinTM.controlAction, Turing.FinTM.emitterBankCfg, Turing.FinTM.emitterBankPart, Turing.FinTM.emitterBankSymbols, Turing.FinTM.emitterClearTM, Turing.FinTM.emitterSlots}
\difficulty{3}
Selecting component $i$ of the assembled bank configuration recovers exactly
the $i$-th cleaner's transition on its own configuration, including all
three independently positioned work heads; a halted component yields the
stationary action.
\end{lemma}

\begin{lemma}[Bank step is componentwise]
\lean{Turing.FinTM.emitterBank_step}
\label{Turing.FinTM.emitterBank_step}
\leanok
\uses{Turing.Action.apply, Turing.Cfg, Turing.Cfg.inputSymbol, Turing.Cfg.workTapeSymbols, Turing.FinTM.controlAction, Turing.FinTM.emitterBankCfg, Turing.FinTM.emitterBankPart, Turing.FinTM.emitterBankTM, Turing.FinTM.emitterBank_part, Turing.FinTM.emitterClearTM, Turing.FinTM.emitterSlots, Turing.MultiTapeTM.step, Turing.moveInputPos_zero}
\difficulty{4}
One step of the simultaneous cleaner bank from an assembled configuration
equals the assembly of one step of every component cleaner; the physical
input and output are preserved.
\end{lemma}

\begin{lemma}[Bank run is componentwise]
\lean{Turing.FinTM.emitterBank_run}
\label{Turing.FinTM.emitterBank_run}
\leanok
\uses{Turing.Cfg, Turing.FinTM.emitterBankCfg, Turing.FinTM.emitterBankTM, Turing.FinTM.emitterBank_step, Turing.FinTM.emitterClearTM, Turing.MultiTapeTM.runFrom, Turing.MultiTapeTM.runFrom_succ_eq_step'}
\difficulty{3}
For every time $t$, the simultaneous cleaner bank's run from an assembled
configuration equals the assembly of each component cleaner's $t$-step run.
\end{lemma}

\begin{lemma}[A returned cleaner is absorbing]
\lean{Turing.FinTM.emitterClear_fixed}
\label{Turing.FinTM.emitterClear_fixed}
\leanok
\uses{Turing.Action.apply, Turing.Cfg, Turing.FinTM.controlAction, Turing.FinTM.controlAction_apply, Turing.FinTM.emitterClearTM, Turing.MultiTapeTM.step, Turing.moveInputPos_zero}
\difficulty{2}
Any configuration of the interval cleaner in its return state is a fixed
point of the step function: every field is preserved by further steps.
\end{lemma}

\begin{lemma}[Simultaneous cleanup of the tracked bank]
\lean{Turing.FinTM.emitterBank_clear}
\label{Turing.FinTM.emitterBank_clear}
\leanok
\uses{Turing.FinTM, Turing.FinTM.bufferTape, Turing.FinTM.emitterBankCfg, Turing.FinTM.emitterBankTM, Turing.FinTM.emitterBank_run, Turing.FinTM.emitterClearCfg, Turing.FinTM.emitterClear_fixed, Turing.FinTM.emitterHi, Turing.FinTM.emitterLo, Turing.FinTM.emitterSpan, Turing.FinTM.emitter_track_clearable, Turing.MultiTapeTM.initCfg, Turing.MultiTapeTM.runFrom, Turing.MultiTapeTM.runFrom_add}
\difficulty{3}
After a source machine has run for $T$ steps, exactly $6T + 7$ steps of the
simultaneous cleaner bank, started on the assembly of all tracked tapes with
their spans and origin markers, clear the entire bank: all data, marker, and
origin tapes become blank with every work head at zero, while the physical
input head and output are unchanged; this includes a source with no work
tapes.
\end{lemma}

\begin{lemma}[A completed bank is absorbing]
\lean{Turing.FinTM.emitterBank_fixed}
\label{Turing.FinTM.emitterBank_fixed}
\leanok
\uses{Turing.Action.apply, Turing.Cfg, Turing.FinTM.controlAction, Turing.FinTM.emitterBankPart, Turing.FinTM.emitterBankTM, Turing.FinTM.emitterClearTM, Turing.FinTM.emitterSlots, Turing.MultiTapeTM.step, Turing.moveInputPos_zero}
\difficulty{3}
Any configuration of the simultaneous cleaner bank whose every component is
in its return state is a fixed point of the step function: no tape, head, or
control field changes on further steps.
\end{lemma}

\begin{lemma}[First return of the simultaneous bank]
\lean{Turing.FinTM.emitterBank_first}
\label{Turing.FinTM.emitterBank_first}
\leanok
\uses{Turing.FinTM, Turing.FinTM.bufferTape, Turing.FinTM.catalogFirstEntry, Turing.FinTM.emitterBankCfg, Turing.FinTM.emitterBankTM, Turing.FinTM.emitterBank_clear, Turing.FinTM.emitterBank_fixed, Turing.FinTM.emitterClearCfg, Turing.FinTM.emitterHi, Turing.FinTM.emitterLo, Turing.FinTM.emitterSpan, Turing.MultiTapeTM.initCfg, Turing.MultiTapeTM.runFrom}
\difficulty{2}
After a source machine has run for $T$ steps, there is a time
$t \le 6T + 7$ such that the simultaneous cleaner bank, started on the
assembly of all tracked tapes, never shows the all-returned state vector
before time $t$ and at time $t$ has reached the fully blank assembled
endpoint; a bank over zero tapes may return at time zero.
\end{lemma}

\begin{definition}[Right-boundary normalization wrapper]
\lean{Turing.FinTM.emitterRightTM}
\label{Turing.FinTM.emitterRightTM}
\leanok
\uses{Turing.FinTM, Turing.FinTM.controlAction}
Given a machine $M$, the right-boundary wrapper is the machine that
simulates $M$ exactly and, after the actual halt of $M$, moves the input
head to the right input boundary without altering the output or work tapes
of $M$; the initial positive move of the scan handles the two distinct
blanks of an empty input correctly.
\end{definition}

\begin{definition}[Wrapper configuration before the scan]
\lean{Turing.FinTM.emitterRightCfg}
\label{Turing.FinTM.emitterRightCfg}
\leanok
\uses{Turing.Cfg, Turing.FinTM, Turing.FinTM.emitterRightTM}
The right-boundary wrapper's configuration corresponding to a source
configuration: all fields are preserved, with a halted source state replaced
by a live administrative state.
\end{definition}

\begin{lemma}[Wrapper follows live source steps]
\lean{Turing.FinTM.emitter_right_step}
\label{Turing.FinTM.emitter_right_step}
\leanok
\uses{Turing.Cfg, Turing.FinTM, Turing.FinTM.emitterRightCfg, Turing.FinTM.emitterRightTM, Turing.MultiTapeTM.step}
\difficulty{3}
For a live source configuration $c$, one step of the right-boundary wrapper
from the wrapper configuration of $c$ equals the wrapper configuration of
the source configuration after one step, including any halting emission;
only the control-state encoding differs.
\end{lemma}

\begin{lemma}[Wrapper run up to the source halt]
\lean{Turing.FinTM.emitter_right_run}
\label{Turing.FinTM.emitter_right_run}
\leanok
\uses{Turing.Cfg, Turing.FinTM, Turing.FinTM.emitterRightCfg, Turing.FinTM.emitterRightTM, Turing.FinTM.emitter_right_step, Turing.MultiTapeTM.runFrom, Turing.MultiTapeTM.runFrom_succ_eq_step'}
\difficulty{3}
If the source stays live for all times before $t$, then $t$ steps of the
right-boundary wrapper from the wrapper configuration of $c$ equal the
wrapper configuration of the source run at time $t$.
\end{lemma}

\begin{definition}[Rightward scan configuration]
\lean{Turing.FinTM.emitterRightScan}
\label{Turing.FinTM.emitterRightScan}
\leanok
\uses{Turing.Cfg, Turing.FinTM, Turing.FinTM.emitterRightTM}
A configuration of the right-boundary wrapper during its final scan: the
completed source's work tapes and output are carried verbatim, and only the
physical input position varies.
\end{definition}

\begin{lemma}[Rightward scan to the boundary]
\lean{Turing.FinTM.emitter_right_scan}
\label{Turing.FinTM.emitter_right_scan}
\leanok
\uses{Turing.Action.apply, Turing.Cfg, Turing.Cfg.inputSymbol, Turing.FinTM, Turing.FinTM.controlAction, Turing.FinTM.controlAction_apply, Turing.FinTM.emitterRightScan, Turing.FinTM.emitterRightTM, Turing.MultiTapeTM.runFrom, Turing.MultiTapeTM.runFrom_succ_eq_step, Turing.MultiTapeTM.runFrom_zero, Turing.MultiTapeTM.step, Turing.inputSymbolInner, Turing.moveInputPos_pos_of_ne_right, Turing.moveInputPos_zero}
\difficulty{4}
Let $x$ be an input and let $c$ be a source configuration on $x$. For all
$j$ and $r$ with $j + r = \abs{x}$, exactly $r + 1$ silent steps of the
right-boundary wrapper from the scan configuration of $c$ at position $j$
reach the halted scan configuration of $c$ at position $\abs{x}$; the scan
never mistakes a blank work cell for the end of the input.
\end{lemma}

\begin{lemma}[Entering the rightward scan]
\lean{Turing.FinTM.emitter_right_finish}
\label{Turing.FinTM.emitter_right_finish}
\leanok
\uses{Turing.Action.apply, Turing.Cfg, Turing.FinTM, Turing.FinTM.controlAction, Turing.FinTM.controlAction_apply, Turing.FinTM.emitterRightCfg, Turing.FinTM.emitterRightScan, Turing.FinTM.emitterRightTM, Turing.FinTM.emitter_right_scan, Turing.MultiTapeTM.runFrom, Turing.MultiTapeTM.runFrom_add, Turing.MultiTapeTM.step, Turing.moveInputPos}
\difficulty{4}
For a halted source configuration $c$, there is a time $t \le \abs{x} + 2$
after which the right-boundary wrapper, started from the wrapper
configuration of $c$, has reached the halted scan configuration at the right
boundary; neither the mandatory initial positive move nor the scan changes
the completed work tapes or output.
\end{lemma}

\begin{lemma}[Endpoint of right-boundary normalization]
\lean{Turing.FinTM.emitter_right_endpoint}
\label{Turing.FinTM.emitter_right_endpoint}
\leanok
\uses{Turing.FinTM, Turing.FinTM.emitterRightCfg, Turing.FinTM.emitterRightScan, Turing.FinTM.emitterRightTM, Turing.FinTM.emitter_right_finish, Turing.FinTM.emitter_right_run, Turing.MultiTapeTM.initCfg, Turing.MultiTapeTM.runFrom, Turing.MultiTapeTM.runFrom_add, Turing.MultiTapeTM.runFrom_of_halt}
\difficulty{4}
If a source machine has halted on input $x$ by time $T$, then
$T + \abs{x} + 2$ steps of its right-boundary wrapper from the initial
configuration reach the halted scan configuration carrying the complete
source configuration at time $T$ --- work tapes included --- with the input
head at the right boundary.
\end{lemma}

\begin{lemma}[Cost of right-boundary normalization]
\lean{Turing.FinTM.emitter_right_computes}
\label{Turing.FinTM.emitter_right_computes}
\leanok
\uses{Turing.FinTM, Turing.FinTM.ComputesInTime, Turing.FinTM.computesInTime_iff, Turing.FinTM.emitterRightTM, Turing.FinTM.emitter_right_endpoint, Turing.MultiTapeTM.initCfg, Turing.MultiTapeTM.runFrom}
\difficulty{2}
If a machine halts on input $x$ with output $o$ within $T$ steps, then its
right-boundary wrapper halts on $x$ with output $o$ within
$T + \abs{x} + 2$ steps, and at that time its input head sits at position
$\abs{x} + 1$.
\end{lemma}

\begin{lemma}[First return of the prepared evaluation]
\lean{Turing.FinTM.emitter_prepared_eval_first}
\label{Turing.FinTM.emitter_prepared_eval_first}
\leanok
\uses{Turing.Cfg.init, Turing.FinTM, Turing.FinTM.ComputesInTime, Turing.FinTM.computesInTime_iff, Turing.FinTM.emitterEvalCfg, Turing.FinTM.emitterEvalTM, Turing.FinTM.emitterHi, Turing.FinTM.emitterLo, Turing.FinTM.emitterRightScan, Turing.FinTM.emitterRightTM, Turing.FinTM.emitterTrackCfg, Turing.FinTM.emitterTrackTM, Turing.FinTM.emitter_eval_first, Turing.FinTM.emitter_right_computes, Turing.FinTM.emitter_right_endpoint, Turing.FinTM.emitter_track_computes, Turing.FinTM.emitter_track_run, Turing.MultiTapeTM.initCfg, Turing.MultiTapeTM.runFrom, Turing.MultiTapeTM.runFrom_add, Turing.MultiTapeTM.runFrom_of_halt, Turing.MultiTapeTM.runFrom_zero}
\difficulty{5}
Let $M$, on the candidate $s$, halt within $T$ steps with output $o$, and
let the composite evaluator capture the tracked, right-normalized simulation
of $M$. Then there is a time $t$ with
$0 < t \le 1 + 2T + \abs{s} + 2$ such that the composite never visits its
return state before $t$ and at time $t$ sits in the evaluator configuration
carrying the full tracked trace banks, the complete captured output, the
candidate head at the right boundary, nothing on the physical output, and
the physical input head fixed at position one.
\end{lemma}

\begin{lemma}[Binary comparison decides the width equation]
\lean{Turing.FinTM.emitter_binary_check}
\label{Turing.FinTM.emitter_binary_check}
\leanok
\uses{Turing.FinTM.emitter_bits_injective}
\difficulty{1}
For a candidate $s$ with $\abs{s} \le \abs{w}$, the binary representation of
$f(\abs{s})$ equals the binary representation of the length of the suffix of
$w$ after position $\abs{s}$ if and only if $\abs{s} + f(\abs{s}) = \abs{w}$;
this includes the case of two empty words.
\end{lemma}

\begin{lemma}[Budget of the width evaluator]
\lean{Turing.FinTM.emitter_width_budget}
\label{Turing.FinTM.emitter_width_budget}
\leanok
\uses{Turing.FinTM, Turing.FinTM.ComputesFunInTime, Turing.FinTM.ComputesInTime, Turing.FinTM.computesInTime_iff, Turing.MultiTapeTM.output_length_le}
\difficulty{2}
Let a machine $E$ compute the binary width word $s \mapsto
\mathrm{bits}(f(\abs{s}))$ within a monotone time bound $T_E$. Then for any
candidate $s$ of length at most $\abs{w}+1$: $E$ halts on $s$ with that
output within $T_E(\abs{s})$ steps, the output's length is at most
$T_E(\abs{s})$, and $T_E(\abs{s}) \le T_E(\abs{w}+1)$.
\end{lemma}

\begin{lemma}[First return of the width evaluation]
\lean{Turing.FinTM.emitter_width_eval_first}
\label{Turing.FinTM.emitter_width_eval_first}
\leanok
\uses{Turing.FinTM, Turing.FinTM.ComputesFunInTime, Turing.FinTM.emitterEvalCfg, Turing.FinTM.emitterEvalTM, Turing.FinTM.emitterHi, Turing.FinTM.emitterLo, Turing.FinTM.emitterRightScan, Turing.FinTM.emitterRightTM, Turing.FinTM.emitterTrackCfg, Turing.FinTM.emitterTrackTM, Turing.FinTM.emitter_prepared_eval_first, Turing.MultiTapeTM.initCfg, Turing.MultiTapeTM.runFrom}
\difficulty{3}
Let a machine compute the binary width word within a monotone bound $T_E$,
and let the candidate $s$ have length at most $\abs{w} + 1$. Then the
composite evaluator has a first return at some positive time
$t \le 3(T_E(\abs{w}+1) + \abs{w} + 2)$, at which it retains its tracked
work banks and the exact captured width word.
\end{lemma}

\begin{definition}[Relocating an action to host slots]
\lean{Turing.FinTM.emitterP2Action}
\label{Turing.FinTM.emitterP2Action}
\leanok
\uses{Turing.Action}
The transform relocating an action of a $k$-tape machine to an arbitrary
fixed set of slots among $l$ host tapes, specified by a partial selection
map: selected tapes perform the original writes and moves, and every
unselected tape is stationary.
\end{definition}

\begin{definition}[Relocated configuration]
\lean{Turing.FinTM.emitterP2Cfg}
\label{Turing.FinTM.emitterP2Cfg}
\leanok
\uses{Turing.Cfg}
The host configuration embedding a $k$-tape configuration into $l$ host
tapes via a selection map: selected slots carry the source tapes and heads,
unselected slots keep given contents and head positions, and the output is
the source's physical output.
\end{definition}

\begin{lemma}[Relocation commutes with one action]
\lean{Turing.FinTM.emitterP2_apply}
\label{Turing.FinTM.emitterP2_apply}
\leanok
\uses{Turing.Action, Turing.Action.apply, Turing.Cfg, Turing.FinTM.emitterP2Action, Turing.FinTM.emitterP2Cfg}
\difficulty{3}
Applying a relocated action to a relocated configuration equals relocating
the result of applying the original action: the write, head motion, and
emission are reproduced on the selected slots, and unselected tape contents
and positions are fixed.
\end{lemma}

\begin{lemma}[Guarded relocated run]
\lean{Turing.FinTM.emitterP2_relocate_run}
\label{Turing.FinTM.emitterP2_relocate_run}
\leanok
\uses{Turing.Cfg, Turing.Cfg.workTapeSymbols, Turing.FinTM.emitterP2Action, Turing.FinTM.emitterP2Cfg, Turing.FinTM.emitterP2_apply, Turing.MultiTapeTM, Turing.MultiTapeTM.runFrom, Turing.MultiTapeTM.runFrom_succ_eq_step', Turing.MultiTapeTM.step, Turing.MultiTapeTM.step_of_halt}
\difficulty{4}
Let the host's transitions on embedded good states equal the relocated
source transitions, with the selection map a left inverse of the tape
indexing. If every state visited by the source before time $t$ is good, then
$t$ host steps from the relocated configuration equal the relocation of the
source's $t$-step run.
\end{lemma}

\begin{definition}[Administrative-word eraser]
\lean{Turing.FinTM.emitterP2EraseTM}
\label{Turing.FinTM.emitterP2EraseTM}
\leanok
\uses{Turing.FinTM, Turing.FinTM.controlAction}
A three-state, one-work-tape machine that erases one contiguous
administrative word: it scans the word to its right blank, erases it right
to left, and crosses the left blank once to restore the head to the origin,
stalling in an absorbing return state; blanks may serve as word boundaries
because the word is complete and canonical.
\end{definition}

\begin{definition}[Eraser configuration]
\lean{Turing.FinTM.emitterP2EraseCfg}
\label{Turing.FinTM.emitterP2EraseCfg}
\leanok
\uses{Turing.Cfg, Turing.FinTM.bufferTape}
Given a physical input $w$, a word $u$, a control state $q$, and a
displacement $h$, the eraser configuration is the configuration of the
administrative-word eraser in state $q$ with $u$ on its single tape, its
head at displacement $h$, the physical input head on the first input cell,
and empty output.
\end{definition}

\begin{lemma}[Forward scan of the eraser]
\lean{Turing.FinTM.emitterP2_erase_scan}
\label{Turing.FinTM.emitterP2_erase_scan}
\leanok
\uses{Turing.Action.apply, Turing.Cfg.workTapeSymbols, Turing.FinTM.bufferTape_nat, Turing.FinTM.emitterP2EraseCfg, Turing.FinTM.emitterP2EraseTM, Turing.MultiTapeTM.runFrom, Turing.MultiTapeTM.runFrom_succ_eq_step', Turing.MultiTapeTM.step, Turing.moveInputPos_zero}
\difficulty{3}
For every $j \le \abs{u}$, $j$ steps of the eraser from the start of the
intact word $u$ move its head to displacement $j$, preserving the word and
emitting nothing.
\end{lemma}

\begin{lemma}[Backward erasure]
\lean{Turing.FinTM.emitterP2_erase_back}
\label{Turing.FinTM.emitterP2_erase_back}
\leanok
\uses{Turing.Action.apply, Turing.Cfg.workTapeSymbols, Turing.FinTM.bufferTape_nat, Turing.FinTM.bufferTape_nil, Turing.FinTM.catalogBuffer_erase, Turing.FinTM.emitterP2EraseCfg, Turing.FinTM.emitterP2EraseTM, Turing.MultiTapeTM.runFrom, Turing.MultiTapeTM.runFrom_succ_eq_step, Turing.MultiTapeTM.runFrom_zero, Turing.MultiTapeTM.step, Turing.moveInputPos_zero}
\difficulty{4}
From the eraser's backward state with the word $u$ and head at displacement
$\abs{u} - 1$, exactly $\abs{u} + 1$ steps erase every cell of $u$ right to
left and restore the head to the origin, reaching the return state with a
blank tape.
\end{lemma}

\begin{lemma}[Complete erasure run]
\lean{Turing.FinTM.emitterP2_erase_run}
\label{Turing.FinTM.emitterP2_erase_run}
\leanok
\uses{Turing.Action.apply, Turing.Cfg.workTapeSymbols, Turing.FinTM.bufferTape_nat, Turing.FinTM.emitterP2EraseCfg, Turing.FinTM.emitterP2EraseTM, Turing.FinTM.emitterP2_erase_back, Turing.FinTM.emitterP2_erase_scan, Turing.MultiTapeTM.runFrom, Turing.MultiTapeTM.runFrom_add, Turing.MultiTapeTM.runFrom_succ_eq_step', Turing.MultiTapeTM.step, Turing.moveInputPos_zero}
\difficulty{3}
Erasing a complete administrative word $u$ costs exactly $2(\abs{u}+1)$
steps of the eraser, after which the tape is blank, the head and native
input position are restored, and nothing has been emitted.
\end{lemma}

\begin{lemma}[First return of the eraser]
\lean{Turing.FinTM.emitterP2_erase_first}
\label{Turing.FinTM.emitterP2_erase_first}
\leanok
\uses{Turing.Action.apply, Turing.FinTM.controlAction, Turing.FinTM.controlAction_apply, Turing.FinTM.emitterP2EraseCfg, Turing.FinTM.emitterP2EraseTM, Turing.FinTM.emitterP2_erase_run, Turing.FinTM.emitter_first_entry, Turing.MultiTapeTM.runFrom, Turing.MultiTapeTM.step, Turing.moveInputPos_zero}
\difficulty{3}
For any word $u$, there is a time $t$ with $0 < t \le 2(\abs{u}+1)$ such
that the eraser never visits its return state before $t$ and at time $t$
has reached the blank-tape return configuration; the return is positive
even for the empty word.
\end{lemma}

\begin{definition}[Candidate-and-suffix preparation machine]
\lean{Turing.FinTM.emitterP2PrepareTM}
\label{Turing.FinTM.emitterP2PrepareTM}
\leanok
\uses{Turing.FinTM, Turing.FinTM.controlAction}
A three-tape machine with nine control phases that prepares the arguments of
the width evaluation: it copies the candidate from its first tape onto a
fresh second tape while advancing the native input head with saturation,
copies the remaining native suffix onto a fresh third tape, rewinds all
work heads along the preserved words and then rewinds the native input; a
finite flag records only whether the candidate ran past the native right
boundary.
\end{definition}

\begin{definition}[Preparation frame]
\lean{Turing.FinTM.emitterP2PrepareCfg}
\label{Turing.FinTM.emitterP2PrepareCfg}
\leanok
\uses{Turing.Cfg, Turing.FinTM.bufferTape, Turing.FinTM.emitterP2PrepareTM}
A configuration of the candidate-and-suffix preparation machine: the three
tapes hold the candidate, its copy, and the scanned native suffix, with the
candidate and copy heads aligned, a separate suffix head, a control phase
with overflow flag, and a given native input position.
\end{definition}

\begin{lemma}[Copying the candidate]
\lean{Turing.FinTM.emitterP2_prepare_candidate}
\label{Turing.FinTM.emitterP2_prepare_candidate}
\leanok
\uses{Turing.Action.apply, Turing.Cfg.inputSymbol, Turing.Cfg.workTapeSymbols, Turing.FinTM.bufferTape, Turing.FinTM.bufferTape_append, Turing.FinTM.bufferTape_nat, Turing.FinTM.emitterP2PrepareCfg, Turing.FinTM.emitterP2PrepareTM, Turing.FinTM.splitPos, Turing.FinTM.splitPos_read, Turing.FinTM.splitPos_succ, Turing.MultiTapeTM.runFrom, Turing.MultiTapeTM.runFrom_succ_eq_step', Turing.MultiTapeTM.step}
\difficulty{4}
For every $j \le \abs{s}$, $j$ steps of the preparation machine from its
start copy the first $j$ bits of the candidate $s$ onto the second tape,
advance the native head to its saturated position for $j$, and set the
overflow flag exactly when $\abs{w} < j$; the candidate itself never
changes.
\end{lemma}

\begin{lemma}[Copying the native suffix]
\lean{Turing.FinTM.emitterP2_prepare_suffix}
\label{Turing.FinTM.emitterP2_prepare_suffix}
\leanok
\uses{Turing.Action.apply, Turing.Cfg.inputSymbol, Turing.Cfg.workTapeSymbols, Turing.FinTM.bufferTape, Turing.FinTM.bufferTape_append, Turing.FinTM.emitterP2PrepareCfg, Turing.FinTM.emitterP2PrepareTM, Turing.FinTM.splitPos, Turing.FinTM.splitPos_read, Turing.FinTM.splitPos_succ, Turing.MultiTapeTM.runFrom, Turing.MultiTapeTM.runFrom_succ_eq_step', Turing.MultiTapeTM.step}
\difficulty{4}
For every $j$ at most the length of the native suffix after position
$\abs{s}$, $j$ steps of the preparation machine in its suffix phase copy the
first $j$ bits of that suffix onto the third tape and advance the native
head accordingly; the other tapes are unchanged.
\end{lemma}

\begin{lemma}[Joint rewind of candidate and copy]
\lean{Turing.FinTM.emitterP2_prepare_rewind_candidate}
\label{Turing.FinTM.emitterP2_prepare_rewind_candidate}
\leanok
\uses{Turing.Action.apply, Turing.Cfg.workTapeSymbols, Turing.FinTM.bufferTape_left, Turing.FinTM.bufferTape_nat, Turing.FinTM.emitterP2PrepareCfg, Turing.FinTM.emitterP2PrepareTM, Turing.MultiTapeTM.runFrom, Turing.MultiTapeTM.runFrom_succ_eq_step, Turing.MultiTapeTM.runFrom_zero, Turing.MultiTapeTM.step, Turing.moveInputPos_zero}
\difficulty{4}
From the preparation machine's candidate-rewind phase with the aligned
candidate and copy heads at displacement $j - 1$, for any $j \le \abs{s}$,
exactly $j+1$ steps restore both heads to the origin without changing
either word; the native input position and the prepared suffix are fixed.
\end{lemma}

\begin{lemma}[Rewind of the suffix tape]
\lean{Turing.FinTM.emitterP2_prepare_rewind_suffix}
\label{Turing.FinTM.emitterP2_prepare_rewind_suffix}
\leanok
\uses{Turing.Action.apply, Turing.Cfg.workTapeSymbols, Turing.FinTM.bufferTape_left, Turing.FinTM.bufferTape_nat, Turing.FinTM.emitterP2PrepareCfg, Turing.FinTM.emitterP2PrepareTM, Turing.MultiTapeTM.runFrom, Turing.MultiTapeTM.runFrom_succ_eq_step, Turing.MultiTapeTM.runFrom_zero, Turing.MultiTapeTM.step, Turing.moveInputPos_zero}
\difficulty{4}
Let $j$ be at most the length of the suffix word. From the preparation
machine's suffix-rewind phase with the suffix head at displacement $j - 1$,
exactly $j+1$ steps yield the frame of the next phase with the suffix head at
the origin; all three words are preserved.
\end{lemma}

\begin{lemma}[Joining exact phase endpoints]
\lean{Turing.FinTM.emitterP2_join}
\label{Turing.FinTM.emitterP2_join}
\leanok
\uses{Turing.Cfg, Turing.MultiTapeTM, Turing.MultiTapeTM.runFrom, Turing.MultiTapeTM.runFrom_add}
\difficulty{1}
If the run from $c$ reaches $d$ in exactly $a$ steps and the run from $d$
reaches $e$ in exactly $b$ steps, then the run from $c$ reaches $e$ in
exactly $a + b$ steps.
\end{lemma}

\begin{lemma}[Complete preparation run]
\lean{Turing.FinTM.emitterP2_prepare_run}
\label{Turing.FinTM.emitterP2_prepare_run}
\leanok
\uses{Turing.Action.apply, Turing.Cfg.inputSymbol, Turing.Cfg.workTapeSymbols, Turing.FinTM.bufferTape_nat, Turing.FinTM.catalogRewind, Turing.FinTM.controlAction, Turing.FinTM.controlAction_apply, Turing.FinTM.emitterP2PrepareCfg, Turing.FinTM.emitterP2PrepareTM, Turing.FinTM.emitterP2_join, Turing.FinTM.emitterP2_prepare_candidate, Turing.FinTM.emitterP2_prepare_rewind_candidate, Turing.FinTM.emitterP2_prepare_rewind_suffix, Turing.FinTM.emitterP2_prepare_suffix, Turing.FinTM.splitPos, Turing.MultiTapeTM.runFrom, Turing.MultiTapeTM.runFrom_succ_eq_step, Turing.MultiTapeTM.runFrom_succ_eq_step', Turing.MultiTapeTM.step, Turing.moveInputPos_zero}
\difficulty{5}
For any input $w$ and candidate $s$, there is a time
$t \le 5(\abs{s} + \abs{w} + 3)$ after which the candidate-and-suffix
preparation machine, started from its initial frame, presents the candidate
$s$ and the native suffix of $w$ after position $\abs{s}$ on clean argument
tapes, with every head and the native input position restored and the
overflow flag recording whether $\abs{w} < \abs{s}$.
\end{lemma}

\begin{lemma}[First return of preparation]
\lean{Turing.FinTM.emitterP2_prepare_first}
\label{Turing.FinTM.emitterP2_prepare_first}
\leanok
\uses{Turing.Action.apply, Turing.FinTM.controlAction, Turing.FinTM.controlAction_apply, Turing.FinTM.emitterP2PrepareCfg, Turing.FinTM.emitterP2PrepareTM, Turing.FinTM.emitterP2_prepare_run, Turing.FinTM.emitter_first_entry, Turing.MultiTapeTM.runFrom, Turing.MultiTapeTM.step, Turing.moveInputPos_zero}
\difficulty{3}
For any input $w$ and candidate $s$, there is a time $t$ with
$0 < t \le 5(\abs{s} + \abs{w} + 3)$ such that the preparation machine never
visits its return phase before $t$ and at time $t$ sits in its completed
frame with the candidate and native suffix prepared; positivity holds even
for empty input and an empty candidate.
\end{lemma}

\begin{lemma}[Relocated phase with anchor exclusion]
\lean{Turing.FinTM.emitterP2_segment}
\label{Turing.FinTM.emitterP2_segment}
\leanok
\uses{Turing.Cfg, Turing.FinTM.emitterP2Action, Turing.FinTM.emitterP2Cfg, Turing.FinTM.emitterP2_relocate_run, Turing.FinTM.splitSafe, Turing.MultiTapeTM, Turing.MultiTapeTM.runFrom}
\difficulty{3}
Under the relocation hypotheses --- the selection map inverts the tape
indexing, no embedded state equals the host anchor, and host transitions on
embedded good states equal the relocated source transitions --- an exact
source run from $c$ to $d$ in $t$ steps, all of whose visited states are
good, relocates to an exact host run between the relocated configurations
whose whole trace avoids the anchor.
\end{lemma}

\begin{lemma}[Entering a clean-call phase]
\lean{Turing.FinTM.emitterP2_call_segment}
\label{Turing.FinTM.emitterP2_call_segment}
\leanok
\uses{Turing.Cfg, Turing.Cfg.workTapeSymbols, Turing.FinTM.emitterP2Action, Turing.FinTM.emitterP2Cfg, Turing.FinTM.emitterP2_apply, Turing.FinTM.emitterP2_relocate_run, Turing.FinTM.splitSafe, Turing.MultiTapeTM, Turing.MultiTapeTM.runFrom, Turing.MultiTapeTM.runFrom_succ_eq_step, Turing.MultiTapeTM.step}
\difficulty{5}
Let a host enter a relocated module through a dedicated entry control state
that performs the module entry state's action unconditionally, and agree
with the relocated module transitions away from the module's exit state. If
the module runs from its entry configuration to $d$ in $t > 0$ steps with no
interior visit to the exit state, then $t$ host steps from the entry
configuration land on the relocation of $d$, with the whole trace avoiding
the host anchor; this covers modules whose entry and exit states coincide.
\end{lemma}

\begin{definition}[Host word layout]
\lean{Turing.FinTM.emitterP2Words}
\label{Turing.FinTM.emitterP2Words}
\leanok
\uses{Turing.stateWord}
The layout of stored words across the $k + l + 1$ host tapes: the preserved
candidate first, then the width-call bank of $k$ tapes, then the
suffix-length-call bank of $l$ tapes; at a seam only the first tape of each
bank carries a word and all other tapes are blank.
\end{definition}

\begin{definition}[Width-bank tape injection]
\lean{Turing.FinTM.emitterP2LeftIndex}
\label{Turing.FinTM.emitterP2LeftIndex}
\leanok
For all bank sizes $k$ and $l$, the width-bank injection sends tape $i$ of the
width-call machine (with $i < k$) to host position $i + 1$, immediately
after the preserved candidate tape at position $0$.
\end{definition}

\begin{definition}[Width-bank selection]
\lean{Turing.FinTM.emitterP2LeftSelect}
\label{Turing.FinTM.emitterP2LeftSelect}
\leanok
The partial selection map recognizing exactly the width-call bank among the
host tapes.
\end{definition}

\begin{definition}[Length-bank tape injection]
\lean{Turing.FinTM.emitterP2RightIndex}
\label{Turing.FinTM.emitterP2RightIndex}
\leanok
For all bank sizes $k$ and $l$, the length-bank injection sends tape $i$ of the
suffix-length-call machine (with $i < l$) to host position $k + i + 1$,
immediately after the width-call bank.
\end{definition}

\begin{definition}[Length-bank selection]
\lean{Turing.FinTM.emitterP2RightSelect}
\label{Turing.FinTM.emitterP2RightSelect}
\leanok
The partial selection map recognizing exactly the suffix-length-call bank
among the host tapes.
\end{definition}

\begin{lemma}[Width-bank selection inverts injection]
\lean{Turing.FinTM.emitterP2_left_inverse}
\label{Turing.FinTM.emitterP2_left_inverse}
\leanok
\uses{Turing.FinTM.emitterP2LeftIndex, Turing.FinTM.emitterP2LeftSelect}
\difficulty{1}
Selecting the host slot at the width-bank injection of index $i$ returns
exactly $i$.
\end{lemma}

\begin{lemma}[Length-bank selection inverts injection]
\lean{Turing.FinTM.emitterP2_right_inverse}
\label{Turing.FinTM.emitterP2_right_inverse}
\leanok
\uses{Turing.FinTM.emitterP2RightIndex, Turing.FinTM.emitterP2RightSelect}
\difficulty{1}
Selecting the host slot at the length-bank injection of index $i$ returns
exactly $i$.
\end{lemma}

\begin{lemma}[Frame of a clean width call]
\lean{Turing.FinTM.emitterP2_left_frame}
\label{Turing.FinTM.emitterP2_left_frame}
\leanok
\uses{Turing.Cfg.ofWords, Turing.FinTM.bufferTape, Turing.FinTM.emitterP2Cfg, Turing.FinTM.emitterP2LeftSelect, Turing.FinTM.emitterP2Words, Turing.stateWord}
\difficulty{3}
Relocating a canonical state-word configuration of the width-call machine
into the host layout --- with the candidate and suffix words preserved on
their ambient tapes --- yields exactly the host's canonical configuration on
the full word layout with the width word installed.
\end{lemma}

\begin{lemma}[Frame of a clean length call]
\lean{Turing.FinTM.emitterP2_right_frame}
\label{Turing.FinTM.emitterP2_right_frame}
\leanok
\uses{Turing.Cfg.ofWords, Turing.FinTM.bufferTape, Turing.FinTM.emitterP2Cfg, Turing.FinTM.emitterP2RightSelect, Turing.FinTM.emitterP2Words, Turing.stateWord}
\difficulty{3}
Relocating a canonical state-word configuration of the suffix-length-call
machine into the host layout --- with the candidate and width words
preserved --- yields exactly the host's canonical configuration on the full
word layout with the length word installed.
\end{lemma}

\begin{lemma}[Clean layout is the loop seam]
\lean{Turing.FinTM.emitterP2_words_clean}
\label{Turing.FinTM.emitterP2_words_clean}
\leanok
\uses{Turing.FinTM.emitterP2Words, Turing.stateWord}
\difficulty{2}
With both argument words erased, the host word layout for candidate $s$
equals the loop's canonical state word on $k + l + 1$ tapes carrying $s$,
including every inactive tape.
\end{lemma}

\begin{definition}[Single-tape selection]
\lean{Turing.FinTM.emitterP2OneSelect}
\label{Turing.FinTM.emitterP2OneSelect}
\leanok
The partial selection map recognizing one designated physical tape, so that
a one-tape phase relocated through it leaves every other tape fixed.
\end{definition}

\begin{lemma}[Single-tape selection inverts injection]
\lean{Turing.FinTM.emitterP2_one_inverse}
\label{Turing.FinTM.emitterP2_one_inverse}
\leanok
\uses{Turing.FinTM.emitterP2OneSelect}
\difficulty{1}
Selecting the designated slot through the single-tape selection map returns
the unique one-tape index.
\end{lemma}

\begin{lemma}[Frame of a one-tape phase]
\lean{Turing.FinTM.emitterP2_one_frame}
\label{Turing.FinTM.emitterP2_one_frame}
\leanok
\uses{Turing.Cfg.ofWords, Turing.FinTM.bufferTape, Turing.FinTM.emitterP2Cfg, Turing.FinTM.emitterP2OneSelect}
\difficulty{3}
For a host slot, a host word layout, a physical input $w$, a word $u$, and a
state $q$, relocating the canonical one-word configuration of a single-tape
phase (state $q$, word $u$) into that slot, with every other tape holding its
layout word at head position $0$, yields the canonical host configuration
whose layout equals the given layout with the slot's word replaced by $u$.
\end{lemma}

\begin{definition}[Administrative-triple injection]
\lean{Turing.FinTM.emitterP2SmallIndex}
\label{Turing.FinTM.emitterP2SmallIndex}
\leanok
\uses{Turing.FinTM.emitterP2LeftIndex, Turing.FinTM.emitterP2RightIndex}
For all positive bank sizes $k$ and $l$, the administrative-triple injection
sends index $0$ to the candidate position $0$, index $1$ to host position
$1$ (the first width-bank tape), and index $2$ to host position $k + 1$ (the
first length-bank tape).
\end{definition}

\begin{definition}[Administrative-triple selection]
\lean{Turing.FinTM.emitterP2SmallSelect}
\label{Turing.FinTM.emitterP2SmallSelect}
\leanok
The partial selection map recognizing the three administrative word slots
--- candidate, width word, and length word --- and leaving all module
scratch tapes inactive.
\end{definition}

\begin{lemma}[Triple selection inverts injection]
\lean{Turing.FinTM.emitterP2_small_inverse}
\label{Turing.FinTM.emitterP2_small_inverse}
\leanok
\uses{Turing.FinTM.emitterP2LeftIndex, Turing.FinTM.emitterP2RightIndex, Turing.FinTM.emitterP2SmallIndex, Turing.FinTM.emitterP2SmallSelect}
\difficulty{2}
For positive bank sizes, selecting the host slot at the
administrative-triple injection of any of the three indices returns that
index.
\end{lemma}

\begin{lemma}[Frame of relocated preparation]
\lean{Turing.FinTM.emitterP2_small_frame}
\label{Turing.FinTM.emitterP2_small_frame}
\leanok
\uses{Turing.Cfg.ofWords, Turing.FinTM.bufferTape_nil, Turing.FinTM.emitterP2Cfg, Turing.FinTM.emitterP2PrepareCfg, Turing.FinTM.emitterP2SmallSelect, Turing.FinTM.emitterP2Words, Turing.stateWord}
\difficulty{3}
At zero head displacements, relocating a preparation frame through the
administrative-triple selection coincides exactly with the host's canonical
configuration on the complete word layout, not merely with its two argument
projections.
\end{lemma}

\begin{definition}[Comparator-pair injection]
\lean{Turing.FinTM.emitterP2PairIndex}
\label{Turing.FinTM.emitterP2PairIndex}
\leanok
\uses{Turing.FinTM.emitterP2LeftIndex, Turing.FinTM.emitterP2RightIndex}
The injection sending the comparator's two tapes to the first tapes of the
two call banks, which hold the two words to be compared.
\end{definition}

\begin{definition}[Comparator-pair selection]
\lean{Turing.FinTM.emitterP2PairSelect}
\label{Turing.FinTM.emitterP2PairSelect}
\leanok
The partial selection map recognizing the two comparison word slots while
keeping the candidate and all scratch tapes fixed.
\end{definition}

\begin{lemma}[Pair selection inverts injection]
\lean{Turing.FinTM.emitterP2_pair_inverse}
\label{Turing.FinTM.emitterP2_pair_inverse}
\leanok
\uses{Turing.FinTM.emitterP2LeftIndex, Turing.FinTM.emitterP2PairIndex, Turing.FinTM.emitterP2PairSelect, Turing.FinTM.emitterP2RightIndex}
\difficulty{2}
For positive bank sizes, selecting the host slot at the comparator-pair
injection of either index returns that index.
\end{lemma}

\begin{lemma}[Frame of the relocated comparison]
\lean{Turing.FinTM.emitterP2_pair_frame}
\label{Turing.FinTM.emitterP2_pair_frame}
\leanok
\uses{Turing.Cfg.ofWords, Turing.FinTM.bufferTape, Turing.FinTM.bufferTape_nil, Turing.FinTM.emitterCompareCfg, Turing.FinTM.emitterP2Cfg, Turing.FinTM.emitterP2PairSelect, Turing.FinTM.emitterP2Words, Turing.stateWord}
\difficulty{3}
Relocating a comparator configuration at zero displacement through the
comparator-pair selection yields the host's canonical configuration whose
layout holds the two whole compared words together with the unchanged
candidate, every module scratch tape being blank.
\end{lemma}

\begin{lemma}[Updating the width word]
\lean{Turing.FinTM.emitterP2_update_left}
\label{Turing.FinTM.emitterP2_update_left}
\leanok
\uses{Turing.FinTM.emitterP2LeftIndex, Turing.FinTM.emitterP2RightIndex, Turing.FinTM.emitterP2Words, Turing.stateWord}
\difficulty{4}
Replacing the word at the first width-bank slot of the host layout for
$(s, u, v)$ by a word $a$ yields exactly the layout for $(s, a, v)$; every
other slot is untouched.
\end{lemma}

\begin{lemma}[Updating the length word]
\lean{Turing.FinTM.emitterP2_update_right}
\label{Turing.FinTM.emitterP2_update_right}
\leanok
\uses{Turing.FinTM.emitterP2LeftIndex, Turing.FinTM.emitterP2RightIndex, Turing.FinTM.emitterP2Words, Turing.stateWord}
\difficulty{4}
Replacing the word at the first length-bank slot of the host layout for
$(s, u, v)$ by a word $a$ yields exactly the layout for $(s, u, a)$; the
candidate and width words are preserved.
\end{lemma}

\begin{lemma}[Updating the candidate word]
\lean{Turing.FinTM.emitterP2_update_candidate}
\label{Turing.FinTM.emitterP2_update_candidate}
\leanok
\uses{Turing.FinTM.emitterP2Words}
\difficulty{2}
For all words $s$, $u$, $v$, and $a$, replacing the word at the candidate
slot of the host layout for $(s, u, v)$ by $a$ yields exactly the layout for
$(a, u, v)$.
\end{lemma}

\begin{definition}[Controller phases of the width-split round]
\lean{Turing.FinTM.EmitterP2State}
\label{Turing.FinTM.EmitterP2State}
\leanok
Given state types $S$ and $T$ of the width and length call machines, the
controller-phase type of the width-split round has an anchor, a preparation
phase, dedicated start states and running phases (carrying a state of $S$,
respectively $T$) for the width call and the suffix-length call, a
comparison phase, two erasure phases carrying the verdict, an advance
(cleanup) phase, and an emission phase.
\end{definition}

\begin{definition}[Width-split round controller]
\lean{Turing.FinTM.emitterP2BodyTM}
\label{Turing.FinTM.emitterP2BodyTM}
\leanok
\uses{Turing.FinTM, Turing.FinTM.EmitterP2State, Turing.FinTM.controlAction, Turing.FinTM.emitterCompareTM, Turing.FinTM.emitterP2Action, Turing.FinTM.emitterP2EraseTM, Turing.FinTM.emitterP2LeftIndex, Turing.FinTM.emitterP2LeftSelect, Turing.FinTM.emitterP2OneSelect, Turing.FinTM.emitterP2PairIndex, Turing.FinTM.emitterP2PairSelect, Turing.FinTM.emitterP2PrepareTM, Turing.FinTM.emitterP2RightIndex, Turing.FinTM.emitterP2RightSelect, Turing.FinTM.emitterP2SmallIndex, Turing.FinTM.emitterP2SmallSelect, Turing.FinTM.splitEmitTM, Turing.FinTM.splitRestoreTM}
Given clean-call machines $C$ (width evaluator) and $D$ (suffix-length
evaluator) with designated entry and exit states, the native controller of
one width-split round: it prepares the candidate copy and native suffix,
runs the two relocated calls on their disjoint banks, compares the two
returned words whole-word, erases both result words while carrying the
verdict, and then either emits the native split or advances the candidate
and returns to the anchor; a past-the-end preparation bypasses both
evaluators, the native input is never overwritten, and no time bound or
width function appears in the transition table.
\end{definition}

\begin{lemma}[Startup of the width-split controller]
\lean{Turing.FinTM.emitterP2_body_start}
\label{Turing.FinTM.emitterP2_body_start}
\leanok
\uses{Turing.Cfg.ofWords, Turing.FinTM, Turing.FinTM.emitterP2BodyTM, Turing.MultiTapeTM.initCfg, Turing.initCfg_ofWords, Turing.stateWord}
\difficulty{2}
On every input $w$, the width-split controller's initial configuration
equals the canonical anchored configuration whose stored words carry the
empty candidate.
\end{lemma}

\begin{lemma}[Silent dispatch on a word seam]
\lean{Turing.FinTM.emitterP2_control}
\label{Turing.FinTM.emitterP2_control}
\leanok
\uses{Turing.Cfg.ofWords, Turing.FinTM.controlAction, Turing.FinTM.controlAction_apply, Turing.MultiTapeTM, Turing.MultiTapeTM.step, Turing.moveInputPos_zero}
\difficulty{2}
If a control state's transition is the unconditional silent dispatch to a
state $r$, then one step from the canonical word-seam configuration in that
state yields the canonical word-seam configuration in state $r$ with every
other component unchanged.
\end{lemma}

\begin{lemma}[Preparation inside the controller]
\lean{Turing.FinTM.emitterP2_body_prepare}
\label{Turing.FinTM.emitterP2_body_prepare}
\leanok
\uses{Turing.Cfg.ofWords, Turing.FinTM, Turing.FinTM.emitterP2BodyTM, Turing.FinTM.emitterP2PrepareCfg, Turing.FinTM.emitterP2PrepareTM, Turing.FinTM.emitterP2SmallIndex, Turing.FinTM.emitterP2SmallSelect, Turing.FinTM.emitterP2Words, Turing.FinTM.emitterP2_prepare_first, Turing.FinTM.emitterP2_segment, Turing.FinTM.emitterP2_small_frame, Turing.FinTM.emitterP2_small_inverse, Turing.FinTM.splitSafe, Turing.MultiTapeTM.runFrom}
\difficulty{4}
For any input $w$ and candidate $s$, there is a time
$t \le 5(\abs{s} + \abs{w} + 3)$ such that the controller, from its
preparation seam carrying $s$, reaches the seam holding the candidate, its
copy, and the native suffix of $w$ after position $\abs{s}$, with the
past-the-end flag recording $\abs{w} < \abs{s}$; the trace avoids the
anchor.
\end{lemma}

\begin{lemma}[Width call inside the controller]
\lean{Turing.FinTM.emitterP2_body_width}
\label{Turing.FinTM.emitterP2_body_width}
\leanok
\uses{Turing.Cfg.ofWords, Turing.FinTM, Turing.FinTM.bufferTape, Turing.FinTM.emitterP2BodyTM, Turing.FinTM.emitterP2LeftIndex, Turing.FinTM.emitterP2LeftSelect, Turing.FinTM.emitterP2Words, Turing.FinTM.emitterP2_call_segment, Turing.FinTM.emitterP2_left_frame, Turing.FinTM.emitterP2_left_inverse, Turing.FinTM.splitSafe, Turing.MultiTapeTM.runFrom, Turing.stateWord}
\difficulty{4}
If the width machine, started on its canonical seam carrying the candidate
copy, reaches its exit seam with result word $u$ in $t > 0$ steps with no
interior exit-state visit, then $t$ controller steps from the width-start
seam install $u$ in the width bank while retaining the original candidate
and prepared suffix; the trace avoids the anchor.
\end{lemma}

\begin{lemma}[Length call inside the controller]
\lean{Turing.FinTM.emitterP2_body_length}
\label{Turing.FinTM.emitterP2_body_length}
\leanok
\uses{Turing.Cfg.ofWords, Turing.FinTM, Turing.FinTM.bufferTape, Turing.FinTM.emitterP2BodyTM, Turing.FinTM.emitterP2RightIndex, Turing.FinTM.emitterP2RightSelect, Turing.FinTM.emitterP2Words, Turing.FinTM.emitterP2_call_segment, Turing.FinTM.emitterP2_right_frame, Turing.FinTM.emitterP2_right_inverse, Turing.FinTM.splitSafe, Turing.MultiTapeTM.runFrom, Turing.stateWord}
\difficulty{4}
If the suffix-length machine, started on its canonical seam carrying the
prepared suffix, reaches its exit seam with result word $z$ in $t > 0$
steps with no interior exit-state visit, then $t$ controller steps from the
length-start seam install $z$ in the length bank while preserving the
candidate and the already installed width word; the trace avoids the
anchor.
\end{lemma}

\begin{lemma}[Comparison inside the controller]
\lean{Turing.FinTM.emitterP2_body_compare}
\label{Turing.FinTM.emitterP2_body_compare}
\leanok
\uses{Turing.Cfg.ofWords, Turing.FinTM, Turing.FinTM.bufferTape, Turing.FinTM.emitterCompareCfg, Turing.FinTM.emitterCompareTM, Turing.FinTM.emitterP2BodyTM, Turing.FinTM.emitterP2PairIndex, Turing.FinTM.emitterP2PairSelect, Turing.FinTM.emitterP2Words, Turing.FinTM.emitterP2_pair_frame, Turing.FinTM.emitterP2_pair_inverse, Turing.FinTM.emitterP2_segment, Turing.FinTM.emitter_compare_first, Turing.FinTM.splitSafe, Turing.MultiTapeTM.runFrom}
\difficulty{4}
For installed result words $u$ and $v$, there is a time
$t \le 2(\max(\abs{u},\abs{v}) + 1)$ after which the controller, from its
comparison seam, records in finite control whether $u = v$, changing no
candidate, buffer, or module scratch cell; the trace avoids the anchor.
\end{lemma}

\begin{lemma}[Erasing the width word inside the controller]
\lean{Turing.FinTM.emitterP2_body_erase_left}
\label{Turing.FinTM.emitterP2_body_erase_left}
\leanok
\uses{Turing.Cfg.ofWords, Turing.FinTM, Turing.FinTM.EmitterP2State, Turing.FinTM.bufferTape, Turing.FinTM.emitterP2BodyTM, Turing.FinTM.emitterP2Cfg, Turing.FinTM.emitterP2EraseCfg, Turing.FinTM.emitterP2EraseTM, Turing.FinTM.emitterP2LeftIndex, Turing.FinTM.emitterP2OneSelect, Turing.FinTM.emitterP2Words, Turing.FinTM.emitterP2_erase_first, Turing.FinTM.emitterP2_one_frame, Turing.FinTM.emitterP2_one_inverse, Turing.FinTM.emitterP2_segment, Turing.FinTM.emitterP2_update_left, Turing.FinTM.splitSafe, Turing.MultiTapeTM.runFrom}
\difficulty{4}
For any verdict bit, there is a time $t \le 2(\abs{u}+1)$ after which the
controller, from its first erasure seam, has blanked the installed width
word $u$ while retaining the verdict in control and preserving every other
tape; the trace avoids the anchor.
\end{lemma}

\begin{lemma}[Erasing the length word inside the controller]
\lean{Turing.FinTM.emitterP2_body_erase_right}
\label{Turing.FinTM.emitterP2_body_erase_right}
\leanok
\uses{Turing.Cfg.ofWords, Turing.FinTM, Turing.FinTM.EmitterP2State, Turing.FinTM.bufferTape, Turing.FinTM.emitterP2BodyTM, Turing.FinTM.emitterP2Cfg, Turing.FinTM.emitterP2EraseCfg, Turing.FinTM.emitterP2EraseTM, Turing.FinTM.emitterP2OneSelect, Turing.FinTM.emitterP2RightIndex, Turing.FinTM.emitterP2Words, Turing.FinTM.emitterP2_erase_first, Turing.FinTM.emitterP2_one_frame, Turing.FinTM.emitterP2_one_inverse, Turing.FinTM.emitterP2_segment, Turing.FinTM.emitterP2_update_right, Turing.FinTM.splitSafe, Turing.MultiTapeTM.runFrom}
\difficulty{4}
For any verdict bit, there is a time $t \le 2(\abs{v}+1)$ after which the
controller, from its second erasure seam, has blanked the installed length
word $v$ while retaining the verdict through the last cleanup phase; after
this all scratch and administrative words are blank. The trace avoids the
anchor.
\end{lemma}

\begin{lemma}[Entry frame of the candidate advance]
\lean{Turing.FinTM.emitterP2_advance_initial}
\label{Turing.FinTM.emitterP2_advance_initial}
\leanok
\uses{Turing.Cfg.ofWords, Turing.FinTM.splitPos, Turing.FinTM.splitRestoreScan}
\difficulty{2}
The one-tape candidate-advance machine's starting scan configuration equals
the canonical word-seam configuration carrying the candidate on its single
tape.
\end{lemma}

\begin{lemma}[One-tape state words are constant]
\lean{Turing.FinTM.emitterP2_stateWord_one}
\label{Turing.FinTM.emitterP2_stateWord_one}
\leanok
\uses{Turing.stateWord}
\difficulty{1}
The canonical state word on one tape carrying $s$ is the constant function
with value $s$.
\end{lemma}

\begin{lemma}[Candidate advance inside the controller]
\lean{Turing.FinTM.emitterP2_body_advance}
\label{Turing.FinTM.emitterP2_body_advance}
\leanok
\uses{Turing.Cfg.ofWords, Turing.FinTM, Turing.FinTM.EmitterP2State, Turing.FinTM.bufferTape, Turing.FinTM.emitterP2BodyTM, Turing.FinTM.emitterP2OneSelect, Turing.FinTM.emitterP2Words, Turing.FinTM.emitterP2_advance_initial, Turing.FinTM.emitterP2_one_frame, Turing.FinTM.emitterP2_one_inverse, Turing.FinTM.emitterP2_segment, Turing.FinTM.emitterP2_stateWord_one, Turing.FinTM.emitterP2_update_candidate, Turing.FinTM.splitRestoreScan, Turing.FinTM.splitRestoreTM, Turing.FinTM.splitRestore_first, Turing.FinTM.splitSafe, Turing.FinTM.splitStep, Turing.MultiTapeTM.runFrom, Turing.stateWord}
\difficulty{4}
There is a time $t \le 2\abs{s} + \abs{w} + 5$ after which the controller,
from its advance seam carrying the candidate $s$, has replaced $s$ by one
split-search step applied to $s$ (appending a $1$ exactly within the input
range, and stalling otherwise) and reached the advance phase's return seam;
the trace avoids the anchor.
\end{lemma}

\begin{lemma}[Entry frame of the payload emitter]
\lean{Turing.FinTM.emitterP2_emit_initial}
\label{Turing.FinTM.emitterP2_emit_initial}
\leanok
\uses{Turing.Cfg.ofWords, Turing.FinTM.splitEmitCfg, Turing.FinTM.splitPos}
\difficulty{2}
The one-tape payload emitter's starting configuration equals the canonical
word-seam configuration carrying the candidate on its single tape.
\end{lemma}

\begin{lemma}[Payload emission inside the controller]
\lean{Turing.FinTM.emitterP2_body_emit}
\label{Turing.FinTM.emitterP2_body_emit}
\leanok
\uses{Turing.Cfg.ofWords, Turing.FinTM, Turing.FinTM.EmitterP2State, Turing.FinTM.bufferTape, Turing.FinTM.emitterP2BodyTM, Turing.FinTM.emitterP2OneSelect, Turing.FinTM.emitterP2Words, Turing.FinTM.emitterP2_emit_initial, Turing.FinTM.emitterP2_one_frame, Turing.FinTM.emitterP2_one_inverse, Turing.FinTM.emitterP2_segment, Turing.FinTM.emitterP2_update_candidate, Turing.FinTM.splitEmitCfg, Turing.FinTM.splitEmitTM, Turing.FinTM.splitEmit_run, Turing.FinTM.splitSafe, Turing.MultiTapeTM.runFrom, Turing.pairEncode}
\difficulty{4}
If the candidate's length is at most $\abs{w}$, then exactly
$\abs{s} + \abs{w} + 3$ controller steps from the emission seam halt with
output the self-delimiting pair encoding of the split of the native input
$w$ at position $\abs{s}$; the emitted bits come from the native input, not
the candidate, unrelated banks stay blank, and the trace avoids the anchor.
\end{lemma}

\begin{lemma}[Extending a safe phase by one dispatch]
\lean{Turing.FinTM.emitterP2_after}
\label{Turing.FinTM.emitterP2_after}
\leanok
\uses{Turing.Cfg, Turing.Cfg.ofWords, Turing.FinTM.controlAction, Turing.FinTM.emitterP2_control, Turing.FinTM.splitSafe, Turing.FinTM.splitSafe_join, Turing.FinTM.splitSafe_one, Turing.MultiTapeTM, Turing.MultiTapeTM.runFrom}
\difficulty{2}
If an anchor-free run of length $t$ ends on a canonical word seam in a
non-anchor state whose transition is the silent dispatch to another
non-anchor state, then the run of length $t+1$ ends on the same word seam
in the new state and is still anchor-free.
\end{lemma}

\begin{lemma}[Testing segment of the round]
\lean{Turing.FinTM.emitterP2_body_test}
\label{Turing.FinTM.emitterP2_body_test}
\leanok
\uses{Turing.Cfg.ofWords, Turing.FinTM, Turing.FinTM.emitterP2BodyTM, Turing.FinTM.emitterP2Words, Turing.FinTM.emitterP2_after, Turing.FinTM.emitterP2_body_compare, Turing.FinTM.emitterP2_body_length, Turing.FinTM.emitterP2_body_width, Turing.FinTM.splitSafe, Turing.FinTM.splitSafe_join, Turing.MultiTapeTM.runFrom, Turing.stateWord}
\difficulty{5}
Suppose the width call completes in $a > 0$ steps with result $u$ and the
suffix-length call completes in $b > 0$ steps with result $v$, each with no
interior exit-state visit. Then there is a time
$t \le a + b + 2(\max(\abs{u},\abs{v}) + 1) + 4$ after which the
controller, from its completed-preparation seam, has run both calls and the
whole-word comparison and sits at the first erasure seam carrying the
verdict of $u = v$; both calls are charged before any relation between
their results is assumed, and the trace avoids the anchor.
\end{lemma}

\begin{lemma}[Finishing segment of the round]
\lean{Turing.FinTM.emitterP2_body_finish}
\label{Turing.FinTM.emitterP2_body_finish}
\leanok
\uses{Turing.Cfg.ofWords, Turing.FinTM, Turing.FinTM.emitterP2BodyTM, Turing.FinTM.emitterP2Words, Turing.FinTM.emitterP2_after, Turing.FinTM.emitterP2_body_advance, Turing.FinTM.emitterP2_body_emit, Turing.FinTM.emitterP2_body_erase_left, Turing.FinTM.emitterP2_body_erase_right, Turing.FinTM.emitterP2_control, Turing.FinTM.splitSafe, Turing.FinTM.splitSafe_add, Turing.FinTM.splitSafe_join, Turing.FinTM.splitStep, Turing.MultiTapeTM.runFrom, Turing.MultiTapeTM.runFrom_add, Turing.MultiTapeTM.runFrom_succ_eq_step', Turing.pairEncode}
\difficulty{5}
Let a verdict bit be given, true only if the candidate's length is at most
$\abs{w}$. Then there is a time
$t \le 2\abs{u} + 2\abs{v} + 2\abs{s} + \abs{w} + 12$ after which the
controller, from the first erasure seam, has erased both installed words
and: on a true verdict, halted with output the pair encoding of the split
of $w$ at position $\abs{s}$; on a false verdict, returned to the canonical
anchored seam carrying one split-search step applied to $s$. No strictly
interior configuration is at the anchor.
\end{lemma}

\begin{lemma}[Strict join of safety intervals]
\lean{Turing.FinTM.emitterP2_strict_join}
\label{Turing.FinTM.emitterP2_strict_join}
\leanok
\uses{Turing.Cfg, Turing.FinTM.splitSafe, Turing.MultiTapeTM, Turing.MultiTapeTM.runFrom, Turing.MultiTapeTM.runFrom_add}
\difficulty{2}
If the run from $c$ reaches $d$ in $a$ anchor-free steps and the run from
$d$ avoids the anchor at all times strictly before $b$, then the run from
$c$ avoids the anchor at all times strictly before $a + b$, even when time
$a + b$ itself returns to the anchor.
\end{lemma}

\begin{lemma}[One complete round of the width split]
\lean{Turing.FinTM.emitterP2_body_round}
\label{Turing.FinTM.emitterP2_body_round}
\leanok
\uses{Turing.Cfg.ofWords, Turing.FinTM, Turing.FinTM.emitterP2BodyTM, Turing.FinTM.emitterP2Words, Turing.FinTM.emitterP2_after, Turing.FinTM.emitterP2_body_finish, Turing.FinTM.emitterP2_body_prepare, Turing.FinTM.emitterP2_body_test, Turing.FinTM.emitterP2_control, Turing.FinTM.emitterP2_strict_join, Turing.FinTM.emitterP2_words_clean, Turing.FinTM.emitterSplitAccept, Turing.FinTM.emitter_binary_check, Turing.FinTM.splitSafe_join, Turing.FinTM.splitStep, Turing.MultiTapeTM.runFrom, Turing.MultiTapeTM.runFrom_add, Turing.MultiTapeTM.step, Turing.pairEncode, Turing.stateWord}
\difficulty{6}
Suppose the width call on the candidate $s$ returns the binary word of
$f(\abs{s})$ in $a > 0$ steps and the suffix-length call returns the binary
word of the suffix length in $b > 0$ steps, each with no interior
exit-state visit. Then there is a time $t$ with
$0 < t \le a + b + 10(\abs{s} + \abs{w} + \ell_u + \ell_v + 10)$, where
$\ell_u, \ell_v$ are the two binary words' lengths, such that the round
from the canonical anchored seam carrying $s$ never revisits the anchor
strictly between $0$ and $t$, and at time $t$: if
$\abs{s} + f(\abs{s}) = \abs{w}$ the controller has halted with output the
pair encoding of the split of $w$ at $\abs{s}$, and otherwise it is back at
the canonical anchored seam carrying the stepped candidate.
\end{lemma}

\begin{lemma}[Closing the width-parametric split contract]
\lean{Turing.FinTM.emitterP2_closed}
\label{Turing.FinTM.emitterP2_closed}
\leanok
\uses{Turing.FinTM, Turing.FinTM.ComputesFunInTime, Turing.FinTM.computesFunInTime_lengthBits, Turing.FinTM.computesInTime_iff, Turing.FinTM.emitterP2BodyTM, Turing.FinTM.emitterP2_body_round, Turing.FinTM.emitterP2_body_start, Turing.FinTM.emitterSplit_of_body, Turing.FinTM.emitter_width_budget, Turing.FinTM.exists_installCallTM, Turing.MultiTapeTM.output_length_le, Turing.pairEncode, Turing.solveSplitWith}
\difficulty{5}
Let a machine $E$ compute the binary representation of $f$ applied to the
input length within a monotone time bound $T_E$. Then there exist a machine
$M$ and a constant $c$ computing, within
$c\,(n+1)(T_E(n+1) + n + 2)$ steps, the function sending $w$ to the pair
encoding of the split of $w$ at the least $i$ with $i + f(i) = \abs{w}$,
and to the empty word when no such $i$ exists.
\end{lemma}

\begin{theorem}[Width-parametric split search, machine contract]
\lean{Turing.FinTM.computesFunInTime_splitSolveWith}
\label{Turing.FinTM.computesFunInTime_splitSolveWith}
\leanok
\uses{Turing.FinTM, Turing.FinTM.ComputesFunInTime, Turing.FinTM.emitterP2_closed, Turing.pairEncode, Turing.solveSplitWith}
\difficulty{1}
Let $f : \mathbb{N} \to \mathbb{N}$ and let a machine $E$ compute the
function sending an input of length $m$ to the little-endian binary
representation of $f(m)$ within a monotone time bound $T_E$. Then there
exist a finite Turing machine $M$ and a constant $c$ such that $M$
computes, within $c\,(n+1)(T_E(n+1) + n + 2)$ steps on inputs of length
$n$, the function that on input $w$ finds the least $i \le \abs{w}$ with
$i + f(i) = \abs{w}$ and, on success, outputs the self-delimiting pair
encoding of the first $i$ bits of $w$ paired with the rest of $w$, and on
failure outputs the empty word.
\end{theorem}

\begin{definition}[Streaming unary-token encoder]
\lean{Turing.FinTM.emitterTokenTM}
\label{Turing.FinTM.emitterTokenTM}
\leanok
\uses{Turing.FinTM}
A six-state machine without work tapes that splits off its input's leading
unary token in one streaming pass: each token bit is emitted doubled, the
token's terminating $0$ is followed by the separator bits $0,1$, and the
remainder is copied verbatim; a right blank during token scanning starts
the separator directly, so an unterminated token and the empty input are
both retained.
\end{definition}

\begin{lemma}[Doubling one token bit]
\lean{Turing.FinTM.emitterToken_double}
\label{Turing.FinTM.emitterToken_double}
\leanok
\uses{Turing.Action.apply, Turing.Cfg.ext_zero_tapes, Turing.FinTM.emitterTokenTM, Turing.FinTM.scanCfg, Turing.FinTM.scanCfg_read, Turing.MultiTapeTM.runFrom, Turing.MultiTapeTM.runFrom_succ_eq_step', Turing.MultiTapeTM.runFrom_zero, Turing.MultiTapeTM.step, Turing.moveInputPos_pos_of_ne_right}
\difficulty{4}
If the input splits as $x = p\, b\, r$ with bit $b$, two transitions of the
token encoder from its scanning state at position $\abs{p}$ emit $b$ twice
and advance the input head once; the encoder continues token scanning when
$b = 1$ and passes to the separator phase when $b = 0$.
\end{lemma}

\begin{lemma}[Separator emission of the token encoder]
\lean{Turing.FinTM.emitterToken_separator}
\label{Turing.FinTM.emitterToken_separator}
\leanok
\uses{Turing.Action.apply, Turing.Cfg.ext_zero_tapes, Turing.FinTM.emitterTokenTM, Turing.FinTM.scanCfg, Turing.MultiTapeTM.runFrom, Turing.MultiTapeTM.runFrom_succ_eq_step', Turing.MultiTapeTM.runFrom_zero, Turing.MultiTapeTM.step}
\difficulty{2}
For every input $x$, position $i \le \abs{x}$, and prior output $o$, two
transitions of the token encoder from its separator phase at position $i$
yield the suffix-copying state at the same position $i$, with output $o$
followed by the pair delimiter bits $0,1$.
\end{lemma}

\begin{lemma}[Main run of the token encoder]
\lean{Turing.FinTM.emitterToken_run}
\label{Turing.FinTM.emitterToken_run}
\leanok
\uses{Turing.Action.apply, Turing.Cfg.ext_zero_tapes, Turing.FinTM.emitterTokenTM, Turing.FinTM.emitterToken_double, Turing.FinTM.emitterToken_separator, Turing.FinTM.scanCfg, Turing.FinTM.scanCfg_read, Turing.FinTM.scanCopy_suffix, Turing.MultiTapeTM.runFrom, Turing.MultiTapeTM.runFrom_add, Turing.MultiTapeTM.runFrom_succ_eq_step', Turing.MultiTapeTM.runFrom_zero, Turing.MultiTapeTM.step, Turing.pairEncode, Turing.unaryTokenSplit}
\difficulty{5}
Let the unconsumed input $r$ split at its leading unary token into the
token $a$ (the maximal leading run of ones with its terminating zero, if
present) and the remainder $b$. Then, after any consumed prefix and prior
output $o$, exactly $2\abs{a} + \abs{b} + 3$ steps of the token encoder
halt with output $o$ followed by the self-delimiting pair encoding of
$(a, b)$, the final step being the blank-reading halt.
\end{lemma}

\begin{lemma}[Token and remainder partition the input]
\lean{Turing.FinTM.emitterToken_length}
\label{Turing.FinTM.emitterToken_length}
\leanok
\uses{Turing.unaryTokenSplit}
\difficulty{3}
For every word $x$, the length of the leading unary token of $x$ plus the
length of the remainder equals $\abs{x}$; this includes the empty input and
an unterminated token.
\end{lemma}

\begin{theorem}[Unary token step, machine contract]
\lean{Turing.FinTM.computesFunInTime_unaryToken}
\label{Turing.FinTM.computesFunInTime_unaryToken}
\leanok
\uses{Turing.Cfg.ext_zero_tapes, Turing.FinTM, Turing.FinTM.ComputesFunInTime, Turing.FinTM.ComputesInTime, Turing.FinTM.ComputesInTime.mono, Turing.FinTM.computesInTime_iff, Turing.FinTM.emitterTokenTM, Turing.FinTM.emitterToken_length, Turing.FinTM.emitterToken_run, Turing.FinTM.scanCfg, Turing.MultiTapeTM.initCfg, Turing.pairEncode, Turing.unaryTokenSplit}
\difficulty{3}
There exist a finite Turing machine $M$ and a constant $c$ such that $M$
computes, within $c\,(n+1)$ steps on inputs of length $n$, the function
splitting its input at the leading unary token --- the maximal leading run
of ones with its terminating zero, if present --- and emitting the
self-delimiting pair encoding of the token and the remainder.
\end{theorem}

\begin{definition}[Bit-appending machine]
\lean{Turing.FinTM.emitterAppendTM}
\label{Turing.FinTM.emitterAppendTM}
\leanok
\uses{Turing.FinTM}
For a fixed bit $b$, the bit-appending machine is the single-state machine
without work tapes that copies each input bit to the output and emits $b$ on
the right-blank halting transition.
\end{definition}

\begin{lemma}[Copy phase of the bit appender]
\lean{Turing.FinTM.emitterAppend_run}
\label{Turing.FinTM.emitterAppend_run}
\leanok
\uses{Turing.Action.apply, Turing.Cfg.ext_zero_tapes, Turing.Cfg.inputSymbol, Turing.FinTM.emitterAppendTM, Turing.FinTM.scanCfg, Turing.FinTM.scanCfg_read, Turing.MultiTapeTM.initCfg, Turing.MultiTapeTM.runFrom, Turing.MultiTapeTM.runFrom_succ_eq_step', Turing.MultiTapeTM.step, Turing.moveInputPos_pos_of_ne_right}
\difficulty{4}
For every $j \le \abs{x}$, after $j$ steps from its initial configuration
on input $x$ the bit appender is at input position $j$ with output exactly
the first $j$ bits of $x$; the appending halt is still ahead.
\end{lemma}

\begin{theorem}[Appending one fixed bit, machine contract]
\lean{Turing.FinTM.computesFunInTime_appendBit}
\label{Turing.FinTM.computesFunInTime_appendBit}
\leanok
\uses{Turing.Action.apply, Turing.Cfg.Halted, Turing.Cfg.inputSymbol, Turing.FinTM, Turing.FinTM.ComputesFunInTime, Turing.FinTM.computesInTime_iff, Turing.FinTM.emitterAppendTM, Turing.FinTM.emitterAppend_run, Turing.FinTM.scanCfg, Turing.FinTM.scanCfg_read, Turing.MultiTapeTM.runFrom_succ_eq_step', Turing.MultiTapeTM.step}
\difficulty{3}
For every fixed bit $b$ there exist a finite Turing machine $M$ and a
constant $c$ such that $M$ computes the function $x \mapsto x b$ (append
the bit $b$ to the input word) within $c\,(n+1)$ steps on inputs of length
$n$.
\end{theorem}
